\documentclass[11pt]{article}
\usepackage[a4paper,margin=22mm]{geometry}
\usepackage[no-math]{fontspec}
\setmainfont{lmroman10-regular.otf}[BoldFont=lmroman10-bold.otf,ItalicFont=lmroman10-italic.otf,BoldItalicFont=lmroman10-bolditalic.otf]
\usepackage{amsmath,amssymb,amsthm,mathtools,graphicx,xcolor,hyperref,xurl,xspace,ifthen}
\setlength{\emergencystretch}{4em}
\SetMathAlphabet{\mathbf}{normal}{OT1}{cmr}{bx}{n}
\SetMathAlphabet{\mathbf}{bold}{OT1}{cmr}{bx}{n}
\newtheorem{theo}{Theorem}[section]
\newtheorem{lemm}[theo]{Lemma}
\newtheorem{theorem}[theo]{Theorem}
\newtheorem*{remas*}{Remarks}
\newtheorem{remas}[theo]{Remarks}
\newtheorem*{exam*}{Example}
\newtheorem{ques}[theo]{Question}
\providecommand{\og}{«\nobreakspace}
\providecommand{\fg}{\nobreakspace»}
\newtheorem{lem}[theo]{Lemma}
\newtheorem{lemma}[theo]{Lemma}
\newtheorem{prop}[theo]{Proposition}
\newtheorem{coro}[theo]{Corollary}
\newtheorem{corol}[theo]{Corollary}
\newtheorem{conj}[theo]{Conjecture}
\newtheorem{Proposition}[theo]{Proposition}
\newtheorem{theoreme}[theo]{Theorem}
\theoremstyle{definition}
\newtheorem{defi}[theo]{Definition}
\newtheorem{exam}[theo]{Example}
\newtheorem{exem}[theo]{Example}
\newtheorem{remark}[theo]{Remark}
\newtheorem{example}[theo]{Example}
\newtheorem{rema}[theo]{Remark}
\newtheorem{nota}[theo]{Notation}
\newtheorem*{rema*}{Remark}
\newtheorem*{defi*}{Definition}
\newtheorem*{theo*}{Theorem}
\newtheorem*{lemm*}{Lemma}
\newtheorem*{prop*}{Proposition}
\newcommand{\enoncename}{}
\newtheorem{corpusenonce}[theo]{\enoncename}
\newenvironment{enonce}[2][]{\renewcommand{\enoncename}{#2}\begin{corpusenonce}}{\end{corpusenonce}}
\newenvironment{enonce*}[2][]{\par\medskip\noindent\textbf{#2.}\itshape}{\par\medskip}
\newenvironment{proofc}[1][\proofname]{\begin{proof}[#1]}{\end{proof}}
\newcommand{\Mc}{\mathcal{M}}
\newcommand{\pointir}{.\ ---\ }
\providecommand{\mathof}[1]{\mathrm{#1}}

\numberwithin{equation}{section}
\numberwithin{figure}{section}
\numberwithin{table}{section}
\def\endappendix{}
%~Mouliné par MaN_auto v.0.21.6 2019-08-22 10:59:50


\usepackage[all]{xy}
\usepackage{mathrsfs}







\newcommand{\Rep}{\mathsf{Rep}}
\newcommand{\Coh}{\mathsf{Coh}}
\newcommand{\Mod}{\mathsf{Mod}}
\newcommand{\Vect}{\mathsf{Vect}}
\newcommand{\Modif}{\mathsf{Modif}}
\newcommand{\HT}{\mathsf{HT}}
\newcommand{\Norm}{\mathsf{Norm}}
\newcommand{\Bun}{\mathsf{Bun}}
\newcommand{\Gra}{\mathsf{Gr}}
\newcommand{\Fil}{\mathsf{Fil}}

\DeclareMathOperator{\inv}{inv}
\DeclareMathOperator{\Spec}{Spec}
\DeclareMathOperator{\Spa}{Spa}
\DeclareMathOperator{\Hom}{Hom}
\DeclareMathOperator{\Ext}{Ext}
\DeclareMathOperator{\coker}{coker}
\DeclareMathOperator{\length}{length}
\DeclareMathOperator{\Tor}{Tor}
\DeclareMathOperator{\rank}{rank}
\DeclareMathOperator{\sk}{sk}
\DeclareMathOperator{\Gr}{Gr}
\DeclareMathOperator{\Gal}{Gal}
\DeclareMathOperator{\Sub}{Sub}
\DeclareMathOperator{\Sym}{Sym}
\DeclareMathOperator{\Jump}{Jump}
\DeclareMathOperator{\HTT}{HT}
\DeclareMathOperator{\loc}{loc}
\DeclareMathOperator{\Frac}{Frac}
\DeclareMathOperator{\Aut}{Aut}
\DeclareMathOperator{\Proj}{Proj}
\DeclareMathOperator{\Exp}{Exp}
\DeclareMathOperator{\Log}{Log}
\DeclareMathOperator{\Pic}{Pic}

\newcommand{\eqd}{\stackrel{\mathrm{def}}{=}}
\newcommand{\pheqd}{\stackrel{\phantom{\mathrm{def}}}{=}}
\newcommand{\sfC}{\mathsf{C}}

\let\oldtilde\tilde
\let\tilde\widetilde

\newcommand*{\dR}{dR}

\graphicspath{{./figures/}}

\newcommand*{\mk}{\mkern -1mu}
\newcommand*{\Mk}{\mkern -2mu}
\newcommand*{\mK}{\mkern 1mu}
\newcommand*{\MK}{\mkern 2mu}

\hypersetup{urlcolor=purple, linkcolor=blue, citecolor=red}

\newcommand*{\romanenumi}{\renewcommand*{\theenumi}{\roman{enumi}}}
\newcommand*{\Romanenumi}{\renewcommand*{\theenumi}{\Roman{enumi}}}
\newcommand*{\alphenumi}{\renewcommand*{\theenumi}{\alph{enumi}}}
\newcommand*{\Alphenumi}{\renewcommand*{\theenumi}{\Alph{enumi}}}
\let\oldforall\forall
\renewcommand*{\forall}{\mathrel{\oldforall}}


\begin{document}
\section*{\raggedright 2690: Uniform convergence and isogeny invariance of renormalized Fargues polygons for finitely presented Breuil–Kisin–Fargues modules over perfectoid extensions of p-adic fields}
\noindent\textbf{Author:} Christophe Cornut; Macarena Peche Irissarry\par\noindent\textbf{Source:} Harder–Narasimhan Filtrations for Breuil–Kisin–Fargues modules. Annales Henri Lebesgue 2 (2019), publisher version, DOI 10.5802/ahl.22. \url{https://ahl.centre-mersenne.org/item/AHL_2019__2__415_0/}
\par\noindent\textbf{License:} CC BY 4.0, \url{https://creativecommons.org/licenses/by/4.0/}. The original publisher PDF has the explicit license notice. See the saved source/license records.
\par\noindent\textbf{Editorial material note (not author text):} The selected problem is Proposition 2.6: uniform convergence, uniform module-dependent Lipschitz control and isogeny invariance of the renormalized Fargues polygons for finitely presented Breuil--Kisin--Fargues modules with M[1/pi] finite free. The exact finite p-adic field, perfectoid extension, ramified Witt ring, Frobenius and xi-prime localization hypotheses are retained; no finite-freeness or HN-type assumption is added. The complete author body retains all local torsion/quasi-abelian categories, devissage, rank/degree, polygon inequalities, convergence and isogeny error arguments. Named p-adic Hodge and category foundations remain explicit prerequisites; later tensor, Hodge, Tate and modification conclusions are unselected context. All proof text and formulas below are editable author TeX. Mathematical wording, language and proof order are retained. Only the article class, title matter and bibliography presentation are adapted for independent compilation; no proof has been generated or repaired.
\par\noindent\textbf{Editorial difficulty:} H2: The author constructs a renormalized Harder–Narasimhan polygon on a non-Noetherian integral p-adic category and proves uniform convergence through quasi-abelian torsion dévissage, polygon inequalities and isogeny error control. This integral p-adic Hodge mechanism exceeds ordinary qualifying algebra and geometry.
\par\noindent\textbf{Editorial retrieval and attribution:} Retrieved 2026-10-01. DOI 10.5802/ahl.22. Source attribution notice: ../licenses/2690-source-license.txt; complete legal text: ../licenses/CC-BY-4.0.txt. Original source excerpts and the existing source-material check are retained in ../sources/2690/.
\medskip\hrule\medskip
\noindent\textbf{Author statement, complete proof and local supporting text follow in source order.}
\section{Introduction}

\subsection{Context}

Cohomology theories provide classifying functors from categories of algebraic varieties to various realisation categories. Grothendieck conjectured that there is a universal such functor, and thus also a universal realisation category, which he called the category of motives. He also worked out an elementary bottom-up construction of this universal functor and its target category, assuming a short list of hard conjectures (the so-called standard conjectures) on which little progress has been made. A top-down approach to Grothendieck's conjecture aims to cut down the elusive category of motives from the various realisation categories of existing cohomology theories, and this first requires assembling them in some ways.

Over an algebraically closed complete extension $C$ of $\mathbb{Q}_{p}$, Bhatt, Morrow and Scholze~\cite{BaMoSc16} have recently defined a new (integral) $p$-adic cohomology theory, which specializes to all other known such theories and nicely explains their relations and pathologies. It takes values in the category of Breuil--Kisin--Fargues modules (hereafter named BKF-modules), a variant of Breuil--Kisin modules due to Fargues~\cite{Fa15}. This new realisation category has various, surprisingly different but nevertheless equivalent incarnations, see~\cite[14.1.1]{ScWe15},~\cite[7.5]{Sc17} or Section~\ref{sec:FunctorsOfFargues}; beyond its obvious relevance for $p$-adic motives, it is also expected to play a role in the reformulation of the $p$-adic Langlands program proposed by Fargues~\cite{Fa16}.

In this paper, we mostly investigate an hidden but implicit structure of these BKF-modules: they are equipped with some sort of Harder--Narasimhan formalism, adapted from either~\cite{Imp16} or~\cite{LeWE16}, which both expanded the original constructions of Fargues~\cite{Fa12} from $p$-divisible groups over $\mathcal{O}_{C}$ to Breuil--Kisin modules.

\subsection{Overview}

In Section~\ref{sec:BKF-Modules}, we define our categories of BKF-modules, review what Bhatt, Morrow and Scholze had to say about them, exhibit the HN-filtrations (which we call Fargues filtrations) and work out their basic properties. In Section~\ref{sec:FunctorsOfFargues}, we turn our attention to the curvy avatar of BKF-modules up to isogenies, namely admissible modifications of vector bundles on the curve, and to their Hodge--Tate realizations. The link between all three incarnations of sthukas with one paw was established by Fargues, according to Scholze who sketched a proof in his lectures at Berkeley\footnote{Between~\cite{ScWe15} and~\cite{BaMoSc16}, the paw was twisted from $A\xi$ to $A\xi'$. We follow the latter convention. No sthukas were harmed in the making of \emph{our} paper, but our valuations have lame normalizations. }. We redo Scholze's proof in slow motion and investigate the Fargues filtration on the curve and Hodge--Tate side, where it tends to be more tractable. We also clarify various issues pertaining to exactness, and introduce some full subcategories where the Fargues filtration is particularly well-behaved. In a subsequent work, we will show that ordinary BKF-modules with $G$-structures factor through these subcategories and compute the corresponding reduction maps, from lattices in the \'etale realization to lattices in the crystalline realization.

\subsection{Results}

We refer to the main body of the paper for all notations.

We define Fargues filtrations $\mathcal{F}_{F}$ and their types $t_{F}$ on $\Mod_{A,t}^{\varphi}$~\eqref{subsec:FarguesOnModAt}, $\Mod_{A,f}^{\varphi,\ast}$~\eqref{subsec:TypeHN}, $\Mod_{\mathcal{O}_{K}^{\flat},f}^{\varphi}$~\eqref{subsec:RisOkflat}, $\Modif_{X}^{ad}$ and $\HT_{E}^{B_{\dR}}$ (\ref{subsec:DefFarguesFiltrOnModif} and~\ref{subsec:HodgeTateModDef}). We show that they are compatible with $\otimes$-product constructions on $\Mod_{A,f}^{\varphi,\ast}$ (Proposition~\ref{prop:catOfHNTypeTensor}), $\Mod_{\mathcal{O}_{K}^{\flat},f}^{\varphi}$ (Proposition~\ref{prop:FarguesonO_Kflattensor}), $\Modif_{X}^{ad}$ and $\HT_{E}^{B_{\dR}}$ (Proposition~\ref{prop:OnHT+AdModif_Farguestensor}). On the isogeny category $\Mod_{A,\ast}^{\varphi}\otimes E$, we only define a type $t_{F,\infty}$~\eqref{subsec:FarguesTypeBKF}, analogous to Fargues's renormalized Harder--Narasimhan function in~\cite{Fa12}. This type matches the Fargues type $t_{F}$ on $\Mod_{A,f}^{\varphi,\ast}\otimes E$ (Proposition~\ref{prop:TypeHN}), and Proposition~\ref{prop:SpecialSubCats} compares it with the Fargues type $t_{F}$ on $\Modif_{X}^{ad}$ and $\HT_{E}^{B_{\dR}}$. We define Hodge filtrations $\mathcal{F}_{H}$ and their types $t_{H}$ on $\Mod_{\mathcal{O}_{K}^{\flat},f}^{\varphi}$~\eqref{subsec:RisOkflat}, $\Mod_{\mathcal{O}_{L},f}^{\varphi}$~\eqref{subsec:RisOL}, $\Mod_{A[\frac{1}{\pi}],f}^{\varphi}$ and $\Mod_{A,\ast}^{\varphi}\otimes E$~\eqref{subsec:RisA=00005B1/pi=00005D}, $\Modif_{X}^{ad}$~\eqref{subsec:AdmModifDef} and $\HT_{E}^{B_{\dR}}$~\eqref{subsec:HodgeTateModDef}. We define opposed Newton filtrations $\mathcal{F}_{N}$ and $\mathcal{F}_{N}^{\iota}$ and their types $t_{N}$ and $t_{N}^{\iota}$ on $\Mod_{L,f}^{\varphi}$~\eqref{subsec:RisL}, and a Newton (or slope) filtration $\mathcal{F}_{N}$ with type $t_{N}$ on $\Bun_{X}$~\eqref{subsec:Newton4BundlesOnCurve} and $\Modif_{X}^{ad}$~\eqref{subsec:AdmModifDef}. The Hodge and Newton filtrations are compatible with $\otimes$-product constructions and satisfy some exactness properties. If $K=C$ is algebraically closed, then for a finite free BKF-module $M\in\Mod_{A,f}^{\varphi}$ mapping to the admissible modification $\underline{\mathcal{E}}\in\Modif_{X}^{ad}$, we establish the following inequalities:
\[
\xymatrix@R=1pc@C=1.5pc{
t_{H}(M\otimes E)\ar@{}[r]|{\geq}\ar@{=}[d] & t_{H}(M\otimes\mathcal{O}_{K}^{\flat})\ar@{}[rr]|{\geq} & & t_{F}(M\otimes\mathcal{O}_{K}^{\flat})\ar@{}[r]|{\geq} & t_{F,\infty}(M)\ar@{:}[dd]_{?}^{?}\\
t_{H}(M\otimes E)\ar@{}[r]|{\geq}\ar@{=}[d] & t_{H}(M\otimes\mathcal{O}_{L})\ar@{}[r]|{\geq} & t_{N}^{\iota}(M\otimes L)\ar@{=}[d]\\
t_{H}(\underline{\mathcal{E}})\ar@{}[rr]|{\geq} & & t_{N}(\underline{\mathcal{E}})\ar@{}[rr]|{\geq} & & t_{F}(\underline{\mathcal{E}}) }
\]
We failed to establish our hope that $t_{F,\infty}(M)=t_{F}(\underline{\mathcal{E}})$ (as did Fargues for $p$-divisible groups in~\cite[Th\'eor\`eme~7]{Fa12} and the second named author for Breuil--Kisin modules in~\cite[Proposition~3.11]{Imp16}), but we nevertheless show in Proposition~\ref{prop:SpecialSubCats} that
\[
\begin{aligned}
t_{F,\infty}(M) & \leq t_{F}(\underline{\mathcal{E}})\quad && \text{if }M\in\Mod_{A,f}^{\varphi,\ast},\\
t_{F,\infty}(M) & \geq t_{F}(\underline{\mathcal{E}}) && \text{if }\underline{\mathcal{E}}\in\Modif_{X}^{ad,\ast}.
\end{aligned}
\]
We also investigate sufficient conditions for the equality $\mathcal{F}_{F}(\underline{\mathcal{E}})=\mathcal{F}_{N}(\underline{\mathcal{E}})$.

\begin{rema}
The definition of the full subcategory $\Mod_{A,f}^{\varphi,\ast}$ of $\Mod_{A,f}^{\varphi}$ is inspired by the notion of $p$-divisible groups of HN-type, due to Fargues~\cite{Fa12}, and expanded to Breuil--Kisin modules in~\cite{Imp16}. The definition of the full subcategory $\Modif_{X}^{ad,\ast}$ of $\Modif_{X}^{ad}$ is new to this paper. We do not know if these subcategories are related under Fargues's equivalence $\underline{\mathcal{E}}:\Mod_{A,f}^{\varphi}\otimes E\stackrel{\simeq}{\longrightarrow}\Modif_{X}^{ad}$ (see Theorem~\ref{thm:FarguesScholze}).
\end{rema}

\subsection{Thanks}

First and foremost, Laurent Fargues, obviously. Then also Matthew Morrow. And Jared Weinstein for his notes, Peter Scholze for his talks.

\subsection{References}

Here are some overall references for the material covered in this paper:~\cite{Co14} or~\cite{Co16} for filtrations, lattices and related invariants, \cite{An09} and~\cite{Co16} for the Harder--Narasimhan formalism, respectively with categories and buildings, \cite{Fa10}, \cite{Fa12} and~\cite{Imp16} or~\cite{LeWE16} for Fargues filtrations on respectively finite flat group schemes, $p$-divisible groups and Breuil--Kisin modules, \cite{FaFo15} and~\cite{Fa15} for the Fargues--Fontaine curve, \cite{ScWe15} and~\cite{BaMoSc16} for BKF-modules, aka~sthukas with one paw. Quasi-abelian categories abound in our paper. These are additive categories with well-behaved kernels and cokernels, which differ from abelian categories by the fact that the canonical morphism from the coimage (cokernel of the kernel) to the image (kernel of the cokernel) may fail to be an isomorphism. They usually show up as the torsion-free part of a cotilting torsion theory on an abelian category~\cite[App.\ B]{BoVdB03}. The Harder--Narasimhan formalism takes as input a quasi-abelian category $\sfC$ equipped with two additive functions $\rank:\sk\sfC\rightarrow\mathbb{N}$ and $\deg:\sk\sfC\rightarrow\mathbb{R}$ subject to various conditions, and outputs a canonical \emph{slope} or \emph{Harder--Narasimhan} filtration on $\sfC$.

\subsection{Notations}

\subsubsection{Types}\label{subsec:Types}

Let $(\Gamma,+,\leq)$ be a totally ordered commutative group. For $r\in\mathbb{N}$, we consider the following submonoid of $\Gamma^{r}$:
\[
\Gamma_{\geq}^{r}\eqd\left\{ (\gamma_{1},\dots,\gamma_{r})\in\Gamma^{r}:\gamma_{1}\geq\dots\geq\gamma_{r}\right\}.
\]
It is equipped with a partial order defined by
\begin{multline*}
(\gamma_{1},\dots,\gamma_{r})\leq(\gamma_{1}^{\prime},\dots,\gamma_{r}^{\prime})\iff\!
\begin{cases}
\forall s\in\{1,\dots,r\}\! & \gamma_{1}+\dots+\gamma_{s}\leq\gamma_{1}^{\prime}+\dots+\gamma_{s}^{\prime},\\
\text{and} & \gamma_{1}+\dots+\gamma_{r}=\gamma_{1}^{\prime}+\dots+\gamma_{r}^{\prime},
\end{cases}\!\!
\end{multline*}
with an involution $\iota:\Gamma_{\geq}^{r}\rightarrow\Gamma_{\geq}^{r}$ and functions $\deg,\max,\min:\Gamma_{\geq}^{r}\rightarrow\Gamma$ defined by
\[
(\gamma_{1},\dots,\gamma_{r})^{\iota}\eqd(-\gamma_{r},\dots,-\gamma_{1}),\qquad
\begin{aligned}
\deg(\gamma_{1},\dots,\gamma_{r}) & \eqd \gamma_{1}+\dots+\gamma_{r},\\
(\gamma_{1},\dots,\gamma_{r})^{\max} & \eqd \gamma_{1},\\
(\gamma_{1},\dots,\gamma_{r})^{\min} & \eqd \gamma_{r}.
\end{aligned}
\]
For $r_{1},r_{2}\in\mathbb{N}$, there is also a ``convex sum'' map
\[
\ast:\Gamma_{\geq}^{r_{1}}\times\Gamma_{\geq}^{r_{2}}\rightarrow\Gamma_{\geq}^{r_{1}+r_{2}}
\]
which concatenates and reorders the elements. We set $\Gamma_{+}:=\{\gamma\in\Gamma:\gamma\geq0\}$ and
\[
\Gamma_{+,\geq}^{\infty}\eqd\underrightarrow{\lim}\,\Gamma_{+,\geq}^{r}=\left\{ (\gamma_{i})_{i=1}^{\infty}\left|
\begin{array}{ll}
\forall i\geq1 & \gamma_{i}\geq\gamma_{i+1}\geq0\\
\forall i\gg1 & \gamma_{i}=0
\end{array}\right.\right\}
\]
where $\Gamma_{+,\geq}^{r}:=\Gamma_{\geq}^{r}\cap\Gamma_{+}^{r}$, with the transition morphisms $\Gamma_{+,\geq}^{r}\hookrightarrow\Gamma_{+,\geq}^{r+1}$ given by $(\gamma_{1},\dots,\gamma_{r})\mapsto(\gamma_{1},\dots,\gamma_{r},0)$. Thus $\Gamma_{+,\geq}^{\infty}$ is yet another partially ordered monoid equipped with a degree function $\deg:\Gamma_{+,\geq}^{\infty}\rightarrow\Gamma_{+}$ and a ``convex sum'' operator
\[
\ast:\Gamma_{+,\geq}^{\infty}\times\Gamma_{+,\geq}^{\infty}\rightarrow\Gamma_{+,\geq}^{\infty}.
\]
If $\Gamma\subset\mathbb{R}$, we will often identify $\Gamma_{\geq}^{r}$ with the monoid of all continuous concave functions $f:[0,r]\rightarrow\mathbb{R}$ such that $f(0)=0$ and $f$ is affine of slope $\gamma_{i}\in\Gamma$ on $[i-1,i]$ for all $i\in\{1,\dots,r\}$. Under this identification, $f\leq f'$ if and only if $f(s)\leq f'(s)$ for all $s\in[0,r]$ with equality for $s=r$, $f^{\iota}(s)=f(r-s)-f(r)$ for all $s\in[0,r]$, $\deg(f)=f(r)$ and finally for $f_{1}\in\Gamma_{\geq}^{r_{1}}$, $f_{2}\in\Gamma_{\geq}^{r_{2}}$ and $s\in[0,r_{1}+r_{2}]$,
\[
f_{1}\ast f_{2}(s)=\max\left\{ f_{1}(s_{1})+f_{2}(s_{2})\:\middle|
\begin{array}{l}
s_{1}\in[0,r_{1}],\,s_{2}\in[0,r_{2}]\\
\text{and }s=s_{1}+s_{2}
\end{array}\!\right\}.
\]
Similarly, we will identify $\Gamma_{+,\geq}^{\infty}$ with the monoid of all continuous concave functions $f:\mathbb{R}_{+}\rightarrow\mathbb{R}_{+}$ such that $f(0)=0$, $f$ is affine of slope $\gamma_{i}\in\Gamma_{+}$ on $[i-1,i]$ for all positive integer $i$, with $\gamma_{i}=0$ for $i\gg0$. Then $f\leq f'$ if and only if $f(s)\leq f'(s)$ for all $s\in\mathbb{R}_{+}$ with equality for $s\gg0$, $\deg(f)=f(s)$ for $s\gg0$ and
\[
f_{1}\ast f_{2}(s)=\max\left\{ f_{1}(t)+f_{2}(s-t)\vert t\in[0,s]\right\}.
\]


\subsubsection{Filtrations}\label{subsec:Filtrations}

In~\cite{Co14}, we defined a notion of $\Gamma$-filtrations for finite free quasi-coherent sheaves (aka vector bundles) over schemes, and in~\cite{Co16} we investigated a notion of $\mathbb{R}$-filtrations on bounded modular lattices of finite length. Here is a common simple framework for $\Gamma$-filtrations and their types. If $(X,\leq)$ is a bounded partially ordered set with smallest element $0_{X}$ and largest element $1_{X}$, then a $\Gamma$-filtration on $X$ is a function $\mathcal{F}:\Gamma\rightarrow X$ which is non-increasing, exhaustive, separated and left-continuous: $\mathcal{F}(\gamma_{1})\geq\mathcal{F}(\gamma_{2})$ for $\gamma_{1}\leq\gamma_{2}$, $\mathcal{F}(\gamma)=1_{X}$ for $\gamma\ll0$, $\mathcal{F}(\gamma)=0_{X}$ for $\gamma\gg0$ and for every $\gamma\in\Gamma$, there is a $\gamma'<\gamma$ such that $\mathcal{F}$ is constant on $]\gamma',\gamma]:=\{\eta\in\Gamma\vert\gamma'<\eta\leq\gamma\}$. If all chains of $X$ are finite, the formula
\[
\mathcal{F}(\gamma)=
\begin{cases}
0_{X} & \text{for }\gamma>\gamma_{1}\\
c_{i} & \text{for }\gamma_{i+1}<\gamma\leq\gamma_{i}\\
1_{X} & \text{for }\gamma\leq\gamma_{s}
\end{cases}
\]
yields a bijection between the set $\mathbf{F}^{\Gamma}(X)$ of all $\Gamma$-filtrations on $X$ and the set of all pairs $(c_{\bullet},\gamma_{\bullet})$ where $c_{\bullet}=\mathcal{F}(\Gamma)=(c_{0}<\dots<c_{s})$ is a (finite) chain of length $s$ in $X$ with $c_{0}=0_{X}$ and $c_{s}=1_{X}$, while $\gamma_{\bullet}=\Jump(\mathcal{F})=(\gamma_{1}>\dots>\gamma_{s})$ is a decreasing sequence in $\Gamma$. We then set $\mathcal{F}_{+}(\gamma):=\max\left\{ \mathcal{F}(\eta):\eta>\gamma\right\} $. If $\rank:X\rightarrow\mathbb{N}$ is an increasing function and $r=\rank(1_{X})$, then all chains of $X$ are finite of length $s\leq r$ and any $\Gamma$-filtration $\mathcal{F}\in\mathbf{F}^{\Gamma}(X)$ has a well-defined type $\mathbf{t}(\mathcal{F})\in\Gamma_{\geq}^{r}$: for any $\gamma\in\Gamma$, the multiplicity of $\gamma$ in $\mathbf{t}(\mathcal{F})$ is equal to $\rank(\mathcal{F}(\gamma))-\rank(\mathcal{F}_{+}(\gamma))$.

If $\sfC$ is an essentially small quasi-abelian category equipped with a rank function $\rank:\sk\,\sfC\rightarrow\mathbb{N}$, as defined in~\cite[3.1]{Co16}, then for every object $X$ of $\sfC$, the partially ordered set $\Sub(X)$ of all strict subobjects of $X$ is a bounded modular lattice of finite length.~A $\Gamma$-filtration on $X$ is then a $\Gamma$-filtration on $\Sub(X)$, and we denote by $\mathbf{F}^{\Gamma}(X)$ the set of all $\Gamma$-filtrations on $X$. For $\mathcal{F}\in\mathbf{F}^{\Gamma}(X)$, we typically write
\[
\mathcal{F}^{\gamma}=\mathcal{F}^{\geq\gamma}=\mathcal{F}(\gamma),\quad\mathcal{F}_{+}^{\gamma}=\mathcal{F}^{>\gamma}=\mathcal{F}_{+}(\gamma)\quad\text{and}\quad\Gr_{\mathcal{F}}^{\gamma}=\mathcal{F}^{\gamma}/\mathcal{F}_{+}^{\gamma}.
\]
If $r=\rank(X)$, the type map $\mathbf{t}:\mathbf{F}^{\Gamma}(X)\rightarrow\Gamma_{\geq}^{r}$ is given by
\[
\mathbf{t}(\mathcal{F})=(\gamma_{1}\geq\dots\geq\gamma_{r})\iff{\forall\gamma\in\Gamma}:\quad\rank\,\Gr_{\mathcal{F}}^{\gamma}=\#\{i:\gamma_{i}=\gamma\}
\]
and the degree map $\deg:\mathbf{F}^{\Gamma}(X)\rightarrow\Gamma$ is given by
\[
\deg(\mathcal{F})=\deg(\mathbf{t}(\mathcal{F}))={\sum_{\gamma\in\Gamma}}\rank\,\Gr_{\mathcal{F}}^{\gamma}\cdot\gamma.
\]
If $0\rightarrow x\rightarrow X\rightarrow y\rightarrow0$ is an exact sequence in $\sfC$, any $\Gamma$-filtration $\mathcal{F}\in\mathbf{F}^{\Gamma}(X)$ induces $\Gamma$-filtrations $\mathcal{F}_{x}\in\mathbf{F}^{\Gamma}(x)$ and $\mathcal{F}_{y}\in\mathbf{F}^{\Gamma}(y)$, and we have
\[
\mathbf{t}(\mathcal{F})=\mathbf{t}(\mathcal{F}_{x})\ast\mathbf{t}(\mathcal{F}_{y})\quad\text{in}\quad\Gamma_{\geq}^{r}.
\]
We denote by $\Gra^{\Gamma}\sfC$ and $\Fil^{\Gamma}\sfC$ the quasi-abelian categories of $\Gamma$-graded and $\Gamma$-filtered objects in $\sfC$. For finite dimensional vector spaces over a field $k$, we set
\[
\Gra_{k}^{\Gamma}\eqd\Gra^{\Gamma}\Vect_{k}\quad\text{and}\quad\Fil_{k}^{\Gamma}\eqd\Fil^{\Gamma}\Vect_{k}.
\]
When $\Gamma=\mathbb{R}$, we simplify our notations to $\mathbf{F}(X):=\mathbf{F}^{\mathbb{R}}(X)$.

\subsubsection{Invariants}\label{subsec:Invariants}

Let $\mathcal{O}$ be a valuation ring with fraction field $K$, maximal ideal $\mathfrak{m}$ and residue field $k$. We denote by $(\Gamma,+,\leq)$ the totally ordered commutative group $(K^{\times}/\mathcal{O}^{\times},\cdot,\leq)$, when we want to view it as an additive group. We extend the total orders to $K/\mathcal{O}^{\times}=K^{\times}/\mathcal{O}^{\times}\cup\{0\}$ and $\Gamma\cup\{-\infty\}$, by declaring that the added elements are smaller than everyone else. We denote by $|{\,\cdot\,}|:K\rightarrow K/\mathcal{O}^{\times}$ the projection. Thus for every $\lambda_{1},\lambda_{2}\in K$, $\left|\lambda_{1}\right|\leq\left|\lambda_{2}\right|\iff\mathcal{O}\lambda_{1}\subset\mathcal{O}\lambda_{2}$. We write
\[
\Exp:\Gamma\cup\{-\infty\}\longleftrightarrow K^{\times}/\mathcal{O}^{\times}\cup\{0\}:\Log
\]
for the corresponding isomorphisms. When the valuation has height $1$, i.e.~when it is given by a genuine absolute value $|{\,\cdot\,}|:K\rightarrow\mathbb{R}_{+}$, we will identify $K^{\times}/\mathcal{O}^{\times}$ with the corresponding subgroup $\left|K^{\times}\right|\subset\mathbb{R}_{+}^{\times}$, and $\Gamma$ with a subgroup of $\mathbb{R}$, using genuine logarithms and exponential maps in some base $b>1$. In this case, we will always normalize the related choices of the absolute value $|{\,\cdot\,}|$ and the base $b$ by requiring that $\log_{b}\left|\pi\right|=-1$ for some specified nonzero element $\pi$ of $\mathfrak{m}$, taking $\pi$ to be a uniformizer whenever $\mathcal{O}$ is a discrete valuation ring. For every element $\gamma\in\Gamma$,
\[
I(\gamma)\eqd\left\{ x\in K:\left|x\right|\leq\Exp(-\gamma)\right\}
\]
is a free, rank one $\mathcal{O}$-submodule of $K$. If $\gamma\in\Gamma_{+}$, it is a principal ideal of $\mathcal{O}$ and
\[
\mathcal{O}(\gamma)\eqd\mathcal{O}/I(\gamma)
\]
is a finitely presented torsion $\mathcal{O}$-module. These modules are the building blocks of the category of finitely presented torsion $\mathcal{O}$-modules, which we denote by $\sfC$.

\begin{lemm}\label{lem:DefInvTorsionOmod}
For any $M\in\sfC$, there is a unique element $(\gamma_{i})_{i=1}^{\infty}$ in $\Gamma_{+,\geq}^{\infty}$ such that $M\simeq\bigoplus_{i=1}^{\infty}\mathcal{O}(\gamma_{i})$. Then $I(\sum_{i=1}^{\infty}\gamma_{i})$ is the Fitting ideal of $M$, $I(\gamma_{i})$ is the annihilator of $\Lambda_{\mathcal{O}}^{i}(M)$ and $\max\{i:\gamma_{i}\neq0\}$ is the minimal number of generators of $M$.
\end{lemm}

\begin{proof}
By~\cite[6.1.14]{GaRa03}, $M\simeq\bigoplus_{i=1}^{r}\mathcal{O}(\gamma_{i})$ for some $r\in\mathbb{N}$ and $\gamma_{1}\geq\dots\geq\gamma_{r}>0$. Plainly, $I(\gamma_{1})$ is the annihilator of $M$, $I(\sum_{i=1}^{r}\gamma_{i})$ is the Fitting ideal of $M$ and $r=\dim_{k}M\otimes_{\mathcal{O}}k$ is the minimal number of generators of $M$. For every $i\geq1$, $\Lambda_{\mathcal{O}}^{i}M\simeq\bigoplus_{I}\mathcal{O}(\gamma_{I})$ where $I$ ranges through the subsets of $\{1,\dots,r\}$ with $i$ elements and $\gamma_{I}:=\min\{\gamma_{i}:i\in I\}$, thus indeed $I(\gamma_{i})$ is the annihilator of $\Lambda_{\mathcal{O}}^{i}M$.
\end{proof}

\begin{defi}
We denote the above invariant by $\inv(M)=(\inv_{i}(M))_{i=1}^{\infty}$ and set
\[
r(M)\eqd\max\{i:\inv_{i}(M)\neq0\},\quad\length(M)\eqd\sum_{i=1}^{\infty}\inv_{i}(M).
\]
\end{defi}

\begin{rema}
When $\mathcal{O}$ is a discrete valuation ring with uniformizer $\pi$, $I(\gamma)=\mathcal{O}\pi^{n}$ for $\gamma=n$ in $\Gamma=\mathbb{Z}$ according to our various conventions, and we thus retrieve the usual invariants of finitely generated torsion $\mathcal{O}$-modules.
\end{rema}

\begin{lemm}\label{lem:InvTorsOModdecUnderSubquo}
Fix $M,N\in\sfC$ and suppose that $N$ is a subquotient of $M$. Then
\[
r(N)\leq r(M)\quad\text{and}\quad\forall i:\;\inv_{i}(N)\leq\inv_{i}(M).
\]
\end{lemm}

\begin{proof}
We just need to establish the second claim when $N$ is either a submodule or a quotient of $M$. For $X\in\sfC$, set $X^{\vee}:=\Hom_{\mathcal{O}}(X,K/\mathcal{O})$. One checks using the previous lemma that this defines an exact duality on $\sfC$, with $\inv(X)=\inv(X^{\vee})$. We may thus even assume that $N$ is a quotient of $M$. Our claim now follows from the previous lemma and the surjectivity of $\Lambda_{\mathcal{O}}^{i}M\twoheadrightarrow\Lambda_{\mathcal{O}}^{i}N$.
\end{proof}

Being a valuation ring, $\mathcal{O}$ is a coherent ring. Thus any finitely generated submodule of a finitely presented $\mathcal{O}$-module is also finitely presented, and the finitely generated submodules of any $M\in\sfC$ also belong to $\sfC$.

\begin{lemm}\label{lem:InvTorsOMod_Formula}
For $M\in\sfC$ and any positive integer $r$,
\[
\sum_{i=1}^{r}\inv_{i}(M)=\max\left\{ \length\left\langle x_{1},\dots,x_{r}\right\rangle :x_{i}\in M\right\}
\]
where $\left\langle x_{1},\dots,x_{r}\right\rangle $ is the $\mathcal{O}$-submodule of $M$ spanned by $\left\{ x_{1},\dots,x_{r}\right\} $.
\end{lemm}
\begin{proof}
It is plainly sufficient to establish that for every submodule $N$ of $M$ generated by $r$ elements, $\length(N)\leq\sum_{i=1}^{r}\inv_{i}(M)$. Now $r(N)\leq r$ by Lemma~\ref{lem:DefInvTorsionOmod}, thus indeed $\length(N)=\sum_{i=1}^{r}\inv_{i}(N)\leq\sum_{i=1}^{r}\inv_{i}(M)$ by Lemma~\ref{lem:InvTorsOModdecUnderSubquo}.
\end{proof}

\begin{lemm}\label{lem:InvTorsOModAndExactSeq}
Let $0\rightarrow M_{1}\rightarrow M_{2}\rightarrow M_{3}\rightarrow0$ be an exact sequence of $\mathcal{O}$-modules. Suppose that two out of $\{M_{1},M_{2},M_{3}\}$ belong to $\sfC.$ Then so does the third one and
\begin{gather*}
\length(M_{1})+\length(M_{3})=\length(M_{2})\quad\text{in}\quad\Gamma_{+},\\
\inv(M_{1})\ast\inv(M_{3})\leq\inv(M_{2})\le\inv(M_{1})+\inv(M_{3})\quad\text{in}\quad\Gamma_{+,\geq}^{\infty}.
\end{gather*}
Moreover, $\inv(M_{1})\ast\inv(M_{3})=\inv(M_{2})$ if and only if the exact sequence splits.
\end{lemm}

\begin{proof}
The first assertion holds for any coherent ring. The additivity of the length comes from~\cite[6.3.1]{GaRa03} and Lemma~\ref{lem:DefInvTorsionOmod}. For $x_{1},\dots,x_{r}\in M_{1}$ and $z_{1},\dots,z_{s}\in M_{3}$, set $y_{i}=x_{i}\in M_{2}$ for $1\leq i\leq r$ and lift $z_{i}\in M_{3}$ to some $y_{r+i}\in M_{2}$ for $1\leq i\leq s$. Then
\begin{align*}
\length\left(\left\langle y_{1},\dots,y_{r+s}\right\rangle \right) & = \length\left(\left\langle y_{1},\dots,y_{r+s}\right\rangle \cap M_{1}\right)+\length\left(\left\langle z_{1},\dots,z_{s}\right\rangle \right)\\
& \geq \length\left(\left\langle x_{1},\dots,x_{r}\right\rangle \right)+\length\left(\left\langle z_{1},\dots,z_{s}\right\rangle \right).
\end{align*}
Lemma~\ref{lem:InvTorsOMod_Formula} now implies that indeed $\inv(M_{1})\ast\inv(M_{3})\leq\inv(M_{2})$. For the second inequality, let $r$ be a positive integer, fix a surjective homomorphism $M_{2}\twoheadrightarrow M_{2}^{\prime}$ where $M_{2}^{\prime}=\bigoplus_{i=1}^{r}\mathcal{O}(\inv_{i}(M_{2}))$, let $M_{1}^{\prime}\subset M_{2}^{\prime}$ be the image of $M_{1}$ and $M_{3}^{\prime}=M_{2}^{\prime}/M_{1}^{\prime}$, so that $M_{i}^{\prime}$ is a finitely presented quotient of $M_{i}$ for $i\in\{1,2,3\}$. Then
\begin{align*}
{\sum_{i=1}^{r}}\inv_{i}(M_{2}) & =\length(M_{2}^{\prime})=\length(M_{1}^{\prime})+\length(M_{3}^{\prime})\\
& ={\sum_{i=1}^{r}}\inv_{i}(M_{1}^{\prime})+{\sum_{i=1}^{r}}\inv_{i}(M_{3}^{\prime})\\
& \leq{\sum_{i=1}^{r}}\inv_{i}(M_{1})+{\sum_{i=1}^{r}}\inv_{i}(M_{3})
\end{align*}
with equality for $r\gg0$, using Lemma~\ref{lem:InvTorsOModdecUnderSubquo} and the aforementioned additivity of $\length$. Therefore indeed $\inv(M_{2})\leq\inv(M_{1})+\inv(M_{3})$ in $\Gamma_{+,\geq}^{\infty}$.

If our exact sequence splits, then plainly $\inv(M_{2})=\inv(M_{1})\ast\inv(M_{3})$. We prove the converse implication by induction on $r(M_{2})$. If $r(M_{2})=0$, there is nothing to prove: $M_{1}=M_{2}=M_{3}=0$. Suppose therefore that $r(M_{2})\geq1$ and $\inv(M_{2})=\inv(M_{1})\ast\inv(M_{3})$, and let $\gamma=\inv_{1}(M_{2})$. Then also $\gamma=\inv_{1}(M_{1})$ or $\gamma=\inv_{1}(M_{3})$. Using the duality $X\mapsto X^{\vee}$ from the proof of Lemma~\ref{lem:InvTorsOModdecUnderSubquo}, we may assume that $\gamma=\inv_{1}(M_{3})$. Write $M_{3}=M_{3}^{\circ}\oplus M_{3}^{\prime}$ with $M_{3}^{\circ}\simeq\mathcal{O}(\gamma)$. This lifts to a splitting $M_{2}=M_{2}^{\circ}\oplus M_{2}^{\prime}$ of the $\mathcal{O}(\gamma)$-module $M_{2}$, with $M_{1}\subset M_{2}^{\prime}$ and $M_{2}^{\circ}\simeq\mathcal{O}(\gamma)$. Since $\inv(M_{i})=\inv(M_{i}^{\circ})\ast\inv(M_{i}^{\prime})$ for $i\in\{2,3\}$, we still have $\inv(M_{2}^{\prime})=\inv(M_{1})\ast\inv(M_{3}^{\prime})$ for the exact sequence $0\rightarrow M_{1}\rightarrow M_{2}^{\prime}\rightarrow M_{3}^{\prime}\rightarrow0$. But $r(M_{2}^{\prime})=r(M_{2})-1$, so this last sequence splits and so does the initial one.
\end{proof}

For every $M\in\sfC$, there is a canonical $\Gamma$-filtration $\mathcal{F}(M)$ on $M\otimes k$ defined by
\[
\mathcal{F}^{\gamma}(M)\eqd
\begin{cases}
\frac{M[I(-\gamma)]+\mathfrak{m}M}{\mathfrak{m}M}\subset\frac{M}{\mathfrak{m}M}=M\otimes k & \text{if }\gamma\leq0,\\
0 & \text{if }\gamma\geq0.
\end{cases}
\]
It depends functorially upon $M$ and one checks easily that we have
\[
\inv(M)=\mathbf{t}^{\iota}(\mathcal{F}(M))\quad\text{in}\quad\Gamma_{+,\geq}^{r}\subset\Gamma_{+,\geq}^{\infty}
\]
where $r=r(M)$. In particular, $\length(M)=-\deg(\mathcal{F}(M))$.

\begin{lemm}\label{lem:InvTorsOModExSeqCaseEqual}
For any exact sequence $0\rightarrow M_{1}\rightarrow M_{2}\rightarrow M_{3}\rightarrow0$ of finitely presented torsion $\mathcal{O}$-modules, the following conditions are equivalent:
\begin{enumerate}
\item \label{1.8.1} The exact sequence splits.
\item \label{1.8.2} For every $\gamma\in\Gamma$, the induced complex of $k$-vector spaces
\[
0\rightarrow\mathcal{F}^{\gamma}(M_{1})\rightarrow\mathcal{F}^{\gamma}(M_{2})\rightarrow\mathcal{F}^{\gamma}(M_{3})\rightarrow0
\]
is a short exact sequence.
\item \label{1.8.3} We have $\inv(M_{2})=\inv(M_{1})\ast\inv(M_{3})$ in $\Gamma_{+,\geq}^{\infty}$.
\end{enumerate}
\end{lemm}

\begin{proof}
Plainly $\eqref{1.8.1}\Rightarrow\eqref{1.8.2}$ and $\eqref{1.8.3}\Rightarrow\eqref{1.8.1}$ by Lemma~\ref{lem:InvTorsOModAndExactSeq}. If \eqref{1.8.2} holds, then
\[
0\rightarrow\left(M_{1}\otimes k,\mathcal{F}(M_{1})\right)\rightarrow\left(M_{2}\otimes k,\mathcal{F}(M_{2})\right)\rightarrow\left(M_{3}\otimes k,\mathcal{F}(M_{3})\right)\rightarrow0
\]
is an exact sequence of $\Gamma$-filtered $k$-vector spaces, thus
\[
\mathbf{t}(\mathcal{F}(M_{2}))=\mathbf{t}(\mathcal{F}(M_{1}))\ast\mathbf{t}(\mathcal{F}(M_{3}))\quad\text{in}\quad\Gamma_{\geq}^{r}
\]
where $r=r(M_{2})=r(M_{1})+r(M_{3})$, therefore
\[
\mathbf{t}^{\iota}(\mathcal{F}(M_{2}))=\mathbf{t}^{\iota}(\mathcal{F}(M_{1}))\ast\mathbf{t}^{\iota}(\mathcal{F}(M_{3}))\quad\text{in}\quad\Gamma_{\geq}^{r}
\]
from which \eqref{1.8.3} immediately follows.
\end{proof}

\subsubsection{Lattices}\label{subsec:Lattices}

An $\mathcal{O}$-lattice in a finite dimensional $K$-vector space $V$ is a finitely generated $\mathcal{O}$-submodule $L$ of $V$ spanning $V$ over $K$. Any such $L$ is finite free over $\mathcal{O}$ by~\cite[VI, \S 3, 6, Lemme~1]{BoAC56}. We denote by $\mathcal{L}(V)$ the set of all $\mathcal{O}$-lattices in $V$. Since $\mathcal{O}$ is an elementary divisor ring~\cite[\S 10]{Ka66}, for every $L_{1},L_{2}\in\mathcal{L}(V)$, there is an $\mathcal{O}$-basis $(e_{1},\dots,e_{r})$ of $L_{1}$ and elements $(x_{1},\dots,x_{r})$ of $K^{\times}$ such that $(x_{1}e_{1},\dots,x_{r}e_{r})$ is an $\mathcal{O}$-basis of $L_{2}$ and $\left|x_{1}\right|\geq\dots\geq\left|x_{r}\right|$ (we say that the basis is adapted to $L_{1}$ and $L_{2}$). If $\gamma_{i}=\Log\left|x_{i}\right|$, then $(\gamma_{1},\dots,\gamma_{r})$ belongs to $\Gamma_{\geq}^{r}$ and does not depend upon the chosen basis. Indeed, one checks using the given adapted basis that the formula
\[
\mathcal{F}^{\gamma}(L_{1},L_{2})\eqd\frac{L_{1}\cap I(\gamma)L_{2}+\mathfrak{m}L_{1}}{\mathfrak{m}L_{1}}\subset\frac{L_{1}}{\mathfrak{m}L_{1}}=L_{1}\otimes k
\]
defines a $\Gamma$-filtration $\mathcal{F}(L_{1},L_{2})$ on $L_{1}\otimes k$, whose type $\mathbf{d}(L_{1},L_{2})\in\Gamma_{\geq}^{r}$ equals $(\gamma_{1},\dots,\gamma_{r})$. In particular, $L_{1}=L_{2}$ if and only if $\mathbf{d}(L_{1},L_{2})=0$. This computation also shows that $\mathbf{d}(L_{2},L_{1})=\mathbf{d}^{\iota}(L_{1},L_{2})$ in $\Gamma_{\geq}^{r}$. If $L_{1}\subset L_{2}$, then $Q=L_{2}/L_{1}$ is a finitely presented torsion $\mathcal{O}$-module, $\mathbf{d}(L_{1},L_{2})\in\Gamma_{+,\geq}^{r}$ and $\mathbf{d}(L_{1},L_{2})=\inv(Q)$ in $\Gamma_{+,\geq}^{\infty}$. If moreover $L_{1}\subset\mathfrak{m}L_{2}$ (i.e.~$\inv_{r}(Q)\neq0$), the projection $L_{2}\twoheadrightarrow Q$ induces an isomorphism $L_{2}\otimes k\simeq Q\otimes k$ mapping $\mathcal{F}(L_{2},L_{1})$ to $\mathcal{F}(Q)$.

\begin{lemm}\label{lem:triangIneq}
For $L_{1},L_{2},L_{3}\in\mathcal{L}(V)$, we have the following triangular inequality:
\[
\mathbf{d}(L_{1},L_{3})\leq\mathbf{d}(L_{1},L_{2})+\mathbf{d}(L_{2},L_{3})\quad\text{in}\quad\Gamma_{\geq}^{r}.
\]
\end{lemm}

\begin{proof}
For any $x\in K^{\times}$ and $L,L'\in\mathcal{L}(V)$, if $\gamma=\Log\left|x\right|$, then
\[
\mathbf{d}(x^{-1}L,L')=\mathbf{d}(L,xL')=\mathbf{d}(L,L')+(\gamma,\dots,\gamma)\quad\text{in}\quad\Gamma_{\geq}^{r}.
\]
Changing $(L_{1},L_{2},L_{3})$ to $(xL_{1},L_{2},x^{-1}L_{3})$ for a suitable~$x$, we may thus assume that $L_{1}\subset L_{2}\subset L_{3}$. Applying Lemma~\ref{lem:InvTorsOModAndExactSeq} to the exact sequence
\[
0\rightarrow L_{2}/L_{1}\rightarrow L_{3}/L_{1}\rightarrow L_{3}/L_{2}\rightarrow0
\]
we obtain the desired inequality.
\end{proof}

\begin{rema}
When $\Gamma\hookrightarrow\mathbb{R}$, the previous lemma also follows from~\cite[5.2.8 \& 6.1]{Co14}.
\end{rema}

\begin{lemm}\label{lem:RelPos=000026ExcSeq}
Let $0\rightarrow V_{1}\rightarrow V_{2}\rightarrow V_{3}\rightarrow0$ be an exact sequence of $K$-vector spaces. For any pair of $\mathcal{O}$-lattices $L_{2},L_{2}^{\prime}\in\mathcal{L}(V_{2})$, their inverse and direct images in $V_{1}$ and $V_{3}$ are $\mathcal{O}$-lattices $L_{1},L'_{1}\in\mathcal{L}(V_{1})$ and $L_{3},L'_{3}\in\mathcal{L}(V_{3})$, and we have
\[
\mathbf{d}(L_{2},L'_{2})\geq\mathbf{d}(L_{1},L'_{1})\ast\mathbf{d}(L_{3},L'_{3})\quad\text{in}\quad\Gamma_{\geq}^{r_{2}}
\]
where $r_{i}=\dim_{K}V_{i}$, with equality if and only if for every $\gamma\in\Gamma$,
\[
0\rightarrow\mathcal{F}^{\gamma}(L_{1},L_{1}^{\prime})\rightarrow\mathcal{F}^{\gamma}(L_{2},L_{2}^{\prime})\rightarrow\mathcal{F}^{\gamma}(L_{3},L_{3}^{\prime})\rightarrow0
\]
is an exact sequence.
\end{lemm}
\begin{proof}
Plainly $0\rightarrow L_{1}\rightarrow L_{2}\rightarrow L_{3}\rightarrow0$ and $0\rightarrow L_{1}^{\prime}\rightarrow L_{2}^{\prime}\rightarrow L_{3}^{\prime}\rightarrow0$ are exact; thus $L_{3}$ and $L_{3}^{\prime}$ are finitely generated over $\mathcal{O}$, in particular they are both $\mathcal{O}$-lattices in $V_{3}$ and free over $\mathcal{O}$; it follows that both exact sequences split, which implies that $L_{1}$ and $L_{1}^{\prime}$ are also (finite free) $\mathcal{O}$-lattices in $V_{1}$. For the remaining claims, we may as above replace $L_{2}^{\prime}$ by $xL_{2}^{\prime}$ for some $x\in K^{\times}$ (which replaces $L_{i}^{\prime}$ by $xL_{i}^{\prime}$ for $i\in\{1,3\}$) to reduce to the case where $L_{i}^{\prime}\subset\mathfrak{m}L_{i}\subset L_{i}$ for all $i\in\{1,2,3\}$. Applying Lemma~\ref{lem:InvTorsOModAndExactSeq} to the resulting exact sequence of finitely presented torsion $\mathcal{O}$-modules
\[
0\rightarrow L_{1}/L_{1}^{\prime}\rightarrow L_{2}/L_{2}^{\prime}\rightarrow L_{3}/L_{3}^{\prime}\rightarrow0
\]
we obtain the inequality $\mathbf{d}^{\iota}(L_{2},L'_{2})\geq\mathbf{d}^{\iota}(L_{1},L'_{1})\ast\mathbf{d}^{\iota}(L_{3},L'_{3})$ in $\Gamma_{\geq}^{r_{2}}$, which is equivalent to the desired inequality $\mathbf{d}(L_{2},L'_{2})\geq\mathbf{d}(L_{1},L'_{1})\ast\mathbf{d}(L_{3},L'_{3})$. Moreover by Lemma~\ref{lem:InvTorsOModExSeqCaseEqual}, equality holds in either one of them if and only if for every $\gamma\in\Gamma$,
\[
0\rightarrow\mathcal{F}^{\gamma}(L_{1},L_{1}^{\prime})\rightarrow\mathcal{F}^{\gamma}(L_{2},L_{2}^{\prime})\rightarrow\mathcal{F}^{\gamma}(L_{3},L_{3}^{\prime})\rightarrow0
\]
is an exact sequence of $k$-vector spaces. This proves the lemma.
\end{proof}

\begin{rema}
When $\Gamma\hookrightarrow\mathbb{R}$, the inequality also follows from~\cite[5.2.10 \& 6.1]{Co14}.
\end{rema}

For $L_{1},L_{2}\in\mathcal{L}(V)$, we denote by $\nu(L_{1},L_{2})\in\Gamma$ the degree of $\mathbf{d}(L_{1},L_{2})$.

\subsubsection{Tensor products}\label{subsec:TensorProd}

There are also compatible notions of tensor products, symmetric and exterior powers for types, objects and $\Gamma$-filtered objects in arbitrary quasi-tannakian categories, and $\mathcal{O}$-lattices in $K$-vector spaces. All of these notions are fairly classical, and their various compatibilities easily checked. For instance if $L$ and $L^{\prime}$ are $\mathcal{O}$-lattices in $V$ and $i\in\mathbb{N}$, then $\Lambda_{\mathcal{O}}^{i}L$ and $\Lambda_{\mathcal{O}}^{i}L'$ are $\mathcal{O}$-lattices in $\Lambda_{K}^{i}V$, $\mathcal{F}(\Lambda_{\mathcal{O}}^{i}L,\Lambda_{\mathcal{O}}^{i}L')$ is the $\Gamma$-filtration $\Lambda_{k}^{i}\mathcal{F}(L,L')$ on $\Lambda_{k}^{i}(L\otimes_{\mathcal{O}}k)$ which is the image of the $\Gamma$-filtration $\mathcal{F}(L,L')^{\otimes i}$ on $(L\otimes_{\mathcal{O}}k)^{\otimes i}$ under the projection $(L\otimes_{\mathcal{O}}k)^{\otimes i}\twoheadrightarrow\Lambda_{k}^{i}(L\otimes_{\mathcal{O}}k)$, where $\mathcal{F}(L,L')^{\otimes i}(\gamma)=\sum_{\gamma_{1}+\dots+\gamma_{i}=\gamma}\mathcal{F}(L,L')(\gamma_{1})\otimes\dots\otimes\mathcal{F}(L,L')(\gamma_{i})$ in $(L\otimes_{\mathcal{O}}k)^{\otimes i}$ for every $\gamma\in\Gamma$. The type $\mathbf{d}(\Lambda_{\mathcal{O}}^{i}L,\Lambda_{\mathcal{O}}^{i}L')=\mathbf{t}(\Lambda_{k}^{i}\mathcal{F}(L,L'))=\Lambda^{i}\mathbf{d}(L,L')$ is obtained from $\mathbf{d}(L,L')=(\gamma_{1},\dots,\gamma_{r})$ (with $r=\dim_{K}V$) by reordering the elements $\gamma_{I}=\sum_{j\in I}\gamma_{j}$ where $I$ ranges through all subsets of $\{1,\dots,r\}$ of cardinality $i$.


\section{Breuil--Kisin--Fargues Modules}\label{sec:BKF-Modules}

\subsection{The rings}\label{subsec:TheRings}

Let $p$ be a prime number, $E$ a finite extension of $\mathbb{Q}_{p}$, $K$ a perfectoid field extension of $E$, $K^{\flat}$ the tilt of $K$. We denote by $\mathcal{O}_{E}$, $\mathcal{O}_{K}$ and $\mathcal{O}_{K}^{\flat}$ the ring of integers in $E$, $K$ and $K^{\flat}$, with maximal ideals $\mathfrak{m}_{E}$, $\mathfrak{m}_{K}$ and $\mathfrak{m}_{K}^{\flat}$, and perfect residue fields $\mathbb{F}_{q}:=\mathcal{O}_{E}/\mathfrak{m}_{E}$ (finite with $q$ elements) and $\mathbb{F}:=\mathcal{O}_{K}/\mathfrak{m}_{K}=\mathcal{O}_{K}^{\flat}/\mathfrak{m}_{K}^{\flat}$. We fix once and for all a uniformizer $\pi$ of $E$. We denote by $W_{\mathcal{O}_{E}}({\,\cdot\,})$ the Witt vector functor with values in $\mathcal{O}_{E}$-algebras, as defined in~\cite[1.2]{FaFo15}. We set
\[
A(\mathcal{O}_{K})\eqd W_{\mathcal{O}_{E}}(\mathcal{O}_{K}^{\flat}),\quad A(K)\eqd W_{\mathcal{O}_{E}}(K^{\flat}),\quad\mathcal{O}_{L}\eqd W_{\mathcal{O}_{E}}(\mathbb{F}),\quad L\eqd\Frac(\mathcal{O}_{L}).
\]
Thus $A(K)$ and $\mathcal{O}_{L}$ are complete discrete valuation rings with uniformizer $\pi$ and residue fields respectively equal to $K^{\flat}$ and $\mathbb{F}$, while our main player $A:=A(\mathcal{O}_{K})$ is a non-noetherian complete local ring with maximal ideal $\mathfrak{m}$ and residue field $\mathbb{F}$. We denote by $\varphi$ the Frobenius $x\mapsto x^{q}$ in characteristic $p$ or its extension to $\mathcal{O}_{E}$-Witt vectors. The ring homomorphisms $\mathcal{O}_{K}^{\flat}\hookrightarrow K^{\flat}$ and $\mathcal{O}_{K}^{\flat}\twoheadrightarrow\mathbb{F}$ induce $\varphi$-equivariant homomorphisms of $\mathcal{O}_{E}$-algebras $A\hookrightarrow A(K)$ and $A\twoheadrightarrow\mathcal{O}_{L}$. The formula $\theta(\sum_{n\geq0}[(x_{i,n})_{i}]\pi^{n})=\sum x_{0,n}\pi^{n}$ defines a surjective ring homomorphism $\theta:A\twoheadrightarrow\mathcal{O}_{K}$ whose kernel is a principal ideal. Here $[-]$ is the Teichm\"uller lift and $(x_{i,n})_{i\geq0}\in\mathcal{O}_{K}^{\flat}$ for all $n\geq0$, i.e.~$x_{i,n}\in\mathcal{O}_{K}$ with $x_{i,n}=x_{i+1,n}^{p}$ for all $i\geq0$. We fix a generator $\xi$ of $\ker(\theta)$ and set $\xi':=\varphi(\xi)$. We write $\varpi$ for the image of $\xi$ in $\mathcal{O}_{K}^{\flat}=A_{1}$, where more generally $A_{n}:=A/\pi^{n}A$ for $n\in\mathbb{N}$. Thus $\varpi$ is a pseudo-uniformizer of $K^{\flat}$, i.e.~a non-zero element of $\mathfrak{m}_{K}^{\flat}$. For an $A$-module $M$ and $n\in\mathbb{N}$, we define\emph{ }
\[
M(K)\eqd M\otimes_{A}A(K),\quad M(\mathcal{O}_{L})\eqd M\otimes_{A}\mathcal{O}_{L},\quad M_{n}\eqd M\otimes_{A}A_{n}=M/\pi^{n}M.
\]
In particular, $M_{1}=M\otimes_{A}\mathcal{O}_{K}^{\flat}$. We normalize the absolute value of $K^{\flat}$ by requiring that $q\left|\varpi^{q}\right|=1$. This weird normalization is meant to simplify some formulas, under the conventions borrowed from~\cite{BaMoSc16} for sthukas with one paw, which differ from those of~\cite{ScWe15}; the latter would have lead us to the normalization $q\left|\varpi\right|=1$. We normalize the discrete absolute value on $E$ and $L$ by requiring that $q\left|\pi\right|=1$, and we similarly normalize the discrete absolute values on the fraction fields $B_{\dR}$ and $B'_{\dR}$ of the completed local rings $B_{\dR}^{+}$ and $B_{\dR}^{\prime+}$ of $A$ at $(\xi)$ and $(\xi')$ by requiring that respectively $q\left|\xi\right|=1$ and $q\left|\xi'\right|=1$.
\begin{rema}
In the sequel, we will often refer to~\cite{BaMoSc16} or~\cite{ScWe15} where $E=\mathbb{Q}_{p}$. We have carefully checked that the results we quote from either of these references also hold when $E$ is an arbitrary finite extension of $\mathbb{Q}_{p}$, with essentially identical proofs.
\end{rema}

\subsection{Categories of \texorpdfstring{$A$}{A}-modules}

\subsubsection{}\label{subsec:ProjDimModA}

For an $A$-module $M$, we denote by $\tilde{M}$ the corresponding quasi-coherent sheaf on $X:=\Spec\,A$. Since $U:=X\setminus\{\mathfrak{m}\}$ is a quasi-compact open subscheme of the affine scheme $X$, there is an exact sequence~\cite[II Corollaire~4]{SGA2r} of $A$-modules
\[
0\rightarrow H_{\{\mathfrak{m}\}}^{0}(X,\tilde{M})\rightarrow M=H^{0}(X,\tilde{M})\rightarrow H^{0}(U,\tilde{M})\rightarrow H_{\{\mathfrak{m}\}}^{1}(X,\tilde{M})\rightarrow0
\]
and for every $i\geq1$, an isomorphism of $A$-modules
\[
H^{i}(U,\tilde{M})\simeq H_{\{\mathfrak{m}\}}^{i+1}(X,\tilde{M}).
\]
Moreover for any sequence of parameters $(a,b)$ spanning an ideal $I$ with $\sqrt{I}=\mathfrak{m}$,
\[
H_{\{\mathfrak{m}\}}^{i}(X,\tilde{M})\simeq H^{i}\left(\left[M\rightarrow M\left[{\frac{1}{a}}\right]\oplus M\left[{\frac{1}{b}}\right]\rightarrow M\left[{\frac{1}{ab}}\right]\right]\right)
\]
by~\cite[II Proposition~5]{SGA2r}, thus $H_{\{\mathfrak{m}\}}^{i}(X,\tilde{M})=H^{i-1}(U,\tilde{M})=0$ for $i\geq3$. Also,
\[
H_{\{\mathfrak{m}\}}^{i}(X,\tilde{M})=\underrightarrow{\lim}\,\Ext_{A}^{i}\left(A/I^{n},M\right)
\]
for any $i\geq0$ if moreover $(a,b)$ is regular by~\cite[II Lemme~9]{SGA2r}. For $M=A$, we find
\[
H_{\{\mathfrak{m}\}}^{i}\left(X,\mathcal{O}_{X}\right)=
\begin{cases}
0 & \text{if }i\neq2\\
\mathcal{E} & \text{if }i=2
\end{cases}\quad\text{with}\quad\mathcal{E}=\frac{A\left[{\frac{1}{\pi[\varpi]}}\right]}{A\left[{\frac{1}{\pi}}\right]+A\left[{\frac{1}{[\varpi]}}\right]}\neq0
\]
using~\cite[4.6]{BaMoSc16} for $i=1$. By~\cite[2.6 \& 2.7]{HaMa07} and with the definition given there,
\[
\text{p-depth}_{A}(M)=\sup\left\{ k\geq0:H_{\{\mathfrak{m}\}}^{i}\left(X,\tilde{M}\right)=0\text{ for all }i<k\right\}.
\]
In particular $\text{p-depth}_{A}(A)=2$. We say that the $A$-module $M$ is perfect if it has a finite resolution by finite free $A$-modules. The Auslander--Buchsbaum theorem of~\cite[Chapter~6, Theorem~2]{No76} then assert that for any such $M$,
\[
\text{proj}.\text{dim}_{A}(M)+\text{p-depth}_{A}(M)=2.
\]
In particular, $\text{proj}.\text{dim}_{A}(M)\leq1$ if and only if the $A$-submodule
\[
M[\mathfrak{m}^{\infty}]\eqd H_{\{\mathfrak{m}\}}^{0}(X,\tilde{M})
\]
of $M=H^{0}(X,\tilde{M})$ is trivial, and $M$ is finite free if and only if moreover
\[
H_{\{\mathfrak{m}\}}^{1}(X,\tilde{M})=\coker\left(M\rightarrow H^{1}(U,\tilde{M})\right)
\]
is trivial.

\subsubsection{}\label{subsec:basicsonModA*}

We denote by $\Mod_{A}$ the abelian category of all $A$-modules. Let $\Mod_{A,\ast}$ be the strictly full subcategory of finitely presented $A$-modules $M$ such that $M[\frac{1}{\pi}]$ is a projective $A[\frac{1}{\pi}]$-module. Any such $M$ is a perfect $A$-module and $M[\frac{1}{\pi}]$ is actually finite and free over $A[\frac{1}{\pi}]$ by~\cite[4.9 \& 4.12]{BaMoSc16}. By~\cite[4.13]{BaMoSc16}, the $A$-dual $M^{\vee}:=\Hom_{A}(M,A)$ is finite free over $A$, so is the bidual $M_{f}:=M^{\vee\vee}$, the kernel of the canonical morphism $M\rightarrow M_{f}$ is the torsion submodule $M[\pi^{\infty}]$ of $M$, it is a finitely presented $A$-module killed by $\pi^{n}$ for $n\gg0$, and the cokernel of $M\rightarrow M_{f}$ is a finitely presented torsion $A$-module $\overline{M}$ supported at $\mathfrak{m}$. We claim that $M[\mathfrak{m}^{\infty}]$ is then also a finitely presented $A$-module (supported at $\mathfrak{m}$). To see this, note that
\[
M[\mathfrak{m}^{\infty}]=\ker\left(M[\pi^{\infty}]\rightarrow M[\pi^{\infty}]\!\left[{\frac{1}{[\varpi]}}\right]\right).
\]
Since $M[\pi^{\infty}][{\frac{1}{[\varpi]}}]\simeq M[\pi^{\infty}](K)$ is a finitely generated torsion module over the complete discrete valuation ring $A(K)$, there is a unique sequence of integers
\[
\inv_{A(K)}\left(M[\pi^{\infty}](K)\right)\eqd(n_{1}\geq\dots\geq n_{r}>0)\quad\text{in}\quad\mathbb{N}_{\geq}^{r}
\]
for some $r\in\mathbb{N}$ such that $M[\pi^{\infty}][{
\frac{1}{[\varpi]}}]$ is isomorphic to
\[
\bigoplus_{i=1}^{r}A_{n_{i}}(K)=\bigoplus_{i=1}^{r}A_{n_{i}}\!\left[{\frac{1}{[\varpi]}}\right]\!.
\]
Chasing denominators, we may modify any such isomorphism into one that fits in a commutative diagram of $\pi^{\infty}$-torsion $A$-modules
\[
\xymatrix@C=2pc@R=1pc{
\bigoplus_{i=1}^{r}A_{n_{i}} \ar@{^{(}->}[r] \ar@{}[d]|{\bigcap}& M[\pi^{\infty}]\ar[d]\\
\bigoplus_{i=1}^{r}A_{n_{i}}[{\frac{1}{[\varpi]}]} \ar[r]^{\simeq} & M[\pi^{\infty}][{\frac{1}{[\varpi]}]}
}
\]
If $\pi^{n}M[\pi^{\infty}]=0$, the cokernel of the top map is a finitely generated $A_{n}$-module $Q$ with $Q[\frac{1}{[\varpi]}]=0$, thus $[\varpi]^{m}Q=0$ for $m\gg0$. Since $(\bigoplus_{i=1}^{r}A_{n_{i}})[\mathfrak{m}^{\infty}]=0$, $M[\mathfrak{m}^{\infty}]$ embeds into $Q$, therefore also $[\varpi]^{m}M[\mathfrak{m}^{\infty}]=0$, i.e.~$M[\mathfrak{m}^{\infty}]$ is the kernel of $[\varpi]^{m}$ acting on the finitely presented $A_{n}$-module $M[\pi^{\infty}]$. It follows that $M[\mathfrak{m}^{\infty}]$ is itself finitely presented over $A_{n}$ and $A$, since $A_{n}$ is a coherent ring by (the easy case of) \cite[Proposition~3.24]{BaMoSc16}. We finally define the subquotient
\[
M_{t}\eqd M[\pi^{\infty}]/M[\mathfrak{m}^{\infty}].
\]
This is a finitely presented $A$-module killed by $\pi^{n}$ for $n\gg0$.

\subsubsection{}\label{subsec:DevissageOfModules}

We will consider the following strictly full subcategories of $\Mod_{A,\ast}$:
\begin{align*}
\Mod_{A,f} & \eqd \left\{ \text{finite free \ensuremath{A}-modules}\right\},\\
\Mod_{A,\pi^{\infty}} & \eqd \left\{ \text{finitely presented \ensuremath{A}-modules killed by \ensuremath{\pi^{n}\,}for \ensuremath{n\gg0}}\right\} \\
& \pheqd \left\{ M\in\Mod_{A,\ast}\text{ such that }M=M[\pi^{\infty}]\right\},\\
\Mod_{A,\mathfrak{m}^{\infty}} & \eqd \left\{ \text{finitely presented \ensuremath{A}-modules killed by \ensuremath{(\pi,[\varpi])^{n}\,}for \ensuremath{n\gg0}}\right\} \\
&\pheqd\left\{M\in\Mod_{A,\ast}\text{ such that }M=M[\mathfrak{m}^{\infty}]\right\},\\
\Mod_{A,t} & \eqd \left\{ \text{finitely presented \ensuremath{A}-modules with \ensuremath{\pi\,}nilpotent and \ensuremath{[\varpi]\,}injective}\right\} \\
& \pheqd \left\{ M\in\Mod_{A,\ast}\text{ such that \ensuremath{M=M_{t}}}\right\}.
\end{align*}
Then any $M\in\Mod_{A,\ast}$ has a canonical and functorial \emph{d\'evissage
\[
\xymatrix@R=1pc{ & M[\mathfrak{m}^{\infty}]\ar@{^{(}->}[d]\\
0\ar[r] & M[\pi^{\infty}]\ar[r]\ar@{->>}[d] & M\ar[r] & M_{f}\ar[r] & \overline{M}\ar[r] & 0\\
& M_{t} }
\]
}with everyone in the relevant subcategory. The projective dimension of the nonzero $A$-modules in $\Mod_{A,f}$, $\Mod_{A,t}$ and $\Mod_{A,\mathfrak{m}^{\infty}}$ are respectively $0$, $1$ and $2$.

\subsubsection{}\label{subsec:noncanseq}

For $n\gg0$, $\pi^{n}$ kills $M[\pi^{\infty}]$ and $\overline{M}$, thus for any $m\in M_{f}$, $\pi^{n}m$ is the image of some $m'\in M$ and $\pi^{n}m'\in M$ only depends upon $m$. This defines an embedding $M_{f}\hookrightarrow M$ whose cokernel $Q$ is a finitely presented $A$-module killed by $\pi^{2n}$:
\[
0\rightarrow M_{f}\rightarrow M\rightarrow Q\rightarrow0.
\]
By construction, the composition $M_{f}\rightarrow M\rightarrow M_{f}$ is multiplication by $\pi^{2n}$ and since $M[\pi^{2n}]=M[\pi^{\infty}]$, we obtain an exact sequence of $A$-modules
\[
0\rightarrow M[\pi^{\infty}]\rightarrow Q\rightarrow M_{f,2n}\rightarrow\overline{M}\rightarrow0.
\]
Thus if $M[\mathfrak{m}^{\infty}]=0$, then also $Q[\mathfrak{m}^{\infty}]=0$, i.e.~$Q\in\Mod_{A,t}$.

\subsubsection{}\label{subsec:FiltOnModAt=000026ModAm}

Any $A$-module $M$ in $\Mod_{A,\pi^{\infty}}$ has yet another canonical and functorial \emph{d\'evissage}, the finite non-decreasing filtration by the finitely presented $A$-submodules $M[\pi^{n}]$ of $M$ whose successive quotients $M[\pi^{n}]/M[\pi^{n-1}]\simeq\pi^{n-1}M[\pi^{n}]$ are finitely presented $\mathcal{O}_{K}^{\flat}$-modules. If $M$ belongs to $\Mod_{A,t}$, these subquotients are torsion free, thus finite free over $\mathcal{O}_{K}^{\flat}$. If $M$ belongs to $\Mod_{A,\mathfrak{m}^{\infty}}$, they are finitely presented torsion $\mathcal{O}_{K}^{\flat}$-modules, thus themselves non-canonically isomorphic to direct sums of modules of the form $\mathcal{O}_{K}^{\flat}(x):=\mathcal{O}_{K}^{\flat}/x\mathcal{O}_{K}^{\flat}$ with $x$ nonzero in $\mathcal{O}_{K}^{\flat}$.

\subsubsection{}\label{subsec:Tor(A(K)LO_L)}


For every $A$-module $N$ and any nonzero $x\in\mathcal{O}_{K}^{\flat}$, the exact sequences
\[
0\rightarrow A\stackrel{\pi}{\rightarrow}A\rightarrow\mathcal{O}_{K}^{\flat}\rightarrow0\quad\text{and}\quad0\rightarrow\mathcal{O}_{K}^{\flat}\stackrel{x}{\rightarrow}\mathcal{O}_{K}^{\flat}\rightarrow\mathcal{O}_{K}^{\flat}(x)\rightarrow0
\]
give $\Tor_{0}^{A}(N,\mathcal{O}_{K}^{\flat})=N/\pi N$, $\Tor_{1}^{A}(N,\mathcal{O}_{K}^{\flat})=N[\pi]$, and an exact sequence
\[
0\rightarrow N[\pi]/xN[\pi]\rightarrow\Tor_{1}^{A}(N,\mathcal{O}_{K}^{\flat}(x))\rightarrow(N/\pi N)[x]\rightarrow0.
\]
It then follows from~\ref{subsec:noncanseq} and~\ref{subsec:FiltOnModAt=000026ModAm} that for every $M\in\Mod_{A,\ast}$,
\[
\Tor_{1}^{A}(A(K),M)=0\quad\text{ and}\quad\Tor_{1}^{A}(L,M)=0
\]
since this holds for $M\in\{A,\mathcal{O}_{K}^{\flat},\mathcal{O}_{K}^{\flat}(x)\}$. If moreover $M[\mathfrak{m}^{\infty}]=0$, then also
\[
\Tor_{1}^{A}(\mathcal{O}_{L},M)=0
\]
since this holds for $M\in\{A,\mathcal{O}_{K}^{\flat}\}$.

\subsubsection{}

The category $\Mod_{A,\ast}$ is stable under extensions in $\Mod_{A}$. The next proposition implies that it inherits from $\Mod_{A}$ the structure of a closed symmetric monoidal category, which just says that~$\Mod_{A,\ast}$ is a $\otimes$-category with internal Homs.
\begin{prop}
For every $M_{1}$ and $M_{2}$ in $\Mod_{A,\ast}$ and any $i\geq0$,
\[
\Tor_{i}^{A}(M_{1},M_{2})\quad\text{and}\quad\Ext_{A}^{i}(M_{1},M_{2})
\]
also belong to $\Mod_{A,\ast}$.
\end{prop}
\begin{proof}
Fix a finite resolution $P_{\bullet}$ of $M_{1}$ by finite free $A$-modules. Then
\[
\Ext_{A}^{i}(M_{1},M_{2})=H^{i}\left(\Hom_{A}(P_{\bullet},M_{2})\right)\quad\text{and}\quad\Tor_{i}^{A}(M_{1},M_{2})=H^{i}\left(P_{\bullet}\otimes_{A}M_{2})\right).
\]
Since $\Hom_{A}(P_{\bullet},M_{2})$ and $P_{\bullet}\otimes_{A}M_{2}$ are perfect complexes, their cohomology groups are finitely presented over $A$. Since moreover $A\rightarrow A[\frac{1}{\pi}]$ is flat,
\[
\Ext_{A}^{i}(M_{1},M_{2})\!\left[{\frac{1}{\pi}}\right]\!=\Mk\Ext_{A[{\frac{1}{\pi}}]}^{i}\!\left(M_{1}\!\left[{\frac{1}{\pi}}\right]\!,M_{2}\!\left[{\frac{1}{\pi}}\right]\right)\!=\!
\begin{cases}
\Hom_{A[{\frac{1}{\pi}}]}(M_{1}[{\frac{1}{\pi}}],M_{2}[{\frac{1}{\pi}}])\!\! & \text{if }i=0,\\
0 & \text{if }i>0,
\end{cases}
\]
since $M_{1}[{\frac{1}{\pi}}]$ is finite free over $A[{
\frac{1}{\pi}}]$, and similarly
\[
\Tor_{i}^{A}(M_{1},M_{2})\!\left[{\frac{1}{\pi}}\right]=\Tor_{i}^{A[{\frac{1}{\pi}}]}\!\left(M_{1}\!\left[{\frac{1}{\pi}}\right],M_{2}\!\left[{\frac{1}{\pi}}\right]\right)=
\begin{cases}
M_{1}[{\frac{1}{\pi}}]\otimes_{A[{\frac{1}{\pi}}]}M_{2}[{\frac{1}{\pi}}] & \text{if }i=0,\\
0 & \text{if }i>0.
\end{cases}
\]
So all of these $A[{\frac{1}{\pi}}]$-modules are indeed finite and free.
\end{proof}

\subsubsection{}\label{subsec:ModAtquasiab}

The categories $\Mod_{A,\pi^{\infty}}$ and $\Mod_{A,\mathfrak{m}^{\infty}}$ are weak Serre subcategories of $\Mod_{A}$: they are stable under kernels, cokernels and extensions. In particular, they are both abelian. The category $\Mod_{A,t}$ is also stable by extensions and kernels in $\Mod_{A}$, but it is only quasi-abelian. In fact, the exact sequence (for $M\in\Mod_{A,\pi^{\infty}}$)
\[
0\rightarrow M[\mathfrak{m}^{\infty}]\rightarrow M\rightarrow M_{t}\rightarrow0
\]
yields a cotilting torsion theory~\cite{BoVdB03} on the abelian category $\Mod_{A,\pi^{\infty}}$ with torsion class $\Mod_{A,\mathfrak{m}^{\infty}}$ and torsion-free class $\Mod_{A,t}$: any $M\in\Mod_{A,\pi^{\infty}}$ is a quotient of $A_{n}^{r}\in\Mod_{A,t}$ for some $n,r\in\mathbb{N}$ (being finitely generated over $A$ and killed by a power of $\pi$), and there is no nonzero morphism from an object in $\Mod_{A,\mathfrak{m}^{\infty}}$ to an object in $\Mod_{A,t}$. The kernel and coimage of a morphism in $\Mod_{A,t}$ are the corresponding kernel and coimage in the abelian category $\Mod_{A,\pi^{\infty}}$ or $\Mod_{A}$. The image and cokernel of $f:M\rightarrow N$ in $\Mod_{A,t}$ are given by
\[
\mathrm{im}_{\Mod_{A,t}}(f)=f(M)^{\mathrm{sat}}\quad\text{and}\quad\coker_{\Mod_{A,t}}(f)=(N/f(M))_{t}=N/f(M)^{\mathrm{sat}}
\]
where
\[
f(M)^{\mathrm{sat}}/f(M)\eqd(N/f(M))[\mathfrak{m}^{\infty}]=(N/f(M))[[\varpi]^{\infty}].
\]
The morphism $f$ is strict if and only if $N/f(M)$ has no $[\varpi]$-torsion. It is a mono-epi if and only if $f$ is injective and $N/f(M)$ is killed by $[\varpi]^{n}$ for $n\gg0$. Finally, short exact sequences in $\Mod_{A,t}$ remain exact in $\Mod_{A}$.

\subsubsection{}

The categories $\Mod_{A,f}$, $\Mod_{A,\pi^{\infty}}$ and $\Mod_{A,\mathfrak{m}^{\infty}}$ are stable under the usual $\Ext$'s and $\Tor$'s in $\Mod_{A}$, and so they are also $\otimes$-categories with internal Homs (but only $\Mod_{A,f}$ has a neutral object). They are also stable under symmetric and exterior powers (of rank $k\geq1$ for the torsion categories).

The category $\Mod_{A,t}$ is stable under the internal Hom of $\Mod_{A}$, but it is not stable under the $\otimes$-product of $\Mod_{A}$. For instance, if $x\neq0$ belongs to $\mathfrak{m}_{K}^{\flat}$, then
\[
M=(\pi,[x])/(\pi^{2})
\]
is a finitely generated ideal of $A_{2}$, so it belongs to $\Mod_{A,t}$, but the image of $\pi$ in
\[
M\otimes_{A}\mathcal{O}_{K}^{\flat}=M/\pi M=(\pi,[x])/(\pi^{2},\pi[x])
\]
is a nonzero element killed by $[x]\in\mathfrak{m}\setminus\pi A$. We can nevertheless equip $\Mod_{A,t}$ with a tensor product compatible with the usual internal Hom, given by
\[
(M_{1},M_{2})\mapsto M_{1}\otimes_{t}M_{2}\eqd\left(M_{1}\otimes_{A}M_{2}\right)_{t}=\left(M_{1}\otimes_{A}M_{2}\right)/\left(M_{1}\otimes_{A}M_{2}\right)[\mathfrak{m}^{\infty}]
\]
With this definition, $\Mod_{A,t}$ becomes yet another $\otimes$-category with internal Homs.

\subsubsection{}

As explained in~\ref{subsec:basicsonModA*} or~\ref{subsec:Invariants}, there is an invariant
\[
\inv_{t}:\sk\,\Mod_{A,\pi^{\infty}}\rightarrow\mathbb{N}_{\geq}^{\infty}
\]
defined as follows: for every $M\in\Mod_{A,\pi^{\infty}}$,
\[
\inv_{t}(M)=(n_{1}\geq\dots\geq n_{s})\iff M(K)\simeq\bigoplus_{i=1}^{s}A_{n_{i}}(K).
\]
Alternatively, $\inv_{t}(M)$ is the unique element $(n_{1}\geq\dots\geq n_{s})$ of $\mathbb{N}_{\geq}^{\infty}$ such that
\[
\forall n\geq1:\qquad\rank_{\mathcal{O}_{K}^{\flat}}\left(M[\pi^{n}]/M[\pi^{n-1}]\right)=\left|\left\{ i:n_{i}\geq n\right\} \right|.
\]
This follows from~\ref{subsec:Tor(A(K)LO_L)}, which indeed implies that for every $n\geq1$,
\[
M(K)[\pi^{n}]/M(K)[\pi^{n-1}]\simeq M[\pi^{n}]/M[\pi^{n-1}]\otimes_{\mathcal{O}_{K}^{\flat}}K^{\flat}.
\]
This invariant yields a function $\rank_{t}:\sk\,\Mod_{A,\pi^{\infty}}\rightarrow\mathbb{N}$ defined by
\[
\rank_{t}(M)=\deg(\inv_{t}(M))=\sum_{i=1}^{s}n_{i}=\length_{A(K)}M(K).
\]
Plainly, $\inv_{t}(M)=\inv_{t}(M_{t})$, $\rank_{t}(M)=\rank_{t}(M_{t})$ and
\[
\begin{array}{rcccl}
\rank_{t}(M)=0 & \iff & \inv_{t}(M)=0 & \iff & {M}(K)=0\\
& \iff & M_{t}=0 & \iff & M=M[\mathfrak{m}^{\infty}].
\end{array}
\]
Moreover by~\ref{subsec:Tor(A(K)LO_L)} and Lemma~\ref{lem:InvTorsOModAndExactSeq}, for every exact sequence
\[
0\rightarrow M_{1}\rightarrow M_{2}\rightarrow M_{3}\rightarrow0
\]
in $\Mod_{A,\pi^{\infty}}$, we have $\rank_{t}(M_{2})=\rank_{t}(M_{1})+\rank_{t}(M_{3})$ and
\[
\inv_{t}(M_{2})\geq\inv_{t}(M_{1})\ast\inv_{t}(M_{3})\quad\text{in}\quad\mathbb{N}_{\geq}^{\infty}
\]
with equality if and only if the exact sequence
\[
0\rightarrow M_{1}(K)\rightarrow M_{2}(K)\rightarrow M_{3}(K)\rightarrow0
\]
of finite length $A(K)$-modules is split. In particular, the function
\[
\rank_{t}:\sk\,\Mod_{A,t}\rightarrow\mathbb{N}
\]
is a rank function on $\Mod_{A,t}$ in the sense of~\cite{Co16}: it is additive on short exact sequences, nonzero on nonzero objects, and constant on mono-epis in $\Mod_{A,t}$.

\subsubsection{}\label{subsec:rankfree/tors}

For $M\in\Mod_{A,\ast}$, let $I$ be the image of $M\rightarrow M_{f}$. For any $n\geq1$, recall that $M_{n}=M/\pi^{n}M$, which is a finitely presented $A_{n}$-module. The d\'evissage of $M$ from~\ref{subsec:DevissageOfModules} yields exact sequences of finitely presented $A_{n}$-modules
\[
0\rightarrow M[\pi^{\infty}]_{n}\rightarrow M_{n}\rightarrow I_{n}\rightarrow0\quad\text{and}\quad0\rightarrow\overline{M}[\pi^{n}]\rightarrow I_{n}\rightarrow M_{f,n}\rightarrow\overline{M}_{n}\rightarrow0.
\]
It follows that $I_{n}(K)\simeq M_{f,n}(K)$ is finite free over $A_{n}(K)$ and
\[
M_{n}(K)\simeq M[\pi^{\infty}]_{n}(K)\oplus M_{f,n}(K)\simeq M_{t,n}(K)\oplus M_{f,n}(K).
\]
In particular, $\inv_{t}M_{n}=\inv_{t}M_{t,n}\ast\inv_{t}M_{f,n}$ and
\[
\rank_{t}M_{n}=\rank_{t}M_{t,n}+n\,\rank_{A}M_{f}
\]
with $n\mapsto\rank_{t}M_{t,n}$ non-decreasing and equal to $\rank_{t}M_{t}$ for $n\gg0$.

\subsubsection{}\label{subsec:defdegonModminfty}

A \emph{good filtration} on a module $M$ in $\Mod_{A,\mathfrak{m}^{\infty}}$ is a sequence
\[
0=M_{0}\subsetneq M_{1}\subsetneq\dots\subsetneq M_{r}=M
\]
of $A$-submodules such that for all $i\in\{1,\dots,r\}$, $M_{i}/M_{i-1}\simeq\mathcal{O}_{K}^{\flat}(x_{i})$ for some nonzero $x_{i}\in\mathcal{O}_{K}^{\flat}$ (thus $M_{i}\in\Mod_{A,\mathfrak{m}^{\infty}}$ for all $i$). We have seen in~\ref{subsec:FiltOnModAt=000026ModAm}~that any $M$ in $\Mod_{A,\mathfrak{m}^{\infty}}$ has such a good filtration. We claim that the principal ideal
\[
\delta(M)\eqd\mathcal{O}_{K}^{\flat}x_{1}\dots x_{r}
\]
does not depend upon the chosen good filtration on $M$. Indeed if
\[
0=M'_{0}\subsetneq M'_{1}\subsetneq\dots\subsetneq M'_{r'}=M
\]
is another good filtration with $M'_{i}/M'_{i-1}\simeq\mathcal{O}_{K}^{\flat}(y_{i})$, $y_{i}\neq0$ in $\mathcal{O}_{K}^{\flat}$, set
\[
M_{i,j}=M_{i-1}+M'_{j}\cap M_{i}\quad\text{and}\quad M'_{j,i}=M'_{j-1}+M_{i}\cap M'_{j}
\]
Then $j\mapsto\overline{M}_{i,j}=M_{i,j}/M_{i-1}$ and $i\mapsto\overline{M}_{j,i}=M'_{j,i}/M'_{j-1}$ are good filtrations on $M_{i}/M_{i-1}$ and $M'_{j}/M'_{j-1}$ respectively, with
\[
\overline{M}_{i,j}/\overline{M}_{i,j-1}\simeq\frac{M'_{j}\cap M_{i}}{M'_{j}\cap M_{i-1}+M'_{j-1}\cap M_{i}}\simeq\overline{M}_{j,i}/\overline{M}_{j,i-1}.
\]
It is therefore sufficient to treat the case where $M=\mathcal{O}_{K}^{\flat}(x)$ for some nonzero $x\in\mathcal{O}_{K}^{\flat}$, which follows from Lemma~\ref{lem:InvTorsOModAndExactSeq}. We thus obtain a generalized length function,
\[
\length_{\mathfrak{m}^{\infty}}:\sk\,\Mod_{A,\mathfrak{m}^{\infty}}\rightarrow\mathbb{R}_{+},\qquad\length_{\mathfrak{m}^{\infty}}(M)\eqd-\log_{q}\left|\delta(M)\right|
\]
which is plainly additive on short exact sequences in $\Mod_{A,\mathfrak{m}^{\infty}}$. Here $\left|\delta\right|=\left|x\right|$ if $\delta=\mathcal{O}_{K}^{\flat}x$. Note also that $\length_{\mathfrak{m}^{\infty}}(M)=0$ if and only if $M=0$.

\subsubsection{}\label{subsec:relatinglength2rank}

For every $a\in A\setminus\pi A$, $M\mapsto M/aM$ is an exact functor from $\Mod_{A,t}$ to $\Mod_{A,\mathfrak{m}^{\infty}}$ which maps $M=\mathcal{O}_{K}^{\flat}$ to $\mathcal{O}_{K}^{\flat}(\overline{a})=\mathcal{O}_{K}^{\flat}/\overline{a}\mathcal{O}_{K}^{\flat}$, with $\overline{a}=a\bmod\pi\in\mathcal{O}_{K}^{\flat}$. It follows that for every $M\in\Mod_{A,t}$, we have the following formula:
\[
\length_{\mathfrak{m}^{\infty}}\left(M/aM\right)=-\log_{q}\left|\overline{a}\right|\cdot\rank_{t}(M).
\]
Since $\log_{q}\left|\overline{\xi'}\right|=\log_{q}\left|\varpi^{q}\right|=-1$, we obtain another formula for the rank on $\Mod_{A,t}$:
\[
\rank_{t}(M)=\length_{\mathfrak{m}^{\infty}}\left(M/\xi'M\right)\quad\text{in}\quad\mathbb{N}\subset\mathbb{R}_{+}.
\]

\subsubsection{}

The functor $M\mapsto M[\frac{1}{\pi}]$ extends to the isogeny categories,
\[
-\left[\frac{1}{\pi}\right]:\Mod_{A,f}\otimes E\rightarrow\Mod_{A,\ast}\otimes E\rightarrow\Mod_{A[\frac{1}{\pi}]}.
\]
The functor $\Mod_{A,f}\otimes E\rightarrow\Mod_{A,\ast}\otimes E$ is an equivalence of categories, with inverse induced by $M\mapsto M_{f}$. The functor $\Mod_{A,\ast}\otimes E\rightarrow\Mod_{A[\frac{1}{\pi}]}$ is fully faithful with essential image the full subcategory $\Mod_{A[\frac{1}{\pi}],f}$ of finite free $A[\frac{1}{\pi}]$-modules.

\subsection{Categories of \texorpdfstring{$\varphi$}{phi}-\texorpdfstring{$A$}{A}-modules}

\subsubsection{}

Let $\Mod_{A}^{\varphi}$ be the category of $A$-modules $M$ equipped with an $A[\xi^{\prime-1}]$-linear isomorphism $\varphi_{M}:(\varphi^{\ast}M)[\xi^{\prime-1}]\rightarrow M[\xi^{\prime-1}]$. A morphism $(M_{1},\varphi_{1})\rightarrow(M_{2},\varphi_{2})$ is an $A$-linear morphism $f:M_{1}\rightarrow M_{2}$ such that the following diagram is commutative:
\[
\xymatrix@R=2pc{(\varphi^{\ast}M_{1})[\xi^{\prime-1}]\ar[r]^{\varphi^{\ast}f}\ar[d]^{\varphi_{1}} & (\varphi^{\ast}M_{2})[\xi^{\prime-1}]\ar[d]^{\varphi_{2}}\\
M_{1}[\xi^{\prime-1}]\ar[r]^{f} & M_{2}[\xi^{\prime-1}] }
\]
Its kernel and cokernels are given by $\left(\ker(f),\varphi_{1}^{\prime}\right)$ and $\left(\coker(f),\varphi_{2}^{\prime}\right)$ with
\[
\xymatrix{(\varphi^{\ast}\ker(f))[\xi^{\prime-1}]\ar@{^{(}->}[r]\ar[d]^{\varphi_{1}^{\prime}} & (\varphi^{\ast}M_{1})[\xi^{\prime-1}]\ar[r]^{\varphi^{\ast}f}\ar[d]^{\varphi_{1}} & (\varphi^{\ast}M_{2})[\xi^{\prime-1}]\ar[d]^{\varphi_{2}}\ar@{->>}[r] & (\varphi^{\ast}\coker(f))[\xi^{\prime-1}]\ar[d]^{\varphi_{2}^{\prime}}\\
\ker(f)[\xi^{\prime-1}]\ar@{^{(}->}[r] & M_{1}[\xi^{\prime-1}]\ar[r]^{f} & M_{2}[\xi^{\prime-1}]\ar@{->>}[r] & \coker[\xi^{\prime-1}] }
\]
commutative. This makes sense since $M\mapsto M[\xi^{\prime-1}]$ and $M\mapsto\varphi^{\ast}M$ are exact. The category $\Mod_{A}^{\varphi}$ is abelian, and it is a $\otimes$-category: using the isomorphisms
\begin{align*}
\varphi^{\ast}(M_{1}\otimes_{A}M_{2})\left[\xi^{\prime-1}\right] & \simeq \left(\varphi^{\ast}(M_{1})\left[\xi^{\prime-1}\right]\right)\otimes_{A\left[\xi^{\prime-1}\right]}\left(\varphi^{\ast}(M_{2})\left[\xi^{\prime-1}\right]\right),\\
\varphi^{\ast}\left(\Sym_{A}^{k}M\right)\left[\xi^{\prime-1}\right] & \simeq \Sym_{A\left[\xi^{\prime-1}\right]}^{k}\left(\varphi^{\ast}\left(M\right)\left[\xi^{\prime-1}\right]\right),\\
\varphi^{\ast}\left(\Lambda_{A}^{k}M\right)\left[\xi^{\prime-1}\right] & \simeq \Lambda_{A\left[\xi^{\prime-1}\right]}^{k}\left(\varphi^{\ast}\left(M\right)\left[\xi^{\prime-1}\right]\right),\\
\varphi^{\ast}(A)\left[\xi^{\prime-1}\right] & \simeq A\left[\xi^{\prime-1}\right],
\end{align*}
the tensor product, symmetric and exterior powers, and neutral object are
\begin{align*}
(M_{1},\varphi_{1})\otimes(M_{2},\varphi_{2}) & \eqd \left(M_{1}\otimes M_{2},\varphi_{1}\otimes\varphi_{2}\right),\\
\Sym^{k}(M,\varphi) & \eqd \left(\Sym^{k}(M),\Sym^{k}(\varphi)\right),\\
\Lambda^{k}(M,\varphi) & \eqd \left(\Lambda^{k}(M),\Lambda^{k}(\varphi)\right),\\
\text{and}\quad A & \eqd (A,\mathrm{Id}).
\end{align*}

\subsubsection{}\label{subsec:DefBKFMod}

A \emph{Breuil--Kisin--Fargues module} or \emph{BKF-module} is an $A$-module $M$ in $\Mod_{A,\ast}$ equipped with an $A[\xi^{\prime-1}]$-linear isomorphism $\varphi_{M}:(\varphi^{\ast}M)[\xi^{\prime-1}]\rightarrow M[\xi^{\prime-1}]$. This defines a strictly full subcategory $\Mod_{A,\ast}^{\varphi}$ of $\Mod_{A}^{\varphi}$. For $\star\in\{f,\pi^{\infty},\mathfrak{m}^{\infty},t\}$, we denote by $\Mod_{A,\star}^{\varphi}$ the strictly full subcategory of $\Mod_{A,\ast}^{\varphi}$ of all BKF-modules $(M,\varphi_{M})$ whose underlying $A$-module $M$ lies in the strictly full subcategory $\Mod_{A,\star}$ of $\Mod_{A}$. Note that $\Mod_{A,\mathfrak{m}^{\infty}}^{\varphi}=\Mod_{A,\mathfrak{m}^{\infty}}$ since $M[\xi^{\prime-1}]=0$ for $M\in\Mod_{A,\mathfrak{m}^{\infty}}$.

Since $M\mapsto\varphi^{\ast}M[\xi^{\prime-1}]$ is exact, the functorial
\emph{d\'evissage} of objects in $\Mod_{A,\ast}$ yields an analogous functorial
\emph{d\'evissage} for any BKF-modules $(M,\varphi_{M})$ in $\Mod_{A,\ast}^{\varphi}$,\emph{
\[
\xymatrix@R=1pc@C=1.5pc{ & \left(M[\mathfrak{m}^{\infty}],\varphi_{M[\mathfrak{m}^{\infty}]}\right)\ar@{^{(}->}[d]\\
0\ar[r] & \left(M[\pi^{\infty}],\varphi_{M[\pi^{\infty}]}\right)\ar[r]\ar@{->>}[d] & \left(M,\varphi_{M}\right)\ar[r] & \left(M_{f},\varphi_{f}\right)\ar[r] & \left(\overline{M},\varphi_{\overline{M}}\right)\ar[r] & 0\\
& \left(M_{t},\varphi_{M_{t}}\right) }
\]
}with everyone in the relevant strictly full subcategory.

\subsubsection{}

The categories $\Mod_{A,\pi^{\infty}}^{\varphi}$ and $\Mod_{A,\mathfrak{m}^{\infty}}^{\varphi}$ are weak Serre subcategories of $\Mod_{A}^{\varphi}$: they are stable under kernels, cokernels and extensions. In particular, they are both abelian. The category $\Mod_{A,t}^{\varphi}$ is also stable under extensions and kernels in $\Mod_{A}^{\varphi}$, but it is only quasi-abelian. This last statement now requires some argument, given below: for every $M\in\Mod_{A,\pi^{\infty}}^{\varphi}$, the exact sequence
\[
0\rightarrow M[\mathfrak{m}^{\infty}]\rightarrow M\rightarrow M_{t}\rightarrow0
\]
yields a torsion theory on the abelian category $\Mod_{A,\pi^{\infty}}^{\varphi}$ with torsion class $\Mod_{A,\mathfrak{m}^{\infty}}^{\varphi}$ and torsion-free class $\Mod_{A,t}^{\varphi}$, but we do not know whether this torsion theory is cotilting (is every object $M$ of $\Mod_{A,\pi^{\infty}}^{\varphi}$ a quotient of some $N$ in $\Mod_{A,t}^{\varphi}$?), and thus we can not appeal to the criterion of~\cite[B.3]{BoVdB03} for quasi-abelian categories, as we did for $\Mod_{A,t}$ in~\ref{subsec:ModAtquasiab}. Plainly, kernels and coimages in $\Mod_{A,t}^{\varphi}$ are the corresponding kernels and coimages in $\Mod_{A}^{\varphi}$. The image and cokernel of a morphism $f:(M,\varphi_{M})\rightarrow(N,\varphi_{N})$ in $\Mod_{A,t}^{\varphi}$ are respectively equal to
\[
\left(f(M)^{\mathrm{sat}},\varphi_{f(M)^{\mathrm{sat}}}\right)\quad\text{and}\quad\left(\left(N/f(M)\right)_{t},\varphi_{(N/f(M))_{t}}\right)
\]
where the Frobeniuses are induced by $\varphi_{N}:\varphi^{\ast}(N)[\xi^{\prime-1}]\rightarrow N[\xi^{\prime-1}]$ on respectively
\begin{align*}
\varphi^{\ast}\left(f(M)^{\mathrm{sat}}\right)[\xi^{\prime-1}] & = \varphi^{\ast}\left(f(M)\right)[\xi^{\prime-1}]\\
\quad\text{and}\quad\varphi^{\ast}\left(\left(N/f(M)\right)_{t}\right)[\xi^{\prime-1}] & = \left(\varphi^{\ast}\left(N\right)/\varphi^{\ast}\left(f(M)\right)\right)[\xi^{\prime-1}].
\end{align*}
Such a morphism is strict if and only if $N/f(M)$ has no $\mathfrak{m}^{\infty}$-torsion, or $[\varpi]$-torsion. It is a monomorphism (resp.~an epimorphism) if and only if $f:M\rightarrow N$ is injective (resp.~$N/f(M)$ is killed by $[\varpi]^{n}$ for $n\gg0$). It is a strict monomorphism (resp.~a strict epimorphism) if and only if $f:M\rightarrow N$ is injective and $N/f(M)$ has no $[\varpi]$-torsion (resp.~$f:M\rightarrow N$ is surjective). We have to show that these classes of morphisms are respectively stable under arbitrary push-outs and pull-backs: this follows from the analogous properties of the quasi-abelian category $\Mod_{A,t}$, since the forgetful functor $\Mod_{A,t}^{\varphi}\rightarrow\Mod_{A,t}$ is strongly exact (i.e.~commutes with kernels and cokernels). Finally, short exact sequences in $\Mod_{A,t}^{\varphi}$ remain exact in $\Mod_{A}^{\varphi}$.

\subsubsection{}

Any BKF-module $(M,\varphi_{M})$ in $\Mod_{A,t}^{\varphi}$ has a canonical functorial filtration by strict subobjects $(M[\pi^{n}],\varphi_{M[\pi^{n}]})$ such that $M[\pi^{n}]/M[\pi^{n-1}]\simeq\pi^{n-1}M[\pi^{n}]$ is a finite free $\mathcal{O}_{K}^{\flat}$-module, and conversely, any $(M,\varphi_{M})$ in $\Mod_{A}^{\varphi}$ which is a successive extension of such BKF-modules belongs to $\Mod_{A,t}^{\varphi}$.

\subsubsection{}

The categories $\Mod_{A,\ast}^{\varphi}$, $\Mod_{A,f}^{\varphi}$, $\Mod_{A,\mathfrak{m}^{\infty}}^{\varphi}$ and $\Mod_{A,\pi^{\infty}}^{\varphi}$ are stable under the tensor product, symmetric and exterior powers of $\Mod_{A}^{\varphi}$. The isomorphism
\[
\varphi^{\ast}\left(\Hom_{A}(M_{1},M_{2})\right)\left[\xi^{\prime-1}\right]\simeq\Hom_{A[\xi^{-1}]}\left(\left(\varphi^{\ast}M_{1}\right)\left[\xi^{\prime-1}\right],\left(\varphi^{\ast}M_{2}\right)\left[\xi^{\prime-1}\right]\right)
\]
which is valid for any finitely presented $M_{1}$ also yields an internal Hom,
\[
\Hom\left(\left(M_{1},\varphi_{1}\right),\left(M_{2},\varphi_{2}\right)\right)\eqd\left(\Hom_{A}(M_{1},M_{2}),\Hom_{A[\xi^{\prime-1}]}\left(\varphi_{1}^{-1},\varphi_{2}\right)\right)
\]
on any of these categories. The subcategory $\Mod_{A,t}^{\varphi}$ of $\Mod_{A,\ast}^{\varphi}$ is stable under this internal Hom, but it is not stable under the tensor product. As for $\Mod_{A,t}$, there is a modified tensor product $(M_{1},\varphi_{1})\otimes_{t}(M_{2},\varphi_{2}):=(M_{1}\otimes_{t}M_{2},\varphi_{1}\otimes_{t}\varphi_{2})$ which turns $\Mod_{A,t}^{\varphi}$ into a genuine $\otimes$-category with internal Hom's.

\subsubsection{}

There is a Tate object $A\{1\}=\left(A\{1\},\varphi_{A\{1\}}\right)$ in $\Mod_{A,f}^{\varphi}$, defined in~\cite[4.24]{BaMoSc16}. The $A$-module $A\{1\}$ is free of rank $1$, and $\varphi_{A\{1\}}:\varphi^{\ast}(A\{1\})[\xi^{\prime-1}]\rightarrow A\{1\}[\xi^{\prime-1}]$ maps $\varphi^{\ast}A\{1\}$ onto $\xi^{\prime-1}A\{1\}$. For any BKF-module $M$ and $n\in\mathbb{Z}$ we set
\[
M\{n\}\eqd M\otimes A\{n\}\quad\text{with}\quad A\{n\}\eqd
\begin{cases}
A\{1\}^{\otimes n} & \text{if }n\geq0,\\
A\{-n\}^{\vee} & \text{if }n\leq0.
\end{cases}
\]
If $M[\mathfrak{m}^{\infty}]=0$, then $M$ and $\varphi^{\ast}M$ have no $\xi^{\prime}$-torsion, thus $M\subset M[\xi^{\prime-1}]$ and $\varphi^{\ast}M\subset\varphi^{\ast}M[\xi^{\prime-1}]$. We then say that $M$ is effective if $\varphi_{M}(\varphi^{\ast}M)\subset M$. Plainly,
\[
M\{-n\}\text{ is effective for every }n\gg0.
\]


\subsection{The Fargues filtration on \texorpdfstring{$\Mod_{A,t}^{\varphi}$}{ModA,tphi}}\label{subsec:FarguesOnModAt}

\subsubsection{}

The rank function on $\Mod_{A,t}$ yields a rank function on $\Mod_{A,t}^{\varphi}$,
\[
\rank_{t}:\sk\,\Mod_{A,t}^{\varphi}\rightarrow\mathbb{N},\qquad\rank_{t}(M,\varphi_{M})\eqd\rank_{t}(M).
\]
In addition, the length function on $\Mod_{A,\mathfrak{m}^{\infty}}$ yields a degree function on $\Mod_{A,t}^{\varphi}$,
\[
\deg_{t}:\sk\,\Mod_{A,t}^{\varphi}\rightarrow\mathbb{R}
\]
which is defined as follows. For every $(M,\varphi_{M})\in\Mod_{A,t}^{\varphi}$ and $n\gg0$, $M\{-n\}$ is effective and $\varphi_{M\{-n\}}$ maps $\varphi^{\ast}M\{-n\}$ injectively into $M\{-n\}$ with cokernel
\[
Q^{n}(M,\varphi_{M})\simeq M/\xi^{\prime n}\varphi_{M}\left(\varphi^{\ast}M\right),\qquad Q^{n}(M,\varphi_{M})\in\Mod_{A,\mathfrak{m}^{\infty}}.
\]
From the short exact sequences
\[
0\rightarrow Q^{n}(M,\varphi_{M})\stackrel{\xi^{\prime}}{\longrightarrow}Q^{n+1}(M,\varphi_{M})\rightarrow M/\xi'M\rightarrow0
\]
and~\ref{subsec:relatinglength2rank}, we thus obtain that
\[
\deg_{t}(M,\varphi_{M})\eqd n\,\rank_{t}(M)-\length_{\mathfrak{m}^{\infty}}Q^{n}(M,\varphi_{M})
\]
does not depend upon $n\gg0$. Plainly,
\[
\deg_{t}(M,\varphi_{M})\{n\}=\deg_{t}(M,\varphi_{M})+n\,\rank_{t}(M)
\]
for every $n\in\mathbb{Z}$. A short exact sequence
\[
0\rightarrow(M_{1},\varphi_{M_{1}})\rightarrow(M_{2},\varphi_{M_{2}})\rightarrow(M_{3},\varphi_{M_{3}})\rightarrow0
\]
in $\Mod_{A,t}^{\varphi}$ yields, for every $n\gg0$, a commutative diagram
\[
\xymatrix@C=1.5pc{ & 0\ar[d] & 0\ar[d] & 0\ar[d]\\
0\ar[r] & \varphi^{\ast}M_{1}\{-n\}\ar[r]\ar[d] & \varphi^{\ast}M_{2}\{-n\}\ar[r]\ar[d] & \varphi^{\ast}M_{3}\{-n\}\ar[r]\ar[d] & 0\\
0\ar[r] & M_{1}\{-n\}\ar[r]\ar[d] & M_{2}\{-n\}\ar[r]\ar[d] & M_{3}\{-n\}\ar[r]\ar[d] & 0\\
0\ar[r] & Q^{n}(M_{1},\varphi_{M_{1}})\ar[r]\ar[d] & Q^{n}(M_{2},\varphi_{M_{2}})\ar[r]\ar[d] & Q^{n}(M_{3},\varphi_{M_{3}})\ar[r]\ar[d] & 0\\
& 0 & 0 & 0 }
\]
with exact rows and columns, from which easily follows that
\[
\deg_{t}(M_{2},\varphi_{M_{2}})=\deg_{t}(M_{1},\varphi_{M_{1}})+\deg_{t}(M_{3},\varphi_{M_{3}}).
\]
Similarly a mono-epi $f:(M,\varphi_{M})\rightarrow(N,\varphi_{N})$ in $\Mod_{A,t}^{\varphi}$ yields an exact sequence
\[
0\rightarrow\ker(\varphi^{\ast}Q\rightarrow Q)\rightarrow Q^{n}(M,\varphi_{M})\rightarrow Q^{n}(N,\varphi_{N})\rightarrow\coker(\varphi^{\ast}Q\rightarrow Q)\rightarrow0
\]
in $\Mod_{A,\mathfrak{m}^{\infty}}$, where $Q=N/f(M)$ and $\varphi^{\ast}Q\rightarrow Q$ is induced by $\xi^{\prime n}\varphi_{N}$. Thus
\begin{align*}
\deg_{t}(N,\varphi_{N})-\deg_{t}(M,\varphi_{M}) & = \length_{\mathfrak{m}^{\infty}}\varphi^{\ast}Q-\length_{\mathfrak{m}^{\infty}}Q\\
& = (q-1)\cdot\length_{\mathfrak{m}^{\infty}}Q
\end{align*}
and $\deg_{t}(M,\varphi_{M})\leq\deg_{t}(N,\varphi_{N})$ with equality if and only if $f$ is an isomorphism.

\subsubsection{}

We may now apply the Harder--Narasimhan formalism of~\cite{An09} or~\cite{Co16} to the quasi-abelian category $\Mod_{A,t}^{\varphi}$, equipped with the rank and degree functions that we have just defined, for the slope function $\mu=\deg_{t}/\rank_{t}$. Specializing the general theory to the case at hand, we obtain the following definitions and results.

A BKF-module $M$ in $\Mod_{A,t}^{\varphi}$ is semi-stable of slope $\mu\in\mathbb{R}$ if and only if for every strict subobject $N$ of $M$, $\deg_{t}(N)\leq\mu\,\rank_{t}(N)$ with equality for $N=M$. With this definition, the trivial BKF-module is semi-stable of slope $\mu$ for all $\mu\in\mathbb{R}$.

The semi-stable BKF-modules of slope $\mu$ form an abelian full subcategory of $\Mod_{A,t}^{\varphi}$, and every BKF-module $M$ in $\Mod_{A,t}^{\varphi}$ has a unique decreasing $\mathbb{R}$-filtration $\mathcal{F}$ by strict subobjects $\mathcal{F}^{\geq\gamma}$ with $\Gr_{\mathcal{F}}^{\gamma}:=\mathcal{F}^{\geq\gamma}/\mathcal{F}^{>\gamma}$ semi-stable of slope $\gamma$ for all $\gamma\in\mathbb{R}$, where $\mathcal{F}^{>\gamma}:=\cup_{\gamma'>\gamma}\mathcal{F}^{\geq\gamma'}$. We call $\mathcal{F}_{F}(M):=\mathcal{F}$ the Fargues filtration of $M$. It depends functorially upon $M$ and there is no nonzero morphism from $M_{1}$ to $M_{2}$ if $M_{1}$ and $M_{2}$ are semi-stable of slope $\mu_{1}$ and $\mu_{2}<\mu_{1}$.

The Fargues type of $M$ is the type $t_{F}(M)\in\mathbb{R}_{\geq}^{r}$ of $\mathcal{F}_{F}(M)$, with $r=\rank_{t}(M)$. Viewed as a concave polygon over $[0,r]$, it is the upper boundary of the convex envelope in $[0,r]\times\mathbb{R}$ of $\{(\rank_{t},\deg_{t})(N):N\subset M\}$, where $N$ ranges either through the set of all subobjects of $M$, or through the subset of all strict subobjects of $M$. Finally, we denote by $\Gr_{F}^{\bullet}(M)$ the associated graded object in $\Mod_{A,t}^{\varphi}$.

\begin{prop}\label{prop:PseudoIsogInModt}
If $M_{1}\rightarrow M_{2}$ is a mono-epi in $\Mod_{A,t}^{\varphi}$ with cokernel $Q$ in $\Mod_{A,\mathfrak{m}^{\infty}}^{\varphi}$ and $r=\rank_{t}M_{1}=\rank_{t}M_{2}$, then for every $s\in[0,r]$,
\[
0\leq t_{F}(M_{2})(s)-t_{F}(M_{1})(s)\leq(q-1)\cdot\length_{\mathfrak{m}^{\infty}}Q
\]
with equality on the left (resp.~right) for $s=0$ (resp.~$s=r$). In particular,
\[
0\leq\left\{\!
\begin{array}{c}
t_{F}^{\max}(M_{2})-t_{F}^{\max}(M_{1})\\
t_{F}^{\min}(M_{2})-t_{F}^{\min}(M_{1})
\end{array}\!\right\} \leq(q-1)\cdot\length_{\mathfrak{m}^{\infty}}Q.
\]
\end{prop}

\begin{proof}
Set $f_{i}=t_{F}(M_{i})$ and $(r_{i},d_{i})(\gamma)=(\rank_{t},\deg_{t})(\mathcal{F}_{F}^{\gamma}(M_{i}))$ for $\gamma\in\mathbb{R}$ and $i\in\{1,2\}$. It is sufficient to show that for every $\gamma\in\mathbb{R}$,
\[
d_{1}(\gamma)\leq f_{2}(r_{1}(\gamma))\quad\text{and}\quad d_{2}(\gamma)\leq f_{1}(r_{2}(\gamma))+(q-1)\cdot\length_{\mathfrak{m}^{\infty}}Q.
\]
For the first inequality, let $\mathcal{F}_{F}^{\gamma}(M_{1})^{\mathrm{sat}}$ be the image of $\mathcal{F}_{F}^{\gamma}(M_{1})$ in $M_{2}$. Then
\[
d_{1}(\gamma)=\deg_{t}\mathcal{F}_{F}^{\gamma}(M_{1})\leq\deg_{t}\mathcal{F}_{F}^{\gamma}(M_{1})^{\mathrm{sat}}\leq f_{2}(r_{1}(\gamma))
\]
since $\mathcal{F}_{F}^{\gamma}(M_{1})\rightarrow\mathcal{F}_{F}^{\gamma}(M_{1})^{\mathrm{sat}}$ is a mono-epi and $\mathcal{F}_{F}^{\gamma}(M_{1})^{\mathrm{sat}}$ is a strict subobject of rank $r_{1}(\gamma)$ in $M_{2}$. For the second inequality, let $\mathcal{F}_{F}^{\gamma}(M_{2})^{\prime}$ and $Q^{\gamma}$ be respectively the kernel and image of $\mathcal{F}_{F}^{\gamma}(M_{2})\rightarrow Q$ in $\Mod_{A,\ast}^{\varphi}$. Then
\begin{align*}
d_{2}(\gamma)=\deg_{t}\mathcal{F}_{F}^{\gamma}(M_{2}) & = \deg_{t}\mathcal{F}_{F}^{\gamma}(M_{2})^{\prime}+(q-1)\cdot\length_{\mathfrak{m}^{\infty}}Q^{\gamma}\\
& \leq f_{1}(r_{2}(\gamma))+(q-1)\cdot\length_{\mathfrak{m}^{\infty}}Q
\end{align*}
since $\mathcal{F}_{F}^{\gamma}(M_{2})^{\prime}$ is a strict subobject of rank $r_{2}(\gamma)$ in $M_{1}$ and $Q^{\gamma}\subset Q$.
\end{proof}

\begin{prop}\label{prop:ExactSeqInModt}
Let $0\rightarrow M_{1}\rightarrow M_{2}\rightarrow M_{3}\rightarrow0$ be an exact sequence in $\Mod_{A,t}^{\varphi}$, set $r_{i}=\rank_{t}M_{i}$ and view $t_{F}(M_{i})$ as a concave function $f_{i}:[0,r_{i}]\rightarrow\mathbb{R}$. Then
\[
f_{1}\ast f_{3}(s)\geq f_{2}(s)\geq
\begin{cases}
f_{1}(s) & \text{if }0\leq s\leq r_{1}\\
f_{1}(r_{1})+f_{3}(s-r_{1}) & \text{if }r_{1}\leq s\leq r_{2}
\end{cases}
\]
with equality for $s=0$ and $s=r_{2}$. In particular,
\begin{gather*}
\begin{aligned}
t_{F}^{\max}(M_{1}) & \leq t_{F}^{\max}(M_{2}) \leq \max\left\{ t_{F}^{\max}(M_{1}),t_{F}^{\max}(M_{3})\right\},\\
t_{F}^{\min}(M_{3}) & \geq t_{F}^{\min}(M_{2})\, \geq \min\left\{ t_{F}^{\min}(M_{1}),t_{F}^{\min}(M_{3})\right\},
\end{aligned}
\\
\text{and}\qquad t_{F}(M_{2})\leq t_{F}(M_{1})\ast t_{F}(M_{3})\quad\text{in}\quad\mathbb{R}_{\geq}^{r_{2}}.
\end{gather*}
Moreover, $t_{F}(M_{2})=t_{F}(M_{1})\ast t_{F}(M_{3})$ if and only if for every $\gamma\in\mathbb{R}$,
\[
0\rightarrow\mathcal{F}_{F}^{\gamma}(M_{1})\rightarrow\mathcal{F}_{F}^{\gamma}(M_{2})\rightarrow\mathcal{F}_{F}^{\gamma}(M_{3})\rightarrow0
\]
is exact.
\end{prop}

\begin{proof}
These are standard properties of Harder--Narasimhan filtrations on quasi-abelian categories, see for instance~\cite[Proposition~21]{Co16} or~\cite[4.4.4]{An09}.
\end{proof}

\begin{prop}\label{prop:tF(twist)onModt}
For every $M\in\Mod_{A,t}^{\varphi}$ of rank $r\in\mathbb{N}$ and any $n\in\mathbb{Z}$,
\[
\mathcal{F}_{F}^{\gamma}(M\{n\})=\mathcal{F}_{F}^{\gamma-n}(M)\{n\}
\]
for every $\gamma\in\mathbb{R}$, hence
\[
t_{F}(M\{n\})=t_{F}(M)+(n,\dots,n)\quad\text{in}\quad\mathbb{R}_{\geq}^{r}.
\]
\end{prop}

\begin{proof}
This is obvious: the map $N\mapsto N\{n\}$ induces a bijection between strict subobjects of $M$ and strict subobjects of $M\{n\}$, with $\mu(N\{n\})=\mu(N)+n$.
\end{proof}

\subsubsection{}\label{subsec:ExactSeqInModpi}

For $M\in\Mod_{A,\pi^{\infty}}^{\varphi}$, we set $t_{F}(M)=t_{F}(M_{t})$. An exact sequence
\[
0\rightarrow M_{1}\rightarrow M_{2}\rightarrow M_{3}\rightarrow0
\]
in $\Mod_{A,\pi^{\infty}}^{\varphi}$ gives rise to three exact sequences:
\begin{gather*}
0\rightarrow M_{1}[\mathfrak{m}^{\infty}]\rightarrow M_{2}[\mathfrak{m}^{\infty}]\rightarrow M_{3}[\mathfrak{m}^{\infty}]\rightarrow Q\rightarrow0
\\
0\rightarrow M_{4}\rightarrow M_{2,t}\rightarrow M_{3,t}\rightarrow0
\\
0\rightarrow M_{1,t}\rightarrow M_{4}\rightarrow Q\rightarrow0
\end{gather*}
with $Q\in\Mod_{A,\mathfrak{m}^{\infty}}^{\varphi}$ and $M_{4}\in\Mod_{A,t}^{\varphi}$. Set $\ell_{i}=\length_{\mathfrak{m}^{\infty}}M_{i}[\mathfrak{m}^{\infty}]$, $r_{i}=\rank_{t}(M_{i})$, $f_{i}=t_{F}(M_{i})$ and $\ell_{Q}=\length_{\mathfrak{m}^{\infty}}Q$. We thus have the following relations:
\begin{gather*}
r_{1}=r_{4},\quad r_{1}+r_{3}=r_{2},\quad\ell_{Q}=\ell_{1}-\ell_{2}+\ell_{3}
\\
f_{4}\ast f_{3}(s)\geq f_{2}(s)\geq
\begin{cases}
f_{4}(s) & \text{for }0\leq s\leq r_{1},\\
f_{4}(r_{1})+f_{3}(s-r_{1}) & \text{for }r_{1}\leq s\leq r_{2}.
\end{cases}
\\
f_{1}(s)\leq f_{4}(s)\leq f_{1}(s)+(q-1)\ell_{Q}\quad\text{for }0\leq s\leq r_{1}.
\end{gather*}
Set $c_{i}=\max\left\{
\left|t_{F}^{\min}(M_{i})\right|,\left|t_{F}^{\max}(M_{i})\right|\right\} $ so that $f_{i}$ is $c_{i}$-Lipschitzian and
\[
c_{2}\leq\max\left\{ c_{3},c_{4}\right\},\quad\left|c_{1}-c_{4}\right|\leq(q-1)\ell_{Q}.
\]
Moreover, we have
\[
f_{4}\ast f_{3}(s)\leq
\begin{cases}
f_{4}(s)+c_{4}r_{3}+\max(f_{3}) & \text{for }0\leq s\leq r_{1},\\
f_{3}(s-r_{1})+c_{3}r_{1}+\max(f_{4}) & \text{for }r_{1}\leq s\leq r_{2}
\end{cases}
\]
We obtain the following inequalities: for $0\leq s\leq r_{1}$,
\[
0\le f_{2}(s)-f_{1}(s)\leq\left(c_{1}+c_{3}\right)r_{3}+(q-1)(\ell_{1}+\ell_{3})(r_{3}+1)
\]
and for $r_{1}\leq s\leq r_{2}$,
\[
-c_{1}r_{1}\leq f_{2}(s)-f_{3}(s-r_{1})\leq\left(c_{1}+c_{3}+(q-1)(\ell_{1}+\ell_{3})\right)r_{1}
\]
which also implies that for $0\leq s\leq r_{3}$,
\[
\left|f_{2}(s)-f_{3}(s)\right|\leq\max\left\{\!
\begin{array}{c}
c_{1}+2c_{3}+(q-1)(\ell_{1}+\ell_{3}),\\
2c_{1}+c_{3}+2(q-1)(\ell_{1}+\ell_{3})
\end{array}\!\right\} \cdot r_{1}.
\]

\subsection{The Fargues type on \texorpdfstring{$\Mod_{A,\ast}^{\varphi}$}{ModA, astphi}}
\label{subsec:FarguesTypeBKF}

\subsubsection{}

For any $M\in\Mod_{A}^{\varphi}$ and $n\geq1$, consider the exact sequence
\[
\xymatrix@C=2pc{0\ar[r] & M[\pi^{n}]\ar[r] & M\ar[r]^{\pi^{n}} & M\ar[r] & M_{n}\ar[r] & 0}.
\]
Suppose that $M$ is a BKF-module, i.e.~belongs to $\Mod_{A,\ast}^{\varphi}$. Then $M_{n}$ and $M[\pi^{n}]$ both belong to $\Mod_{A,\pi^{\infty}}^{\varphi}$. Moreover, $\rank_{t}M_{n}\geq n\,\rank_{A}M$ by~\ref{subsec:rankfree/tors}. Viewing $t_{F}(M_{n})$ as a concave function on $[0,\rank_{t}M_{n}]$, we may thus define
\[
t_{F,n}(M):[0,\rank_{A}M]\rightarrow\mathbb{R},\qquad t_{F,n}(M)(s)={\frac{1}{n}}t_{F}(M_{n})(ns).
\]

\begin{prop}\label{prop:DeftFinfty}
There is a constant $C(M)$ such that the functions $t_{F,n}(M)$ are $C(M)$-Lipschitzian. They converge uniformly to a continuous concave function
\[
t_{F,\infty}(M):[0,\rank_{A}M]\rightarrow\mathbb{R}.
\]
If $M_{1},M_{2}\in\Mod_{A,\ast}^{\varphi}$ become isomorphic in the isogeny category $\Mod_{A,\ast}^{\varphi}\otimes E$, then
\[
t_{F,\infty}(M_{1})=t_{F,\infty}(M_{2}).
\]
\end{prop}

\begin{proof}
Let $r=\rank_{A}M=\rank_{A}M_{f}$ and set $f_{n}=f_{n}(M)=t_{F,n}(M)$.

Suppose first that $M$ is free. Then for every $n,m\geq1$, the exact sequence
\[
0\rightarrow M_{n}\stackrel{\pi^{m}}{\longrightarrow}M_{n+m}\rightarrow M_{m}\rightarrow0
\]
in $\Mod_{A,t}^{\varphi}$ gives the inequality
\[
t_{F}(M_{n+m})\leq t_{F}(M_{n})\ast t_{F}(M_{m})\quad\text{in}\quad\mathbb{R}_{\geq}^{r(n+m)}.
\]
It follows that for every $n,k\geq1$ and $0\leq s\leq r$,
\[
f_{nk}(s)\leq f_{n}(s)\quad\text{with equality for }s\in\{0,r\}.
\]
In particular, $f_{n}(s)\leq f_{1}(s)$ with equality for $s\in\{0,r\}$, and the slopes of the continuous piecewise linear functions $f_{n}$ are uniformly bounded by the constant
\[
C=C(M)=\max\left\{ \left|t_{F}^{\min}(M_{1})\right|,\left|t_{F}^{\max}(M_{1})\right|\right\}.
\]
Fix $n_{0},n\geq1$. For $n=n_{0}q_{n}+r_{n}$ with $q_{n}\geq0$ and $0\leq r_{n}<n_{0}$, we have
\[
t_{F}(M_{n})\leq t_{F}(M_{n_{0}q_{n}})\ast t_{F}(M_{r_{n}})\leq t_{F}(M_{n_{0}})^{\ast^{q_{n}}}\ast t_{F}(M_{r_{n}})
\]
from which we obtain that for $0\leq s\leq r$,
\[
f_{n}(s)\leq\left(1-\frac{r_{n}}{n}\right)f_{n_{0}}(s^{\prime})+{\frac{r_{n}}{n}}f_{r_{n}}(s^{\prime\prime})
\]
for some $s^{\prime},s^{\prime\prime}\in[0,r]$ with $n_{0}q_{n}s^{\prime}+r_{n}s^{\prime\prime}=ns$. But then $s^{\prime}-s=\frac{r_{n}}{n}(s^{\prime}-s^{\prime\prime})$, thus
\[
f_{n}(s)\leq\left(1-\frac{r_{n}}{n}\right)f_{n_{0}}(s)+{\frac{r_{n}}{n}}\left(2rC\left(1-\frac{r_{n}}{n}\right)+\sup(f_{1})\right).
\]
Therefore $\limsup f_{n}(s)\leq f_{n_{0}}(s)$ and this being true for all $n_{0}\geq1$,
\[
\limsup f_{n}(s)\leq\liminf f_{n}(s)
\]
i.e.~$f_{n}(s)$ converges to some limit $f_{\infty}(s)\in\mathbb{R}$. Since all the $f_{n}$'s are $C$-Lipschitzian concave, so is $f_{\infty}=f_{\infty}(M)$ and the convergence is uniform.

Suppose next that $M$ is torsion free, so that $0\rightarrow M\rightarrow M_{f}\rightarrow\overline{M}\rightarrow0$ is exact and for $n\gg0$ (such that $\pi^{n}\overline{M}=0$), we obtain an exact sequence
\[
0\rightarrow\overline{M}\rightarrow M_{n}\rightarrow M_{f,n}\rightarrow\overline{M}\rightarrow0
\]
which identifies $\overline{M}$ and $M_{n}[\mathfrak{m}^{\infty}]$ (since $M_{f,n}[\mathfrak{m}^{\infty}]=0$), i.e.
\[
0\rightarrow M_{n,t}\rightarrow M_{f,n}\rightarrow\overline{M}\rightarrow0
\]
is exact. Our claim now follows from Proposition~\ref{prop:PseudoIsogInModt}, with the constant
\[
C(M)=C(M_{f})+(q-1)\length_{\mathfrak{m}^{\infty}}\overline{M}
\]
and the limit $f_{\infty}(M)=f_{\infty}(M_{f})$.

For the general case, let $I$ be the image of $M\rightarrow M_{f}$, so that $I$ is a torsion free BKF-module. This time for $n\gg0$, we have an exact sequence
\[
0\rightarrow M[\pi^{\infty}]\rightarrow M_{n}\rightarrow I_{n}\rightarrow0.
\]
We have just seen that $I_{n}[\mathfrak{m}^{\infty}]=\overline{I}=\overline{M}$ for $n\gg0$. Our claim now follows from the discussion of Section~\ref{subsec:ExactSeqInModpi} with the constant
\[
C(M)=\max\left\{ C(I),C(M_{t})+(q-1)\length_{\mathfrak{m}^{\infty}}M[\mathfrak{m}^{\infty}]\oplus\overline{M}\right\}
\]
and the limit $f_{\infty}(M)=f_{\infty}(I)=f_{\infty}(M_{f})$. Here
\begin{align*}
C(I) & = \max\left\{ \left|t_{F}^{\min}(M_{f,1})\right|,\left|t_{F}^{\max}(M_{f,1})\right|\right\} +(q-1)\length_{\mathfrak{m}^{\infty}}\overline{M},\\
C(M_{t}) & = \max\left\{ \left|t_{F}^{\min}(M_{t})\right|,\left|t_{F}^{\max}(M_{t})\right|\right\}.
\end{align*}

It remains to establish that $M\mapsto t_{F,\infty}(M)$ is constant on isogeny classes, and we already know that $t_{F,\infty}(M)=t_{F,\infty}(M_{f})$. We thus have to show that if
\[
0\rightarrow M_{1}\rightarrow M_{2}\rightarrow Q\rightarrow0
\]
is an exact sequence in $\Mod_{A}^{\varphi}$ with $M_{1}$, $M_{2}$ finite free and $Q$ torsion, then $t_{F,\infty}(M_{1})$ equals $t_{F,\infty}(M_{2})$. For $n\gg0$ (such that $\pi^{n}Q=0$), we obtain exact sequences
\[
0\rightarrow Q\rightarrow M_{1,n}\rightarrow M_{2,n}\rightarrow Q\rightarrow0.
\]
Splitting them in two short exact sequences and using again the computations of Section~\ref{subsec:ExactSeqInModpi} yields the desired equality.
\end{proof}

\begin{prop}\label{prop:tFinfty(twist)onMod*}
For any BKF-module $M$ of rank $r\in\mathbb{N}$ and any $n\in\mathbb{Z}$,
\[
\forall s\in[0,r]:\quad t_{F,\infty}(M\{n\})(s)=t_{F,\infty}(M)(s)+sn.
\]
\end{prop}

\begin{proof}
This follows from Proposition~\ref{prop:tF(twist)onModt}.
\end{proof}

\subsubsection{}\label{subsec:TypeHN}

The first part of the proof of Proposition~\ref{prop:DeftFinfty} shows that

\begin{prop}\label{prop:ComptFinftytFn}
For a finite free BKF-module $M$ of rank $r\in\mathbb{N}$,
\[
t_{F,\infty}(M)(s)\leq t_{F,n}(M)(s)\leq t_{F,1}(M)(s)
\]
for every $n\geq1$ and $s\in[0,r]$ with equality for $s\in\{0,r\}$.
\end{prop}

\begin{defi}\label{Def:HNType}
We say that a finite free BKF-module $M$ is of HN-type if
\[
t_{F,\infty}(M)=t_{F,1}(M).
\]
\end{defi}

Thus if $M$ is of HN-type with rank $r\in\mathbb{N}$, then $t_{F,\infty}(M)\in\mathbb{R}_{\geq}^{r}$.

\begin{prop}\label{prop:TypeHN}
Let $M$ be a finite free BKF-module of HN-type. Then
\begin{enumerate}
\item \label{2.10.1} For every $\gamma\in\mathbb{R}$ and $n,m\geq1$, the exact sequence
\[
\xymatrix{0\ar[r] & M_{n}\ar[r]^{\pi^{m}} & M_{n+m}\ar[r] & M_{m}\ar[r] & 0}
\]
induces an exact sequence
\[
\xymatrix{0\ar[r] & \mathcal{F}_{F}^{\gamma}(M_{n})\ar[r]^{\pi^{m}} & \mathcal{F}_{F}^{\gamma}(M_{n+m})\ar[r] & \mathcal{F}_{F}^{\gamma}(M_{m})\ar[r] & 0}
\]
\item \label{2.10.2} The formula $\mathcal{F}_{F}^{\gamma}(M)=\underleftarrow{\lim}\mathcal{F}_{F}^{\gamma}(M_{n})$ defines an $\mathbb{R}$-filtration on $M$ by finite free BKF-submodules whose underlying $A$-submodules are direct summands: the quotient $\Gr_{F}^{\gamma}(M)=\mathcal{F}_{F}^{\geq\gamma}(M)/\mathcal{F}_{F}^{>\gamma}(M)$ is a finite free BKF-module.
\item \label{2.10.3} For every $\gamma\in\mathbb{R}$ and $n\geq1$,
\[
\mathcal{F}_{F}^{\gamma}(M)_{n}=\mathcal{F}_{F}^{\gamma}(M_{n})\quad\text{ and}\quad\Gr_{F}^{\gamma}(M)_{n}=\Gr_{F}^{\gamma}(M_{n}).
\]
In particular, the type of the $\mathbb{R}$-filtration $\mathcal{F}_{F}^{\bullet}(M)$ is given by
\[
\mathbf{t}\left(\mathcal{F}_{F}^{\bullet}(M)\right)=t_{F}(M_{1})=t_{F,1}(M)=t_{F,\infty}(M).
\]
\end{enumerate}
\end{prop}

\begin{proof}
\eqref{2.10.1} Since $t_{F,\infty}(M)=t_{F,1}(M)$, also $t_{F,n}(M)=t_{F,1}(M)$ for every $n\geq1$, thus
\[
t_{F}(M_{n+m})=t_{F}(M_{n})\ast t_{F}(M_{m})
\]
for every $n,m\geq1$, from which~\eqref{2.10.1} immediately follows by Proposition~\ref{prop:ExactSeqInModt}.

\eqref{2.10.2} and~\eqref{2.10.3}: This follows from~\eqref{2.10.1} by a standard argument: consider for $n,m\geq1$ and $\gamma\in\mathbb{R}$ the commutative diagram with exact rows and columns
\[
\xymatrix{ & 0\ar[d] & 0\ar[d] & 0\ar[d]\\
0\ar[r] & \mathcal{F}_{F}^{\gamma}(M_{n})\ar[d]\ar[r]^{\pi^{m}} & \mathcal{F}_{F}^{\gamma}(M_{n+m})\ar[d]\ar[r] & \mathcal{F}_{F}^{\gamma}(M_{m})\ar[d]\ar[r] & 0\\
0\ar[r] & M_{n}\ar[d]\ar[r]^{\pi^{m}} & M_{n+m}\ar[d]\ar[r] & M_{m}\ar[d]\ar[r] & 0\\
0\ar[r] & \mathcal{G}_{F}^{\gamma}(M_{n})\ar[d]\ar[r]^{\pi^{m}} & \mathcal{G}_{F}^{\gamma}(M_{n+m})\ar[d]\ar[r] & \mathcal{G}_{F}^{\gamma}(M_{m})\ar[d]\ar[r] & 0\\
& 0 & 0 & 0 }
\]
Taking the projective limit over $n$, and since every one is Mittag--Leffler surjective, we obtain a commutative diagram of $A$-modules with exact rows and columns
\[
\xymatrix{ & 0\ar[d] & 0\ar[d] & 0\ar[d]\\
0\ar[r] & \mathcal{F}_{F}^{\gamma}(M)\ar[d]\ar[r]^{\pi^{m}} & \mathcal{F}_{F}^{\gamma}(M)\ar[d]\ar[r] & \mathcal{F}_{F}^{\gamma}(M_{m})\ar[d]\ar[r] & 0\\
0\ar[r] & M\ar[d]\ar[r]^{\pi^{m}} & M\ar[d]\ar[r] & M_{m}\ar[d]\ar[r] & 0\\
0\ar[r] & \mathcal{G}_{F}^{\gamma}(M)\ar[d]\ar[r]^{\pi^{m}} & \mathcal{G}_{F}^{\gamma}(M)\ar[d]\ar[r] & \mathcal{G}_{F}^{\gamma}(M_{m})\ar[d]\ar[r] & 0\\
& 0 & 0 & 0 }
\]
Since $\mathcal{F}_{F}^{\gamma}(M)=\underleftarrow{\lim}\,\mathcal{F}_{F}^{\gamma}(M_{n})$, the first row tells us that $\mathcal{F}_{F}^{\gamma}(M)$ is separated and complete in the $\pi$-adic topology, with $\mathcal{F}_{F}^{\gamma}(M)_{1}\simeq\mathcal{F}_{F}^{\gamma}(M_{1})$ finite free over $A_{1}=\mathcal{O}_{K}^{\flat}$, say of rank $s\in\mathbb{N}$. Pick a morphism $\alpha:A^{s}\rightarrow\mathcal{F}_{F}^{\gamma}(M)$ reducing to an isomorphism modulo $\pi$. By the topological version of Nakayama's lemma, $\alpha$ is surjective, and $\mathcal{F}_{F}^{\gamma}(M)$ is finitely generated over $A$. Playing the same game with the third row, we obtain a surjective morphism $\beta:A^{s'}\twoheadrightarrow\mathcal{G}_{F}^{\gamma}(M)$ reducing to an isomorphism modulo $\pi$. But now the kernel $N$ of $\beta$ has to be finitely generated over $A$ since $\mathcal{G}_{F}^{\gamma}(M)$ is finitely presented over $A$ by the second column. Applying $\Tor_{\bullet}^{A}(-,\mathcal{O}_{K}^{\flat})$ to the resulting short exact sequence $0\rightarrow N\rightarrow A^{s^{\prime}}\rightarrow\mathcal{G}_{F}^{\gamma}(M)\rightarrow0$, we find that
\[
N\otimes\mathcal{O}_{K}^{\flat}\simeq\Tor_{1}^{A}\left(\mathcal{G}_{F}^{\gamma}(M),\mathcal{O}_{K}^{\flat}\right)\simeq\mathcal{G}_{F}^{\gamma}(M)[\pi],
\]
which is trivial by the third row, thus $N=0$ by the classical version of Nakayama's lemma. It follows that $\beta$ is an isomorphism, $\mathcal{G}_{F}^{\gamma}(M)$ is free, the middle column is split (in $\Mod_{A}$), and $\mathcal{F}_{F}^{\gamma}(M)$ is also free, being finite projective over the local ring $A$. The remaining assertions of \eqref{2.10.2} and~\eqref{2.10.3} easily follow.
\end{proof}

\begin{rema}
For a finite free BKF-module $M$ of HN-type and $n\in\mathbb{Z}$, the Tate twist $M\{n\}$ is also of HN-type and $\mathcal{F}_{F}^{\gamma}(M\{n\})=\mathcal{F}_{F}^{\gamma-n}(M)\{n\}$ by Proposition~\ref{prop:tF(twist)onModt}.
\end{rema}

\begin{defi}
We say that a finite free BKF-module $M$ is semi-stable (of slope $\gamma\in\mathbb{R}$) if $M_{1}$ is semi-stable (of slope $\gamma\in\mathbb{R}$).
\end{defi}

\begin{example}
Any finite free BKF-module $M$ of rank $1$ is semi-stable of slope $\deg_{t}(M_{1})$, thus $A$ is semi-stable of slope $0$ and $A\{1\}$ is semi-stable of slope $1$.
\end{example}

By~Proposition~\ref{prop:ComptFinftytFn}, a finite free BKF-module $M$ is semi-stable (of slope $\gamma$) if and only if $M_{n}$ is semi-stable (of slope $\gamma$) for every $n\geq1$, in which case $M$ is of HN-type and $t_{F,\infty}(M)=t_{F,1}(M)=t_{F}(M_{1})$ is isoclinic (of slope $\gamma$). By Proposition~\ref{prop:TypeHN}, a finite free BKF-module $M$ of HN-type has a canonical filtration $\mathcal{F}_{F}(M)$ whose graded pieces are finite free semi-stable BKF-modules with decreasing slopes. Conversely, any finite free BKF-module which has such a filtration is of HN-type (by uniqueness of the Fargues filtration on $\Mod_{A,t}^{\varphi}$) and its filtration is the canonical one.

\subsubsection{}\label{subsec:DefFullSubCatHNType}

We denote by $\Mod_{A,f}^{\varphi,\ast}$ the strictly full subcategory of $\Mod_{A,f}^{\varphi}$ whose objects are the finite free BKF-modules of HN-type. The functoriality of the Fargues filtration on $\Mod_{A,t}^{\varphi}$ implies that $M\mapsto\mathcal{F}_{F}(M)$ is functorial on $\Mod_{A,f}^{\varphi,\ast}$.

\begin{prop}\label{prop:catOfHNTypeTensor}
The subcategory $\Mod_{A,f}^{\varphi,\ast}$ of $\Mod_{A,f}^{\varphi}$ is stable under $\otimes$-products and inner Homs and the $\mathbb{R}$-filtration $\mathcal{F}_{F}$ on $\Mod_{A,f}^{\varphi,\ast}$ is compatible with them.
\end{prop}
\begin{proof}
The Fargues filtrations on $M,M'\in\Mod_{A,f}^{\varphi,\ast}$ induce $\mathbb{R}$-filtrations on $M\otimes M'$ and $\Hom(M,M')$ whose graded pieces are the finite free BKF-modules
\[
\bigoplus_{\gamma_{1}+\gamma_{2}=\gamma}\Gr_{F}^{\gamma_{1}}(M)\otimes\Gr_{F}^{\gamma_{2}}(M')\quad\text{and}\quad\bigoplus_{\gamma_{2}-\gamma_{1}=\gamma}\Hom(\Gr_{F}^{\gamma_{1}}(M),\Gr_{F}^{\gamma_{2}}(M')).
\]
We thus have to show that if $M$ and $M'$ are semi-stable of slope $\gamma$ and $\gamma'$, then $P=M\otimes M'$ and $H=\Hom(M,M')$ are semi-stable of slope $\gamma+\gamma'$ and $\gamma'-\gamma$. Since $P_{1}=M_{1}\otimes M_{1}^{\prime}$ and $H_{1}=\Hom(M_{1},M_{1}^{\prime})$, we need to establish the analogous statement for BKF-modules which are finite free over $A_{1}=\mathcal{O}_{K}^{\flat}$. This is a special case of~\cite[\S 5.3]{Co16}, see also Section~\ref{subsec:RisOkflat} below.
\end{proof}

\subsection{Categories of \texorpdfstring{$\varphi$}{phi}-\texorpdfstring{$R$}{R}-modules}

For any $A$-algebra $R$ equipped with a ring isomorphism $\varphi:R\rightarrow R$ compatible with $\varphi:A\rightarrow A$, we may analogously define the abelian $\otimes$-category $\Mod_{R}^{\varphi}$ and its full $\otimes$-subcategories $\Mod_{R,\ast}^{\varphi}$ and $\Mod_{R,f}^{\varphi}$. They come equipped with $\otimes$-functors $\Mod_{A,?}^{\varphi}\rightarrow\Mod_{R,?}^{\varphi}$ for $?\in\{\emptyset,\ast,f\}$, which are exact when $A\rightarrow R$ is flat. In this section, we discuss the following cases:
\[
R\in\left\{ \mathcal{O}_{K}^{\flat},L,\mathcal{O}_{L},A(K),A\left[{\frac{1}{\pi}}\right]\right\}.
\]

\subsubsection{\texorpdfstring{$R=\mathcal{O}_{K}^{\flat}$}{R=OK flat}}\label{subsec:RisOkflat}

In this case, $\Mod_{R,f}^{\varphi}$ is the full subcategory of $\Mod_{A,t}^{\varphi}$ made of all BKF-modules killed by $\pi$. This is the quasi-abelian category of all finite free $\mathcal{O}_{K}^{\flat}$-modules $M$ equipped with an isomorphism $\varphi_{M}:\varphi^{\ast}M\otimes K^{\flat}\rightarrow M\otimes K^{\flat}$, or equivalently, with a $\varphi$-semilinear isomorphism $\phi_{M}:M\otimes K^{\flat}\rightarrow M\otimes K^{\flat}$. As a subcategory of $\Mod_{A,\ast}^{\varphi}$, it is stable under tensor products, internal Homs, symmetric and exterior powers, and it has a neutral object of its own. Using the isomorphisms
\begin{align*}
\varphi^{\ast}(M_{1}\otimes M_{2})\otimes K^{\flat} & \simeq \left(\varphi^{\ast}(M_{1})\otimes K^{\flat}\right)\otimes_{K^{\flat}}\left(\varphi^{\ast}(M_{2})\otimes K^{\flat}\right),\\
\varphi^{\ast}\left(\Hom_{\mathcal{O}_{K}^{\flat}}(M_{1},M_{2})\right)\otimes K^{\flat} & \simeq \Hom_{K^{\flat}}\left(\varphi^{\ast}(M_{1})\otimes K^{\flat},\varphi^{\ast}(M_{2})\otimes K^{\flat}\right),\\
\varphi^{\ast}(\mathcal{O}_{K}^{\flat})\otimes K^{\flat} & \simeq K^{\flat},
\end{align*}
the tensor products, internal Homs and neutral object in $\Mod_{R,f}^{\varphi}$ are given by
\begin{align*}
(M_{1},\varphi_{1})\otimes(M_{2},\varphi_{2}) & \eqd \left(M_{1}\otimes M_{2},\varphi_{1}\otimes\varphi_{2}\right),\\
\Hom\left((M_{1},\varphi_{1}),(M_{2},\varphi_{2})\right) & \eqd \left(\Hom_{\mathcal{O}_{K}^{\flat}}(M_{1},M_{2}),\Hom_{K^{\flat}}(\varphi_{1}^{-1},\varphi_{2})\right),\\
\mathcal{O}_{K}^{\flat} & \eqd (\mathcal{O}_{K}^{\flat},\mathrm{Id}).
\end{align*}
The rank and degree functions on $\Mod_{A,t}^{\varphi}$ induce rank and degree functions on $\Mod_{R,f}^{\varphi}$, and the corresponding Harder--Narasimhan (Fargues) filtrations $\mathcal{F}_{F}$ are compatible since the essential image of $\Mod_{R,f}^{\varphi}\hookrightarrow\Mod_{A,t}^{\varphi}$ is stable under strict subobjects. The rank of $(M,\varphi_{M})\in\Mod_{R,f}^{\varphi}$ is the usual rank of the finite free $\mathcal{O}_{K}^{\flat}$-module $M$, and its degree is the degree of the Hodge $\mathbb{R}$-filtration
\[
\mathcal{F}_{H}(M,\varphi_{M})\eqd\mathcal{F}\left(M,\varphi_{M}\left(\varphi^{\ast}M\right)\right)
\]
which is induced by the $\mathcal{O}_{K}^{\flat}$-lattice $\varphi_{M}(\varphi^{\ast}M)$ of $M\otimes K^{\flat}$ on the residue $M\otimes\mathbb{F}$ of $M$. The Hodge type of $(M,\varphi_{M})$ is the type $t_{H}(M,\varphi_{M})$ of $\mathcal{F}_{H}(M,\varphi_{M})$, so that
\[
t_{H}(M,\varphi_{M})=\mathbf{d}\left(M,\varphi_{M}(\varphi^{\ast}M)\right)\quad\text{in}\quad\mathbb{R}_{\geq}^{r}
\]
where $r=\rank\,M$. The next proposition then follows from~\cite[\S 5.3]{Co16}:

\begin{prop}\label{prop:FarguesonO_Kflattensor}
The restriction of the Fargues filtration to the subcategory $\Mod_{R,f}^{\varphi}$ of $\Mod_{A,t}^{\varphi}$ is compatible with tensor products, duals, symmetric and exterior powers.
\end{prop}

\begin{prop}\label{prop:HodgeonO_Kflat}
The Hodge filtration $\mathcal{F}_{H}:\Mod_{R,f}^{\varphi}\rightarrow\Fil_{\mathbb{F}}^{\mathbb{R}}$ is compatible with tensor products, duals, symmetric and exterior powers. For every exact sequence
\[
0\rightarrow(M_{1},\varphi_{1})\rightarrow(M_{2},\varphi_{2})\rightarrow(M_{3},\varphi_{3})\rightarrow0
\]
in $\Mod_{R,f}^{\varphi}$ with $r_{i}=\rank M_{i}$ (so that $r_{2}=r_{1}+r_{3}$), we have
\[
t_{H}(M_{1},\varphi_{1})\ast t_{H}(M_{3},\varphi_{3})\leq t_{H}(M_{2},\varphi_{2})\quad\text{in}\quad\mathbb{R}_{\geq}^{r_{2}}
\]
with equality if and only if for every $\gamma\in\mathbb{R}$, the complex of $\mathbb{F}$-vector spaces
\[
0\rightarrow\mathcal{F}_{H}^{\gamma}(M_{1},\varphi_{1})\rightarrow\mathcal{F}_{H}^{\gamma}(M_{2},\varphi_{2})\rightarrow\mathcal{F}_{H}^{\gamma}(M_{3},\varphi_{3})\rightarrow0
\]
is exact.
\end{prop}

\begin{proof}
This follows from~\ref{subsec:TensorProd} and Lemma~\ref{lem:RelPos=000026ExcSeq}.
\end{proof}

\begin{coro}\label{cor:tF=000026tH4M1}
For every $(M,\varphi)$ in $\Mod_{R,f}^{\varphi}$ of rank $r\in\mathbb{N}$,
\[
t_{F}(M,\varphi)\leq t_{H}(M,\varphi)\quad\text{in}\quad\mathbb{R}_{\geq}^{r}.
\]
\end{coro}

\begin{proof}
Let $X^{\bullet}=\bigoplus_{\gamma}X^{\gamma}$ be the $\mathbb{R}$-graded object of $\Mod_{R,f}^{\varphi}$ attached to the Fargues filtration of $X=(M,\varphi)$. Then by Propositions~\ref{prop:ExactSeqInModt} and~\ref{prop:HodgeonO_Kflat},
\[
t_{F}(X)=t_{F}(X^{\bullet})=\ast_{\gamma}t_{F}(X^{\gamma})\quad\text{and}\quad t_{H}(X)\geq t_{H}(X^{\bullet})=\ast_{\gamma}t_{H}(X^{\gamma}).
\]
We may thus assume that $X$ is semi-stable, in which case the result is obvious since the concave polygons $t_{F}(X)$ and $t_{H}(X)$ have the same terminal points.
\end{proof}

Let $\mathcal{O}_{K}^{\flat}\{n\}:=A\{n\}\otimes\mathcal{O}_{K}^{\flat}$ and $\mathbb{F}\{n\}:=A\{n\}\otimes\mathbb{F}$, so that $M\{n\}=M\otimes\mathcal{O}_{K}^{\flat}\{n\}$ for every $M$ in $\Mod_{R,f}^{\varphi}$. The map $X\mapsto X\{n\}=X\otimes\mathbb{F}\{n\}$ then induces a bijection between $\mathbb{F}$-subspaces $X$ of $M\otimes\mathbb{F}$ and $X\{n\}$ of $M\{n\}\otimes\mathbb{F}=M\otimes\mathbb{F}\{n\}$.

\begin{prop}\label{prop:tH(twist)4O_K^flat}
For every $M\in\Mod_{R,f}^{\varphi}$ of rank $r\in\mathbb{N}$ and any $n\in\mathbb{Z}$,
\[
\mathcal{F}_{H}^{\gamma}(M\{n\})=\mathcal{F}_{H}^{\gamma-n}(M)\{n\}\quad\text{inside}\quad M\{n\}\otimes\mathbb{F}
\]
for every $\gamma\in\mathbb{R}$, hence
\[
t_{H}(M\{n\})=t_{H}(M)+(n,\dots,n)\quad\text{in}\quad\mathbb{R}_{\geq}^{r}.
\]
\end{prop}
\begin{proof}
By definition, $\mathcal{F}_{H}^{\gamma}(M\{n\})$ equals
\[
\frac{M\otimes\mathcal{O}_{K}^{\flat}\{n\}\cap\left(I^{\gamma}\cdot\varphi_{M}\left(\varphi^{\ast}M\right)\otimes(\xi^{\prime}\bmod\pi)^{-n}\mathcal{O}_{K}^{\flat}\{n\}\right)+\mathfrak{m}_{K}^{\flat}\cdot M\otimes\mathcal{O}_{K}^{\flat}\{n\}}{\mathfrak{m}_{K}^{\flat}\cdot M\otimes\mathcal{O}_{K}^{\flat}\{n\}}
\]
where $I^{\gamma}=\left\{ x\in K^{\flat}:\left|x\right|\leq q^{-\gamma}\right\} $, and $\xi^{\prime}\bmod\pi=\varpi^{q}$, i.e.
\begin{align*}
\mathcal{F}_{H}^{\gamma}(M\{n\}) & = \frac{M\cap\left(I^{\gamma}\varpi^{-qn}\cdot\varphi_{M}\left(\varphi^{\ast}M\right)\right)+\mathfrak{m}_{K}^{\flat}\cdot M}{\mathfrak{m}_{K}^{\flat}\cdot M}\otimes\mathbb{F}\{n\}\\
& = \mathcal{F}_{H}^{\gamma-n}(M)\{n\}
\end{align*}
since $\left|\varpi^{-qn}\right|=q^{n}$ by our convention in~\ref{subsec:TheRings} for the normalization of the absolute value $|{\,\cdot\,}|$ on $\mathcal{O}_{K}^{\flat}$. This proves the proposition.
\end{proof}

\subsubsection{\texorpdfstring{$R=L$}{R=L}}\label{subsec:RisL}

Here, $\Mod_{L,\ast}^{\varphi}=\Mod_{L,f}^{\varphi}$ is the $E$-linear tannakian category of $E$-isocrystals over the perfect field $\mathbb{F}$, i.e.~finite dimensional vector spaces $D$ over $L$ equipped with an isomorphism $\varphi_{D}:\varphi^{\ast}D\rightarrow D$, or equivalently, with a $\varphi$-semilinear automorphism $\phi_{D}:D\rightarrow D$. The Dieudonn\'e-Manin classification~\cite{Ma63,De86}, extended from algebraically closed to perfect fields in~\cite{Zi84} (but with coefficients in $E$ instead of $\mathbb{Q}_{p}$ as in~\cite{Ko85} or~\cite{FaFo15}), gives a slope decomposition~\cite[VIII \S 1]{DaOrRa10}
\[
(D,\varphi_{D})=\bigoplus_{\lambda\in\mathbb{Q}}(D_{\lambda},\varphi_{D_{\lambda}}).
\]
For $\lambda=\frac{d}{h}$ with $d\in\mathbb{Z}$ and $h\in\mathbb{N}^{\ast}$ relatively prime, $D_{\lambda}$ is the union of the finitely generated $\mathcal{O}_{L}$-submodules $X$ of $D$ such that $\phi_{D}^{h}(X)=\pi^{d}X$. This Newton decomposition is functorial, compatible with all tensor product constructions, thus
\[
\mathcal{G}_{N}:\Mod_{L,f}^{\varphi}\rightarrow\Gra_{L}^{\mathbb{Q}},\qquad\mathcal{G}_{N}^{\lambda}(D,\varphi_{D})\eqd D_{\lambda}
\]
is an exact $\otimes$-functor, and so are the corresponding opposed Newton $\mathbb{Q}$-filtrations
\[
\mathcal{F}_{N},\mathcal{F}_{N}^{\iota}:\Mod_{L,f}^{\varphi}\rightarrow\Fil_{L}^{\mathbb{Q}}
\]
which are given by the usual formulas
\[
\mathcal{F}_{N}^{\lambda}(D,\varphi_{D})\eqd\bigoplus_{\lambda^{\prime}\geq\lambda}D_{\lambda^{\prime}}\quad\text{and}\quad\mathcal{F}_{N}^{\iota\lambda}(D,\varphi_{D})\eqd\bigoplus_{\lambda^{\prime}\geq\lambda}D_{-\lambda^{\prime}}.
\]
We denote by $t_{N}(D,\varphi_{D})$ and $t_{N}^{\iota}(D,\varphi_{D})$ the corresponding opposed types. Both Newton filtrations are Harder--Narasimhan filtrations, for the obvious rank function on $\Mod_{L,f}^{\varphi}$ and for the opposed degree functions which are respectively given by
\[
\deg_{N}(D,\varphi_{D})\eqd\deg\mathcal{F}_{N}(D,\varphi_{D})\quad\text{and}\quad\deg_{N}^{\iota}(D,\varphi_{D})\eqd\deg\mathcal{F}_{N}^{\iota}(D,\varphi_{D}).
\]
These degree functions are $\mathbb{Z}$-valued! If the residue field $\mathbb{F}$ is algebraically closed, the category $\Mod_{L,f}^{\varphi}$ is even semi-simple, with one simple object $D_{\lambda}^{\circ}$ for each slope $\lambda\in\mathbb{Q}$. If $\lambda=\frac{d}{h}$ as above, then $\rank(D_{\lambda}^{\circ})=h$ and $\deg_{N}(D_{\lambda}^{\circ})=d=-\deg_{N}^{\iota}(D_{\lambda}^{\circ})$.

Since $\varphi_{A\{1\}}(\varphi^{\ast}A\{1\})=\xi^{\prime-1}A\{1\}$ and $\xi'$ maps to a uniformizer in $L$, we have $\phi_{L\{1\}}(\mathcal{O}_{L}\{1\})=\pi^{-1}\mathcal{O}_{L}\{1\}$ for the $\mathcal{O}_{L}$-lattice $\mathcal{O}_{L}\{1\}:=A\{1\}\otimes\mathcal{O}_{L}$ in the Tate object $L\{1\}:=A\{1\}\otimes L$ of $\Mod_{L,f}^{\varphi}$. It follows that
\[
\deg_{N}(L\{1\})=-1\quad\text{and}\quad\deg_{N}^{\iota}(L\{1\})=+1.
\]
For $D$ in $\Mod_{L,f}^{\varphi}$ and $n\in\mathbb{Z}$, we set $D\{n\}:=D\otimes L\{1\}^{\otimes n}$ as usual. Then
\[
\mathcal{G}_{N}^{\gamma}(D\{n\})=\mathcal{G}_{N}^{\gamma+n}(D)\{n\}
\]
for every $\gamma\in\mathbb{Q}$, therefore
\[
\mathcal{F}_{N}^{\gamma}(D\{n\})=\mathcal{F}_{N}^{\gamma+n}(D)\{n\}\quad\text{and}\quad\mathcal{F}_{N}^{\iota\gamma}(D\{n\})=\mathcal{F}_{N}^{\iota\gamma-n}(D)\{n\}.
\]
In particular, we have the following equalities in $\mathbb{Q}_{\geq}^{r}$ where $r=\dim_{L}D$:
\[
t_{N}(D\{n\})=t_{N}(D)-(n,\dots,n)\quad\text{and}\quad t_{N}^{\iota}(D\{n\})=t_{N}^{\iota}(D)+(n,\dots,n).
\]

\subsubsection{\texorpdfstring{${R=\mathcal{O}_{L}}$}{R=OL}}\label{subsec:RisOL}

The category $\Mod_{\mathcal{O}_{L},f}^{\varphi}$ is now the category of $\mathcal{O}_{E}$-crystals over $\mathbb{F}$, or $\mathcal{O}_{L}$-lattices in $E$-isocrystals over $\mathbb{F}$, whose objects are finite free $\mathcal{O}_{L}$-modules $M$ equipped with an isomorphism $\varphi_{M}:\varphi^{\ast}M\otimes L\rightarrow M\otimes L$. It is a quasi-abelian $\mathcal{O}_{E}$-linear rigid $\otimes$-category, with an exact faithful $\otimes$-functor $-\otimes L:\Mod_{\mathcal{O}_{L},f}^{\varphi}\rightarrow\Mod_{L,f}^{\varphi}$. Since $\mathcal{O}_{L}$ is a discrete valuation ring, there is also a Hodge $\mathbb{Z}$-filtration, defined by
\[
\mathcal{F}_{H}(M,\varphi_{M})\eqd\mathcal{F}\left(M,\varphi_{M}(\varphi^{\ast}M)\right),
\]
a $\mathbb{Z}$-filtration on $M\otimes\mathbb{F}$, whose type will be denoted by $t_{H}(M,\varphi_{M})$, so that
\[
t_{H}(M,\varphi_{M})=\mathbf{d}\left(M,\varphi_{M}(\varphi^{\ast}M)\right)\quad\text{in}\quad\mathbb{Z}_{\geq}^{r}
\]
where $r=\rank(M)$. As before, we have the following proposition:

\begin{prop}\label{prop:HodgeonO_L}
The Hodge filtration $\mathcal{F}_{H}:\Mod_{\mathcal{O}_{L},f}^{\varphi}\rightarrow\Fil_{\mathbb{F}}^{\mathbb{Z}}$ is compatible with tensor products, duals, symmetric and exterior powers. For every exact sequence
\[
0\rightarrow(M_{1},\varphi_{1})\rightarrow(M_{2},\varphi_{2})\rightarrow(M_{3},\varphi_{3})\rightarrow0
\]
in $\Mod_{\mathcal{O}_{L},f}^{\varphi}$ with $r_{i}=\rank\,M_{i}$ (so that $r_{2}=r_{1}+r_{3}$), we have
\[
t_{H}(M_{1},\varphi_{1})\ast t_{H}(M_{3},\varphi_{3})\leq t_{H}(M_{2},\varphi_{2})\quad\text{in}\quad\mathbb{Z}_{\geq}^{r_{2}}
\]
with equality if and only if for every $\gamma\in\mathbb{Z}$, the complex of $\mathbb{F}$-vector spaces
\[
0\rightarrow\mathcal{F}_{H}^{\gamma}(M_{1},\varphi_{1})\rightarrow\mathcal{F}_{H}^{\gamma}(M_{2},\varphi_{2})\rightarrow\mathcal{F}_{H}^{\gamma}(M_{3},\varphi_{3})\rightarrow0
\]
is exact.
\end{prop}

\begin{coro}[Mazur's inequality]\label{cor:Mazurinequality}
For every $X$ in $\Mod_{\mathcal{O}_{L},f}^{\varphi}$ of rank $r\in\mathbb{N}$,
\[
t_{N}^{\iota}(X\otimes L)\leq t_{H}(X)\quad\text{in}\quad\mathbb{Q}_{\geq}^{r}.
\]
\end{coro}

\begin{proof}
We first show that $\mathcal{F}_{N}^{\iota}(X\otimes L)$ and $\mathcal{F}_{H}(X)$ have the same degree. Since both filtrations are compatible with exterior powers, we may assume that the rank of $X=(M,\varphi_{M})$ equals $1$. Then $\varphi_{M}(\varphi^{\ast}M)=\pi^{-d}M$ for some $d\in\mathbb{Z}$, thus indeed $\deg\mathcal{F}_{N}^{\iota}(X\otimes L)=d=\deg\mathcal{F}_{H}(X)$. Returning to the general case, both polygons thus have the same terminal points. We now follow the proof of Corollary~\ref{cor:tF=000026tH4M1}. Let $X^{\bullet}=\bigoplus_{\gamma}X^{\gamma}$ be the $\mathbb{Q}$-graded object of $\Mod_{\mathcal{O}_{L},f}^{\varphi}$ attached to the filtration on $X$ induced by $\mathcal{F}_{N}^{\iota}(X\otimes L)$. Then by exactness of $\mathcal{F}_{N}^{\iota}$ and the previous proposition
\[
t_{N}^{\iota}(X\otimes L)=t_{N}^{\iota}(X^{\bullet}\otimes L)=\ast_{\gamma}t_{N}^{\iota}(X^{\gamma}\otimes L)\quad\text{and}\quad t_{H}(X)\geq t_{H}(X^{\bullet})=\ast_{\gamma}t_{H}(X^{\gamma}).
\]
We may thus assume that $X\otimes L$ is semi-stable~(i.e.~isoclinic), in which case the result is obvious since $t_{N}^{\iota}(X\otimes L)$ and $t_{H}(X)$ have the same terminal points.
\end{proof}

\begin{rema}
For another (metric) proof of a generalization of the corollary, see~\cite[Th\'eor\`eme~7]{CoNi16} when $\mathbb{F}$ is algebraically closed and~\cite[\S 3.2]{Co15} in the general case.
\end{rema}

We have already defined the Tate object $\mathcal{O}_{L}\{1\}=A\{1\}\otimes\mathcal{O}_{L}$, giving rise to Tate twists $M\{n\}:=M\otimes\mathcal{O}_{L}\{1\}^{\otimes n}$ for every $M\in\Mod_{\mathcal{O}_{L},f}^{\varphi}$ and $n\in\mathbb{Z}$, with a bijection~$X\mapsto X\{n\}:=X\otimes\mathbb{F}\{n\}$ between $\mathbb{F}$-subspaces of $M\otimes\mathbb{F}$ and $M\{n\}\otimes\mathbb{F}=M\otimes\mathbb{F}\{n\}$.

\begin{prop}\label{prop:tH(twist)4O_L}
For every $M\in\Mod_{\mathcal{O}_{L},f}^{\varphi}$ of rank $r\in\mathbb{N}$ and $n\in\mathbb{Z}$,
\[
\mathcal{F}_{H}^{\gamma}(M\{n\})=\mathcal{F}_{H}^{\gamma-n}(M)\{n\}\quad\text{ inside}\quad M\{n\}\otimes\mathbb{F}
\]
for every $\gamma\in\mathbb{Z}$, hence
\[
t_{H}(M\{n\})=t_{H}(M)+(n,\dots,n)\quad\text{in}\quad\mathbb{Z}_{\geq}^{r}.
\]
\end{prop}

\begin{proof}
This is similar to Proposition~\ref{prop:tH(twist)4O_K^flat}. It also follows from the compatibility of the Hodge filtration with tensor products (Proposition~\ref{prop:HodgeonO_L}), along with the formula
\[
\varphi_{L\{1\}}(\varphi^{\ast}\mathcal{O}_{L}\{1\})=\pi^{-1}\mathcal{O}_{L}\{1\}
\]
which shows that $\mathcal{F}_{H}(\mathcal{O}_{L}\{1\})$ has a single jump at $1$.
\end{proof}

\subsubsection{\texorpdfstring{${R=A(K)=W_{\mathcal{O}_{E}}(K^{\flat})}$}{R=A(K)=WOE(K flat)}}\label{subsec:RisW(Kflat)}

Then $\Mod_{R,\ast}^{\varphi}$ is the abelian $\otimes$-category of finitely generated $R$-modules $M$ equipped with an isomorphism $\varphi_{M}:\varphi^{\ast}M\rightarrow M$, or equivalently, with a $\varphi$-semilinear automorphism $\phi_{M}:M\rightarrow M$. If $\overline{K}^{\flat}$ is an algebraic closure of $K^{\flat}$ with Galois group $\Gamma=\Gal(\overline{K}^{\flat}/K^{\flat})$ and $\overline{R}=W_{\mathcal{O}_{E}}(\overline{K}^{\flat})$, the formulas
\begin{align*}
(M,\varphi_{M}) & \mapsto (T,\rho)=\left(\left(M\otimes_{R}\overline{R}\right)^{\phi_{M}\otimes\varphi=1},\mathrm{Id}\otimes\rho_{\overline{R}}\right)\\
(T,\rho) & \mapsto (M,\varphi_{M})=\left(\left(T\otimes_{\mathcal{O}_{E}}\overline{R},\rho\otimes\rho_{\overline{R}}\right)^{\Gamma},\mathrm{Id}\otimes\varphi\right)
\end{align*}
yield equivalences of $\otimes$-categories between $\Mod_{R,\ast}^{\varphi}$ and the category $\Rep_{\mathcal{O}_{E},\ast}(\Gamma$) of continuous representations $\rho:\Gamma\rightarrow\Aut_{\mathcal{O}_{E}}(T)$ on finitely generated $\mathcal{O}_{E}$-modules $T$~\cite[1.2.6]{Fo90}. Here $\rho_{\overline{R}}:\Gamma\rightarrow\Aut_{R,\varphi}(\overline{R})$ is induced by the functoriality of $W_{\mathcal{O}_{E}}({\,\cdot\,})$.

\subsubsection{\texorpdfstring{${{R=A[\frac{1}{\pi}]}}$}{R=A[1/pi]}}\label{subsec:RisA=00005B1/pi=00005D}

The category $\Mod_{R,\ast}^{\varphi}=\Mod_{R,f}^{\varphi}$ is the rigid $E$-linear $\otimes$-category of finite free $A[\frac{1}{\pi}]$-modules $M$ with an isomorphism $\varphi_{M}:\varphi^{\ast}M[\xi^{\prime-1}]\rightarrow M[\xi^{\prime-1}]$. Since $\xi'=\varphi(\xi)$, the Frobenius of $A$ induces an isomorphism $\varphi:B_{\dR}^{+}\rightarrow B_{\dR}^{\prime+}$ of discrete valuation rings between the completion of the local rings of $A[\frac{1}{\pi}]$ at the maximal ideals $A[\frac{1}{\pi}]\xi=\ker(\theta:A[\frac{1}{\pi}]\twoheadrightarrow K)$ and $A[\frac{1}{\pi}]\xi'=\varphi(A[\frac{1}{\pi}]\xi)$, along with the induced isomorphisms $\varphi:K\rightarrow K'$ and $\varphi:B_{\dR}\rightarrow B_{\dR}^{\prime}$ between the residue and fraction fields of $B_{\dR}^{+}$ and $B_{\dR}^{\prime+}$. For $(M,\varphi_{M})$ in $\Mod_{R,f}^{\varphi}$, the commutative diagram
\[
\xymatrix@C=4pc{M\left[\xi^{-1}\right]\ar[d]_{\varphi}\ar[r]^{(\varphi^{-1})^{\ast}(\varphi_{M})} & ((\varphi^{-1})^{\ast}M)\left[\xi^{-1}\right]\ar[d]^{\varphi}\\
(\varphi^{\ast}M)\left[\xi^{\prime-1}\right]\ar[r]^{\varphi_{M}} & M\left[\xi^{\prime-1}\right] }
\]
extends to a commutative diagram
\[
\xymatrix@C=4pc{M\otimes B_{\dR}\ar[d]_{\varphi\otimes\varphi}\ar[r]^{(\varphi^{-1})^{\ast}(\varphi_{M})} & ((\varphi^{-1})^{\ast}M)\otimes B_{\dR}\ar[d]^{\varphi\otimes\varphi}\\
(\varphi^{\ast}M)\otimes B_{\dR}^{\prime}\ar[r]^{\varphi_{M}} & M\otimes B_{\dR}^{\prime} }
\]
Then $M\otimes B_{\dR}^{+}$ is a $B_{\dR}^{+}$-lattice in $M\otimes B_{\dR}$ and similarly for the other three vertices. Each line of our diagram thus yields a pair of $\mathbb{Z}$-filtrations on the residue (over $K$ or $K'$) of its vertices, which have opposed types in $\mathbb{Z}_{\geq}^{r}$ where $r$ is the rank of $M$, and the two pairs match along the $\varphi$-equivariant isomorphisms which are induced by the vertical maps. In particular, the Hodge $\mathbb{Z}$-filtrations
\begin{align*}
\mathcal{F}_{H}^{\iota}(M,\varphi_{M}) & \eqd \mathcal{F}\left(M\otimes B_{\dR}^{+}\,,\left((\varphi^{-1})^{\ast}(\varphi_{M})\right)^{-1}\left(\left((\varphi^{-1})^{\ast}M\right)\otimes B_{\dR}^{+}\right)\right)\\
\text{and}\quad\mathcal{F}_{H}(M,\varphi_{M}) & \eqd \mathcal{F}\left(M\otimes B_{\dR}^{\prime+}\,,\varphi_{M}\left((\varphi^{\ast}M)\otimes B_{\dR}^{\prime+}\right)\right)
\end{align*}
on respectively $M\otimes_{A}K$ and $M\otimes_{A}K'$ have opposed types
\[
t_{H}^{\iota}(M,\varphi_{M})\quad\text{and}\quad t_{H}(M,\varphi_{M})\quad\text{in}\quad\mathbb{Z}_{\geq}^{r}.
\]

\begin{prop}\label{prop:Hodge4A=00005B1/pi=00005D_tensorexact}
The Hodge filtration $\mathcal{F}_{H}:\Mod_{R,f}^{\varphi}\rightarrow\Fil_{K^{\prime}}^{\mathbb{Z}}$ is compatible with tensor products, duals, symmetric and exterior powers. For every exact sequence
\[
0\rightarrow(M_{1},\varphi_{1})\rightarrow(M_{2},\varphi_{2})\rightarrow(M_{3},\varphi_{3})\rightarrow0
\]
in $\Mod_{R,f}^{\varphi}$ with $r_{i}=\rank M_{i}$ (so that $r_{2}=r_{1}+r_{3}$), we have
\[
t_{H}(M_{1},\varphi_{1})\ast t_{H}(M_{3},\varphi_{3})\leq t_{H}(M_{2},\varphi_{2})\quad\text{in}\quad\mathbb{Z}_{\geq}^{r_{2}}
\]
with equality if and only if for every $\gamma\in\mathbb{R}$, the complex of $K'$-vector spaces
\[
0\rightarrow\mathcal{F}_{H}^{\gamma}(M_{1},\varphi_{1})\rightarrow\mathcal{F}_{H}^{\gamma}(M_{2},\varphi_{2})\rightarrow\mathcal{F}_{H}^{\gamma}(M_{3},\varphi_{3})\rightarrow0
\]
is exact. The Hodge filtration $\mathcal{F}_{H}^{\iota}:\Mod_{R,f}^{\varphi}\rightarrow\Fil_{K}^{\mathbb{Z}}$ has analogous properties.
\end{prop}

For the Tate object $A[\frac{1}{\pi}]\{1\}:=A\{1\}[\frac{1}{\pi}]$, $t_{H}=1=-t_{H}^{\iota}$, thus again:

\begin{prop}\label{prop:t_H(twist)4A=00005B1/pi=00005D}
For every $M$ in $\Mod_{R,f}^{\varphi}$ of rank $r\in\mathbb{N}$ and $n\in\mathbb{Z}$,
\[
\mathcal{F}_{H}^{\gamma}(M\{n\})=\mathcal{F}_{H}^{\gamma-n}(M)\{n\}\quad\text{and}\quad\mathcal{F}_{H}^{\iota\gamma}(M\{n\})=\mathcal{F}_{H}^{\iota\gamma+n}(M)\{n\}
\]
for every $\gamma\in\mathbb{Z}$, therefore
\[
\begin{aligned}
t_{H}(M\{n\}) & = t_{H}(M)+(n,\dots,n)\\
t_{H}^{\iota}(M\{n\}) & = t_{H}^{\iota}(M)-(n,\dots,n)
\end{aligned}\quad\text{in}\quad\mathbb{Z}_{\geq}^{r}.
\]
\end{prop}

The $\otimes$-functor $\Mod_{A,\ast}^{\varphi}\rightarrow\Mod_{R,f}^{\varphi}$ identifies the isogeny category $\Mod_{A,\ast}^{\varphi}\otimes E$ with a full subcategory of $\Mod_{R,f}^{\varphi}$. We may thus unambiguously denote by $X\mapsto X\otimes E$ or $X[\frac{1}{\pi}]$ the $\otimes$-functor from $\Mod_{A,\ast}^{\varphi}$ to either $\Mod_{A,\ast}^{\varphi}\otimes E$ or $\Mod_{R,f}^{\varphi}$.

\begin{prop}\label{prop:CompHodgeA=00005B1/pi=00005DOLA1}
For a finite free BKF-module $M$ of rank $r\in\mathbb{N}$,
\[
\begin{aligned}
t_{H}(M\otimes\mathcal{O}_{L}) & \leq t_{H}\left(M\otimes E\right) && \text{in}\quad\mathbb{Z}_{\geq}^{r},\\
t_{H}(M\otimes\mathcal{O}_{K}^{\flat}) & \leq t_{H}\left(M\otimes E\right) && \text{in}\quad\mathbb{R}_{\geq}^{r}.
\end{aligned}
\]
\end{prop}

\begin{proof}
Using the compatibility with Tate twists (Propositions~\ref{prop:tH(twist)4O_K^flat}, \ref{prop:tH(twist)4O_L},\linebreak and~\ref{prop:t_H(twist)4A=00005B1/pi=00005D}), we may assume that $M\subset M'=\varphi_{M}(\varphi^{\ast}M)$ in $M[\xi^{\prime-1}]$. Then $Q=M'/M$ is a perfect $A$-module of projective dimension $\leq1$ which is killed by a power of $\xi^{\prime}$, say $\xi^{\prime n}Q=0$. For $0\leq i\leq n$, let $M^{i}$ be the inverse image of $Q^{i}=Q[\xi^{\prime i}]$ in $M'$, so~that
\[
M=M^{0}\subset M^{1}\subset\dots\subset M^{n}=M'\quad\text{with}\quad M^{i}/M^{0}=Q^{i}.
\]
We first claim that each $M^{i}$ is finite free over $A$. By descending induction on $i$, it is sufficient to establish that the following $A$-module has projective dimension~$1$:
\[
X^{i}=M^{i}/M^{i-1}\simeq Q^{i}/Q^{i-1}\simeq\xi^{\prime i-1}Q[\xi^{\prime i}]\subset Q[\xi^{\prime}]\subset Q.
\]
We will show that it is finite free over $A(1):=A/A\xi^{\prime}$. Since $A(1)\simeq A/A\xi\simeq\mathcal{O}_{K}$ is a valuation ring, we just have to verify that $X^{i}$ is finitely generated and torsion-free over $A(1)$. Since $Q$ is finitely presented over $A$, it is finitely presented over $A(n):=A/A\xi^{\prime n}$, which is a coherent ring by~\cite[3.24]{BaMoSc16}, thus $Q^{i}=Q[\xi^{\prime i}]$ is finitely presented over $A(n)$ and $A$ for all $i$, and so is $X^{i}\simeq Q^{i}/Q^{i-1}$. On the other hand, $Q[\mathfrak{m}^{\infty}]=0$ by~\ref{subsec:ProjDimModA}, thus also $X^{i}[\mathfrak{m}^{\infty}]=0$, which means that $X^{i}$ is indeed torsion-free as an $A(1)$-module. We denote by $x_{i}$ the rank of $X^{i}$ over $A(1)$.

Let $S$ be any one of the valuations rings $B_{\dR}^{\prime+}$, $\mathcal{O}_{K}^{\flat}$ or $\mathcal{O}_{L}$. Then $\Tor_{1}^{A}(X^{i},S)=0$ since $\Tor_{1}^{A}(A(1),S)=S[\xi^{\prime}]=0$. We thus obtain a sequence of $S$-lattices
\[
M\otimes S=M^{0}\otimes S\subset M^{1}\otimes S\subset\dots\subset M^{n}\otimes S=M^{\prime}\otimes S
\]
inside $M\otimes\Frac(S)$. The triangular inequality of Lemma~\ref{lem:triangIneq} then yields
\[
\mathbf{d}\left(M\otimes S,M'\otimes S\right)\leq\sum_{i=1}^{n}\mathbf{d}\left(M^{i-1}\otimes S,M^{i}\otimes S\right)\quad\text{in}\quad\mathbb{R}_{\geq}^{r}.
\]
Since $M^{i}\otimes S/M^{i-1}\otimes S\simeq X^{i}\otimes S\simeq(S/\xi_{S}^{\prime}S)^{x_{i}}$ where $\xi_{S}^{\prime}$ is the image of $\xi^{\prime}$ in $S$ and since also $\left|\xi_{S}^{\prime}\right|=q^{-1}$ in all three cases for the normalized absolute value on $S$,
\[
\mathbf{d}\left(M^{i-1}\otimes S,M^{i}\otimes S\right)=(1,\dots,1,0,\dots,0)\quad\text{in}\quad\mathbb{Z}_{\geq}^{r}\subset\mathbb{R}_{\geq}^{r}
\]
with exactly $x_{i}$ one's. Now observe that by definition of our various Hodge types,
\[
\mathbf{d}\left(M\otimes S,M'\otimes S\right)=
\begin{cases}
t_{H}\left(M\otimes E\right) & \text{for }S=B_{\dR}^{\prime+},\\
t_{H}(M\otimes\mathcal{O}_{L}) & \text{for }S=\mathcal{O}_{L},\\
t_{H}(M\otimes\mathcal{O}_{K}^{\flat}) & \text{for }S=\mathcal{O}_{K}^{\flat}.
\end{cases}
\]
To establish the proposition, it is now sufficient to show that for $S=B_{\dR}^{\prime+}$, actually
\[
\mathbf{d}\left(M\otimes S,M'\otimes S\right)=\sum_{i=1}^{n}\mathbf{d}\left(M^{i}\otimes S,M^{i-1}\otimes S\right)\quad\text{in}\quad\mathbb{Z}_{\geq}^{r}.
\]
Being the completion of the Noetherian local ring $A_{(\xi')}$ of $A$, $S$ is flat over $A$, thus
\[
\frac{M^{i}\otimes S}{M^{0}\otimes S}=Q[\xi^{\prime i}]\otimes S=Q\otimes S[\xi_{S}^{\prime i}]=\frac{M'\otimes S}{M\otimes S}[\xi_{S}^{\prime i}],
\]
which means that $M^{i}\otimes S=(M^{\prime}\otimes S)\cap\xi_{S}^{\prime-i}(M\otimes S)$ in $M\otimes B_{\dR}^{\prime}$. If
\[
\mathbf{d}\left(M\otimes S,M'\otimes S\right)=(n_{1}\geq\dots\geq n_{r})\quad\text{in}\quad\mathbb{Z}_{\geq}^{r},
\]
there exists an $S$-basis $(e_{1},\dots,e_{r})$ of $M\otimes S$ such that $(\xi_{S}^{\prime-n_{1}}e_{1},\dots,\xi_{S}^{\prime-n_{r}}e_{r})$ is an $S$-basis of $M^{\prime}\otimes S$. Then $(\xi_{S}^{\prime-\min(n_{1},i)}e_{1},\dots,\xi_{S}^{\prime-\min(n_{r},i)}e_{r})$ is an $S$-basis of $M^{i}\otimes S$, $x_{i}=\max\{j:n_{j}\geq i\}$ and indeed $n_{j}=\sharp\{i:x_{i}\geq j\}$ for all $j\in\{1,\dots,r\}$.
\end{proof}
\begin{rema}
With notations as above (and for a finite free BKF-module $M$ such that $M^{\vee}$ is effective), the proof shows that we have equality when $n=1$, i.e.~$\xi'Q=0$, i.e.~$t_{H}(M\otimes E)=(1,\dots,1,0,\dots0)$ is minuscule. More generally for $S\in\{\mathcal{O}_{L},\mathcal{O}_{K}^{\flat}\}$, $t_{H}(M\otimes E)=t_{H}(M\otimes S)$ if $\Tor_{j}^{A}(S,Q/\xi^{\prime i}Q)=0$ for $1\leq i\leq n$ and $j\in\{1,2\}$.
\end{rema}

\begin{rema}
For a finite free BKF-module $M\in\Mod_{A,f}^{\varphi}$ of rank $r\in\mathbb{N}$, we thus have
\[
t_{F,\infty}(M)\leq t_{F}(M\otimes\mathcal{O}_{K}^{\flat})\leq t_{H}(M\otimes\mathcal{O}_{K}^{\flat})\leq t_{H}\!\left(M\!\left[{\frac{1}{\pi}}\right]\right)\geq t_{H}(M\otimes\mathcal{O}_{L})\geq t_{N}^{\iota}(M\otimes L)
\]
by Propositions~\ref{prop:ComptFinftytFn}, \ref{prop:CompHodgeA=00005B1/pi=00005DOLA1} and Corollaries~\ref{cor:tF=000026tH4M1} and~\ref{cor:Mazurinequality}. In particular,
\[
\begin{aligned}t_{F,\infty}(M)(r) & =\deg_{t}(M\otimes\mathcal{O}_{K}^{\flat})=\deg\mathcal{F}_{H}(M\otimes\mathcal{O}_{K}^{\flat})\\
& =\deg\mathcal{F}_{H}\!\left(M\!\left[\frac{1}{\pi}\right]\right)\\
& =\deg\mathcal{F}_{H}(M\otimes\mathcal{O}_{L})=\deg_{N}^{\iota}(M\otimes L)
\end{aligned}
\]
and this a priori real number actually belongs to $\mathbb{Z}$. We call it the degree of $M$.
\end{rema}


\section{The functors of Fargues}\label{sec:FunctorsOfFargues}

Suppose from now on that $K=C$ is algebraically closed. In this section, we will define and study the following commutative diagram of covariant $\otimes$-functors:
\[
\xymatrix{\Mod_{A,f}^{\varphi}\ar[rr]^{\HTT'}\ar@{^{(}->}[d] & & \HT_{\mathcal{O}_{E}}^{B_{\dR}}\ar@{^{(}->}[d]\\
\Mod_{A,f}^{\varphi}\otimes E\ar[r]^{\underline{\mathcal{E}}}\ar@{^{(}->}[d] & \Modif_{X}^{ad}\ar[r]^{\HTT}\ar@{^{(}->}[d] & \HT_{E}^{B_{\dR}}\\
\Mod_{A[\frac{1}{\pi}]}^{\varphi}\ar[r]^{\underline{\mathcal{E}}} & \Modif_{X} }
\]
In this diagram, the first two lines are equivalences of $\otimes$-categories, the top vertical arrows are faithful and the bottom ones fully faithful. The construction of $\underline{\mathcal{E}}$ which is given below is a covariant version of the analytic construction of~\cite{Fa15}. A slightly twisted version of it was sketched in Scholze's course~\cite{ScWe15} (for stukhas with one paw at $\mathfrak{m}=A\xi$). Our variant is meant to match the normalized construction of $\HTT'$ in~\cite{BaMoSc16}, where the paw was twisted from $\mathfrak{m}$ to $\mathfrak{m}'=A\xi'$. Following~\cite{BaMoSc16}, we fix a compatible system of $p$-power roots of unity, $\zeta_{p^{r}}\in\mathcal{O}_{C}^{\times}$ for $r\geq1$, and set
\begin{gather*}
\epsilon=(1,\zeta_{p},\zeta_{p^{2}},\dots)\in\mathcal{O}_{C}^{\flat,\times},\qquad\mu=[\epsilon]-1\in\mathfrak{m}\subset A,
\\
\xi={\frac{\mu}{\varphi^{-1}(\mu)}}=1+[\epsilon^{1/q}]+\dots+[\epsilon^{1/q}]^{q-1}\in\mathfrak{m}\subset A,
\\
\xi^{\prime}=\varphi(\xi)={\frac{\varphi(\mu)}{\mu}}=1+[\epsilon]+\dots+[\epsilon]^{q-1}\in\mathfrak{m}\subset A,
\\
\varpi=\xi\bmod\pi=1+\epsilon^{1/q}+\dots+(\epsilon^{1/q})^{q-1}\in\mathfrak{m}_{C}^{\flat}\subset\mathcal{O}_{C}^{\flat},
\\
\varpi^{q}=\xi^{\prime}\bmod\pi=1+\epsilon+\dots+\epsilon^{q-1}\in\mathfrak{m}_{C}^{\flat}\subset\mathcal{O}_{C}^{\flat}.
\end{gather*}
As suggested by the notations, $\xi$ is a generator of $\ker(\theta:A\twoheadrightarrow\mathcal{O}_{C})$. We have
\[
\varphi^{-1}(\mu)\mid\mu\mid\varphi(\mu)\,\,\text{in }A\quad\text{thus}\quad A\!\left[\frac{1}{\varphi^{-1}(\mu)}\right]\subset A\!\left[\frac{1}{\mu}\right]\subset A\!\left[\frac{1}{\varphi(\mu)}\right].
\]
Moreover, $\theta(\varphi^{-1}(\mu))=\zeta_{q}-1\neq0$, and therefore $\xi\nmid\varphi^{-1}(\mu)$ and $\xi^{\prime}\nmid\mu$.

\subsection{Modifications of vector bundles on the curve}

\subsubsection{The Fargues--Fontaine curve}\label{subsec:The-Fargues--Fontaine-curve}

Let $X=X_{C^{\flat},E}$ be the Fargues--Fontaine curve attached to $(C^{\flat},E)$~\cite{FaFo15}. This is an integral noetherian regular $1$-dimensional scheme over $E$ which is a complete curve in the sense of~\cite[5.1.3]{FaFo15}: the degree function on divisors factors through a degree function on the Picard group, $\deg:\Pic(X)\rightarrow\mathbb{N}$. We denote by $\eta$ the generic point of $X$ and by $E(X)=\mathcal{O}_{X,\eta}$ the field of rational functions on $X$. In addition, there is a distinguished closed point $\infty\in\left|X\right|$ whose completed local ring $\mathcal{O}_{X,\infty}^{\wedge}$ is canonically isomorphic to the ring $B_{\dR}^{+}$ of Section~\ref{subsec:RisA=00005B1/pi=00005D}, see also~\ref{subsec:CompHodge} below.

\subsubsection{Vector bundles on the curve}

Let $\Bun_{X}$ be the $E$-linear $\otimes$-category of vector bundles $\mathcal{E}$ on $X$. Since $X$ is a regular curve, it is a quasi-abelian category whose short exact sequences remain exact in the larger category of all sheaves on $X$, and the generic fiber $\mathcal{E}\mapsto\mathcal{E}_{\eta}$ yields an exact and faithful $\otimes$-functor
\[
({\,\cdot\,})_{\eta}:\Bun_{X}\rightarrow\Vect_{E(X)}
\]
which induces an isomorphism between the poset $\Sub(\mathcal{E})$ of strict subobjects of $\mathcal{E}$ in $\Bun_{X}$ and the poset $\Sub(\mathcal{E}_{\eta})$ of $E(X)$-subspaces of $\mathcal{E}_{\eta}$.

\subsubsection{Newton slope filtrations}\label{subsec:Newton4BundlesOnCurve}

The usual rank and degree functions
\[
\rank:\sk\,\Bun_{X}\rightarrow\mathbb{N}\quad\text{ and}\quad\deg:\sk\,\Bun_{X}\rightarrow\mathbb{Z}
\]
are additive on short exact sequences in $\Bun_{X}$, and they are respectively constant and non-decreasing on mono-epis in $\Bun_{X}$. More precisely, if $f:\mathcal{E}_{1}\rightarrow\mathcal{E}_{2}$ is a mono-epi, then $\rank(\mathcal{E}_{1})=\rank(\mathcal{E}_{2})$ and $\deg(\mathcal{E}_{1})\leq\deg(\mathcal{E}_{2})$ with equality if and only if $f$ is an isomorphism. These functions yield a Harder--Narasimhan filtration on $\Bun_{X}$, the Newton filtration $\mathcal{F}_{N}$ with slopes $\mu=\deg/\rank$ in $\mathbb{Q}$ (introduced in~\cite[5.5]{FaFo15}). The filtration $\mathcal{F}_{N}(\mathcal{E})$ on $\mathcal{E}\in\Bun_{X}$ is non-canonically split. More precisely for every $\mu\in\mathbb{Q}$, the full subcategory of semi-stable vector bundles of slope $\mu$ is abelian, equivalent to the category of right $D_{\mu}$-vector spaces, where $D_{\mu}$ is the semi-simple division $E$-algebra whose invariant is the class of $\mu$ in $\mathbb{Q}/\mathbb{Z}$. We denote by $\mathcal{O}_{X}(\mu)$ its unique simple object~\cite[5.6.22]{FaFo15}. Then for every vector bundle $\mathcal{E}$ on $X$, there is unique sequence $\mu_{1}\geq\dots\geq\mu_{s}$ in $\mathbb{Q}$ for which there is a (non-unique) isomorphism $\bigoplus_{i=1}^{s}\mathcal{O}_{X}(\mu_{i})\simeq\mathcal{E}$, and any such isomorphism maps $\bigoplus_{i:\mu_{i}\geq\gamma}\mathcal{O}_{X}(\mu_{i})$ to $\mathcal{F}_{N}^{\gamma}(\mathcal{E})$ for every $\gamma\in\mathbb{Q}$ (see~\cite[Chapter~$8$]{FaFo15}, particularly its main theorem~\cite[8.2.10]{FaFo15}). We denote by $t_{N}(\mathcal{E})\in\mathbb{Q}_{\geq}^{r}$ the type of $\mathcal{F}_{N}(\mathcal{E})$, where $r=\rank(\mathcal{E})$.

\begin{prop}
The Newton filtration is compatible with tensor products, duals, symmetric and exterior powers in $\Bun_{X}$. For any exact sequence in $\Bun_{X}$,
\[
0\rightarrow\mathcal{E}_{1}\rightarrow\mathcal{E}_{2}\rightarrow\mathcal{E}_{3}\rightarrow0
\]
set $r_{i}=\rank\,\mathcal{E}_{i}$ and view $t_{N}(\mathcal{E}_{i})$ as a concave function $f_{i}:[0,r_{i}]\rightarrow\mathbb{R}$. Then
\[
f_{1}\ast f_{3}(s)\geq f_{2}(s)\geq
\begin{cases}
f_{1}(s) & \text{if }0\leq s\leq r_{1}\\
f_{1}(r_{1})+f_{3}(s-r_{1}) & \text{if }r_{1}\leq s\leq r_{2}
\end{cases}
\]
with equality for $s=0$ and $s=r_{2}$. In particular,
\begin{gather*}
\begin{aligned}
t_{N}^{\max}(\mathcal{E}_{1}) & \leq t_{N}^{\max}(\mathcal{E}_{2}) \leq \max\left\{ t_{N}^{\max}(\mathcal{E}_{1}),t_{N}^{\max}(\mathcal{E}_{3})\right\},\\
t_{N}^{\min}(\mathcal{E}_{3}) & \geq t_{N}^{\min}(\mathcal{E}_{2}) \,\geq \min\left\{ t_{N}^{\min}(\mathcal{E}_{1}),t_{N}^{\min}(\mathcal{E}_{3})\right\},
\end{aligned}
\\
\text{and}\qquad t_{N}(\mathcal{E}_{2})\leq t_{N}(\mathcal{E}_{1})\ast t_{N}(\mathcal{E}_{3})\quad\text{in}\quad\mathbb{Q}_{\geq}^{r_{2}}.
\end{gather*}
Moreover, $t_{N}(\mathcal{E}_{2})=t_{N}(\mathcal{E}_{1})\ast t_{N}(\mathcal{E}_{3})$ if and only if for every $\gamma\in\mathbb{Q}$,
\[
0\rightarrow\mathcal{F}_{N}^{\gamma}(\mathcal{E}_{1})\rightarrow\mathcal{F}_{N}^{\gamma}(\mathcal{E}_{2})\rightarrow\mathcal{F}_{N}^{\gamma}(\mathcal{E}_{3})\rightarrow0
\]
is exact.
\end{prop}

\begin{proof}
The compatibility of $\mathcal{F}_{N}$ with $\otimes$-products and duals comes from~\cite[5.6.23]{FaFo15}. Since $\Bun_{X}$ is an $E$-linear category, the compatibility of $\mathcal{F}_{N}$ with symmetric and exterior powers follows from its additivity and compatibility with $\otimes$-products. For the remaining assertions, see~\cite[Proposition~21]{Co16} or~\cite[4.4.4]{An09}.
\end{proof}

\subsubsection{Modifications of vector bundles}

We denote by $\Modif_{X}$ the category of triples
\[
\underline{\mathcal{E}}=(\mathcal{E}_{1},\mathcal{E}_{2},f)
\]
where $\mathcal{E}_{1}$ and $\mathcal{E}_{2}$ are vector bundles on $X$ while $f$ is an isomorphism
\[
f:\mathcal{E}_{1}\vert_{X\setminus\{\infty\}}\rightarrow\mathcal{E}_{2}\vert_{X\setminus\{\infty\}}.
\]
A morphism $F:\underline{\mathcal{E}}\rightarrow\underline{\mathcal{E}}'$ is a pair of morphisms $F_{i}:\mathcal{E}_{i}\rightarrow\mathcal{E}_{i}^{\prime}$ with $F_{2}\circ f=f'\circ F_{1}$. This defines a quasi-abelian $E$-linear rigid $\otimes$-category with a Tate twist. The kernels and cokernels are induced by those of $\Bun_{X}$. The neutral object is the trivial modification $\underline{\mathcal{O}}_{X}=(\mathcal{O}_{X},\mathcal{O}_{X},\mathrm{Id})$, the tensor product and duals are given by
\[
\underline{\mathcal{E}}\otimes\underline{\mathcal{E}}^{\prime}\eqd(\mathcal{E}_{1}\otimes\mathcal{E}_{1}^{\prime},\mathcal{E}_{2}\otimes\mathcal{E}_{2}^{\prime},f\otimes f')\quad\text{and}\quad\underline{\mathcal{E}}^{\vee}\eqd(\mathcal{E}_{1}^{\vee},\mathcal{E}_{2}^{\vee},f^{\vee-1}).
\]
The Tate twist is $\underline{\mathcal{E}}\{i\}:=\underline{\mathcal{E}}\otimes\underline{\mathcal{O}}_{X}\{i\}$ where $\underline{\mathcal{O}}_{X}\{i\}:=\underline{\mathcal{O}}_{X}\{1\}^{\otimes i}$ with
\[
\underline{\mathcal{O}}_{X}\{1\}\eqd\left(\mathcal{O}_{X}\otimes_{E}E(1),\mathcal{O}_{X}(1)\otimes_{E}E(1),\mathrm{can}\otimes\mathrm{Id}\right).
\]
Here $E(1):=E\otimes_{\mathbb{Z}_{p}}\mathbb{Z}_{p}(1)$ with $\mathbb{Z}_{p}(1):=\underleftarrow{\lim}\,\mu_{p^{n}}(C)$ and $\mathrm{can}:\mathcal{O}_{X}\hookrightarrow\mathcal{O}_{X}(1)$ is the canonical morphism, dual to the embedding $\mathcal{I}(\infty)\hookrightarrow\mathcal{O}_{X}$. There are also symmetric and exterior powers, given by the following formulae: for every $k\geq0$,
\begin{align*}
\Sym^{k}(\underline{\mathcal{E}}) & \eqd (\Sym^{k}\mathcal{E}_{1},\Sym^{k}\mathcal{E}_{2},\Sym^{k}f),\\
\Lambda^{k}(\underline{\mathcal{E}}) & \eqd (\Lambda^{k}\mathcal{E}_{1},\Lambda^{k}\mathcal{E}_{2},\Lambda^{k}f).
\end{align*}
The generic fiber $\underline{\mathcal{E}}\mapsto\mathcal{E}_{1,\eta}$ yields an exact faithful $\otimes$-functor
\[
({\,\cdot\,})_{1,\eta}:\Modif_{X}\rightarrow\Vect_{E(X)}
\]
which induces an isomorphism between the poset $\Sub(\underline{\mathcal{E}})$ of strict subobjects of $\underline{\mathcal{E}}$ in $\Modif_{X}$ and the poset $\Sub(\mathcal{E}_{1,\eta})$ of $E(X)$-subspaces of $\mathcal{E}_{1,\eta}$. We say that a modification $\underline{\mathcal{E}}=(\mathcal{E}_{1},\mathcal{E}_{2},f)$ is effective if $f$ extends to a (necessarily unique) morphism $f:\mathcal{E}_{1}\rightarrow\mathcal{E}_{2}$, which is then a mono-epi in $\Bun_{X}$. For every $\underline{\mathcal{E}}$ in $\Modif_{X}$,
\[
\underline{\mathcal{E}}\{i\}\text{ is effective for }i\gg0.
\]


\subsubsection{Hodge and Newton filtrations}\label{subsec:Hodge-and-NewtonForModif}

For $\underline{\mathcal{E}}=(\mathcal{E}_{1},\mathcal{E}_{2},f)$ as above, we denote by
\[
f_{\dR}:\mathcal{E}_{1,\dR}^{+}[\xi^{-1}]\rightarrow\mathcal{E}_{2,\dR}^{+}[\xi^{-1}]
\]
the $B_{\dR}$-isomorphism induced by $f$, where $\mathcal{E}_{i,\dR}^{+}=\mathcal{E}_{i,\infty}^{\wedge}$ is the completed local stalk at $\infty$. For $i\in\{1,2\}$, the Hodge filtration $\mathcal{F}_{H,i}(\underline{\mathcal{E}})$ is the $\mathbb{Z}$-filtration induced by $\mathcal{E}_{3-i,\dR}^{+}$ on the residue $\mathcal{E}_{i}(\infty)=\mathcal{E}_{i,\dR}^{+}/\xi\mathcal{E}_{i,\dR}^{+}$ of $\mathcal{E}_{i}$. Thus for every $\gamma\in\mathbb{Z}$,
\[
\mathcal{F}_{H,1}^{\gamma}\eqd\frac{f_{\dR}^{-1}(\xi^{\gamma}\mathcal{E}_{2,\dR}^{+})\cap\mathcal{E}_{1,\dR}^{+}+\xi\mathcal{E}_{1,\dR}^{+}}{\xi\mathcal{E}_{1,\dR}^{+}},\quad\mathcal{F}_{H,2}^{\gamma}\eqd\frac{f_{\dR}(\xi^{\gamma}\mathcal{E}_{1,\dR}^{+})\cap\mathcal{E}_{2,\dR}^{+}+\xi\mathcal{E}_{2,\dR}^{+}}{\xi\mathcal{E}_{2,\dR}^{+}}.
\]
These are filtrations with opposed types $t_{H,i}(\underline{\mathcal{E}})\in\mathbb{\mathbb{Z}}_{\geq}^{r}$, where
\[
r=\rank(\underline{\mathcal{E}})=\rank(\mathcal{E}_{1})=\rank(\mathcal{E}_{2}).
\]
We denote by $\mathcal{F}_{N,i}(\underline{\mathcal{E}})$ the Newton filtration on $\mathcal{E}_{i}$ with type $t_{N,i}(\mathcal{\underline{E}})\in\mathbb{Q}_{\geq}^{r}$. Thus
\[
\begin{aligned}
t_{N,1}\left(\underline{\mathcal{E}}^{\vee}\right) & = t_{N,1}\left(\underline{\mathcal{E}}\right)^{\iota}\\
t_{N,2}\left(\underline{\mathcal{E}}^{\vee}\right) & = t_{N,2}\left(\underline{\mathcal{E}}\right)^{\iota}\\
t_{H,1}\left(\underline{\mathcal{E}}^{\vee}\right) & = t_{H,1}\left(\underline{\mathcal{E}}\right)^{\iota}\\
t_{H,2}\left(\underline{\mathcal{E}}^{\vee}\right) & = t_{H,2}\left(\underline{\mathcal{E}}\right)^{\iota}
\end{aligned}\quad\text{and}\quad
\begin{aligned}
t_{N,1}\left(\underline{\mathcal{E}}\{i\}\right) & = t_{N,1}\left(\underline{\mathcal{E}}\right)\\
t_{N,2}\left(\underline{\mathcal{E}}\{i\}\right) & = t_{N,2}\left(\underline{\mathcal{E}}\right)+(i,\dots,i)\\
t_{H,1}\left(\underline{\mathcal{E}}\{i\}\right) & = t_{H,1}\left(\underline{\mathcal{E}}\right)+(i,\dots,i)\\
t_{H,2}\left(\underline{\mathcal{E}}\{i\}\right) & = t_{H,2}\left(\underline{\mathcal{E}}\right)-(i,\dots,i)
\end{aligned}
\]
The filtrations $\mathcal{F}_{N,i}$ and $\mathcal{F}_{H,i}$ are compatible with tensor products, duals, symmetric and exterior powers. In particular for every $0\leq k\leq r$,
\[
t_{H,i}^{\max}(\Lambda^{k}\underline{\mathcal{E}})=t_{H,i}(\underline{\mathcal{E}})(k)\quad\text{and}\quad t_{N,i}^{\max}(\Lambda^{k}\underline{\mathcal{E}})=t_{N,i}(\underline{\mathcal{E}})(k)
\]
viewing the right hand side terms as functions on $[0,r]$. Also, $\underline{\mathcal{E}}$ is effective if and only if the slopes of $\mathcal{F}_{H,1}$ (resp. $\mathcal{F}_{H,2}$) are non-negative (resp.~non-positive), in which case $t_{H,1}(\underline{\mathcal{E}})$ is the type $t(\mathcal{Q})$ of the torsion $\mathcal{O}_{X}$-module $\mathcal{Q}=\mathcal{E}_{2}/f(\mathcal{E}_{1})$ supported at $\infty$, which means that if $t_{H,1}=(n_{1}\geq\dots\geq n_{r})\in\mathbb{N}_{\geq}^{r}$, then
\[
\mathcal{E}_{2}/f(\mathcal{E}_{1})\simeq\mathcal{O}_{X,\infty}/\mathfrak{m}_{\infty}^{n_{1}}\oplus\dots\oplus\mathcal{O}_{X,\infty}/\mathfrak{m}_{\infty}^{n_{r}}\simeq B_{\dR}^{+}/\xi^{n_{1}}B_{\dR}^{+}\oplus\dots\oplus B_{\dR}^{+}/\xi^{n_{r}}B_{\dR}^{+}.
\]

\begin{prop}\label{prop:CompNewtonHodgeModif}
For every modification $\underline{\mathcal{E}}$ on $X$ of rank $r\in\mathbb{N}$,
\[
t_{N,2}(\underline{\mathcal{E}})\leq t_{N,1}(\underline{\mathcal{E}})+t_{H,1}(\underline{\mathcal{E}})\quad\text{in}\quad\mathbb{Q}_{\geq}^{r}.
\]
\end{prop}

\begin{proof}
Using a Tate twist, we may assume that $\underline{\mathcal{E}}$ is effective. The left and right-hand side concave polygons then already have the same terminal point, since $\deg(\mathcal{E}_{2})=\deg(\mathcal{E}_{1})+\deg(\mathcal{Q})$ where $\mathcal{Q}=\mathcal{E}_{2}/f(\mathcal{E}_{1})$. By the formula for the exterior powers, it is then sufficient to establish that
\[
t_{N}^{\max}(\mathcal{E}_{2})\leq t_{N}^{\max}(\mathcal{E}_{1})+t^{\max}(\mathcal{Q}).
\]
Let $\underline{\mathcal{E}}^{\prime}=(\mathcal{E}_{1}^{\prime},\mathcal{E}_{2}^{\prime},f^{\prime})$ where $\mathcal{E}_{2}^{\prime}$ is the first (smallest) step of $\mathcal{F}_{N}(\mathcal{E}_{2})$, $\mathcal{E}_{1}^{\prime}=f^{-1}(\mathcal{E}_{2}^{\prime})$ and $f^{\prime}=f\vert\mathcal{E}_{1}^{\prime}$. Set $\mathcal{Q}'=\mathcal{E}_{2}^{\prime}/f'(\mathcal{E}_{1}^{\prime})$. Then $\mathcal{E}_{2}^{\prime}$ is semi-stable of slope $\mu=t_{N}^{\max}(\mathcal{E}_{2})$ and $\deg\mathcal{E}_{2}^{\prime}=\deg\mathcal{E}_{1}^{\prime}+\deg\mathcal{Q}'$, thus $t_{N}(\mathcal{E}_{2}^{\prime})\leq t_{N}(\mathcal{E}_{1}^{\prime})+t(\mathcal{Q}')$ by concavity of the sum and equality of the terminal points. Considering the first (largest) slopes, we find that $\mu\leq t_{N}^{\max}(\mathcal{E}_{1}^{\prime})+t^{\max}(\mathcal{Q}')$. But $\mathcal{E}_{1}^{\prime}\subset\mathcal{E}_{1}$ and $\mathcal{Q}'\subset\mathcal{Q}$, thus
\[
t_{N}^{\max}(\mathcal{E}_{1}^{\prime})\leq t_{N}^{\max}(\mathcal{E}_{1})\quad\text{and}\quad t^{\max}(\mathcal{Q}')\leq t^{\max}(\mathcal{Q}).
\]
This yields the desired inequality.
\end{proof}

\subsubsection{Admissible modifications}\label{subsec:AdmModifDef}

Let $\Modif_{X}^{ad}$ be the full subcategory of $\Modif_{X}$ whose objects are the modifications $\underline{\mathcal{E}}=(\mathcal{E}_{1},\mathcal{E}_{2},f)$ such that $\mathcal{E}_{1}$ is semi-stable of slope $0$, i.e.~$t_{N,1}(\underline{\mathcal{E}})=t_{N}(\mathcal{E}_{1})=0$. This is a quasi-abelian $E$-linear rigid $\otimes$-category with Tate twists. The kernels, cokernels, duals, $\otimes$-products, Tate twist, symmetric and exterior powers are induced by those of $\Modif_{X}$. On $\Modif_{X}^{ad}$, we set
\[
\mathcal{F}_{N}\eqd\mathcal{F}_{N,2},\quad\mathcal{F}_{H}\eqd\mathcal{F}_{H,1},\quad t_{N}\eqd t_{N,2}\quad\text{and}\quad t_{H}\eqd t_{H,1}.
\]

\begin{prop}\label{prop:OnAdmModif_NewtonHodge}
For every admissible modification $\underline{\mathcal{E}}$ of rank $r\in\mathbb{N}$,
\[
t_{N}(\underline{\mathcal{E}})\leq t_{H}(\underline{\mathcal{E}})\quad\text{in}\quad\mathbb{Q}_{\geq}^{r}.
\]
\end{prop}

\begin{proof}
This is the special case of Proposition~\ref{prop:CompNewtonHodgeModif} where $t_{N,1}(\underline{\mathcal{E}})=0$.
\end{proof}

The restriction of the generic fiber functor $({\,\cdot\,})_{1,\eta}:\Modif_{X}\rightarrow\Vect_{E(X)}$ to the full subcategory $\Modif_{X}^{ad}$ of $\Mod_{X}$ descends to an exact $E$-linear faithful $\otimes$-functor
\[
\omega:\Modif_{X}^{ad}\rightarrow\Vect_{E},\qquad\omega(\underline{\mathcal{E}})=\Gamma(X,\mathcal{E}_{1})
\]
inducing an isomorphism between the poset $\Sub^{ad}(\underline{\mathcal{E}})\subset\Sub(\underline{\mathcal{E}})$ of strict subobjects of $\underline{\mathcal{E}}$ in $\Modif_{X}^{ad}$ and the poset $\Sub(\omega(\underline{\mathcal{E}}))\subset\Sub(\mathcal{E}_{1,\eta})$ of $E$-subspaces of $\omega(\underline{\mathcal{E}})$.

\subsubsection{The Fargues filtration}\label{subsec:DefFarguesFiltrOnModif}

The rank and degree functions
\[
\rank:\sk\,\Modif_{X}^{ad}\rightarrow\mathbb{N}\quad\text{ and}\quad\deg:\sk\,\Modif_{X}^{ad}\rightarrow\mathbb{Z}
\]
which are respectively defined by
\[
\begin{aligned}\rank(\underline{\mathcal{E}}) & \eqd\rank(\mathcal{E}_{1})=\rank(\mathcal{E}_{2})=\dim_{E}\omega(\underline{\mathcal{E}})\\
\deg(\underline{\mathcal{E}}) & \eqd\deg\mathcal{E}_{2}=\deg\mathcal{F}_{N}(\underline{\mathcal{E}})=\deg\mathcal{F}_{H}(\underline{\mathcal{E}})
\end{aligned}
\]
$ $are additive on short exact sequences in $\Modif_{X}^{ad}$, and they are respectively constant and non-decreasing on mono-epis in $\Modif_{X}^{ad}$. More precisely if $F=(F_{1},F_{2})$ is a mono-epi $F:\underline{\mathcal{E}}\rightarrow\underline{\mathcal{E}}^{\prime},$ then $F_{1}:\mathcal{E}_{1}\rightarrow\mathcal{E}'_{1}$ is an isomorphism and $F_{2}:\mathcal{E}_{2}\rightarrow\mathcal{E}_{2}^{\prime}$ is a mono-epi in $\Bun_{X}$, thus $\deg(\underline{\mathcal{E}})=\deg(\mathcal{E}_{2})\leq\deg(\mathcal{E}_{2}^{\prime})=\deg(\underline{\mathcal{E}}^{\prime})$ with equality if and only if $F_{2}$ is an isomorphism in $\Bun_{X}$, which amounts to $F=(F_{1},F_{2})$ being an isomorphism in $\Modif_{X}^{ad}$. These rank and degree functions thus induce a Harder--Narasimhan filtration on $\Modif_{X}^{ad}$, the Fargues filtration $\mathcal{F}_{F}$ with slopes $\mu=\deg/\rank$ in $\mathbb{Q}$, and the full subcategory of $\Modif_{X}^{ad}$ of semi-stable objects of slope $\mu$ is abelian. We denote by $t_{F}(\underline{\mathcal{E}})$ the type of $\mathcal{F}_{F}(\underline{\mathcal{E}})$.

\begin{prop}\label{prop:OnAdmModif_FarguesExactness}
Let $0\rightarrow\underline{\mathcal{E}}_{1}\rightarrow\underline{\mathcal{E}}_{2}\rightarrow\underline{\mathcal{E}}_{3}\rightarrow0$ be an exact sequence in $\Modif_{X}^{ad}$, set $r_{i}=\rank\,\underline{\mathcal{E}}_{i}$ and view $t_{F}(M_{i})$ as a concave function $f_{i}:[0,r_{i}]\rightarrow\mathbb{R}$. Then
\[
f_{1}\ast f_{3}(s)\geq f_{2}(s)\geq
\begin{cases}
f_{1}(s) & \text{if }0\leq s\leq r_{1}\\
f_{1}(r_{1})+f_{3}(s-r_{1}) & \text{if }r_{1}\leq s\leq r_{2}
\end{cases}
\]
with equality for $s=0$ and $s=r_{2}$. In particular,
\begin{gather*}
\begin{aligned}
t_{F}^{\max}(\underline{\mathcal{E}}_{1}) & \leq t_{F}^{\max}(\underline{\mathcal{E}}_{2}) \leq \max\left\{ t_{F}^{\max}(\underline{\mathcal{E}}_{1}),t_{F}^{\max}(\underline{\mathcal{E}}_{3})\right\},\\
t_{F}^{\min}(\underline{\mathcal{E}}_{3}) & \geq t_{F}^{\min}(\underline{\mathcal{E}}_{2}) \, \geq \min\left\{ t_{F}^{\min}(\underline{\mathcal{E}}_{1}),t_{F}^{\min}(\underline{\mathcal{E}}_{3})\right\},
\end{aligned}
\\
\text{and}\quad t_{F}(\underline{\mathcal{E}}_{2})\leq t_{F}(\underline{\mathcal{E}}_{1})\ast t_{F}(\underline{\mathcal{E}}_{3})\quad\text{in}\quad\mathbb{Q}_{\geq}^{r_{2}}.
\end{gather*}
Moreover, $t_{F}(\underline{\mathcal{E}}_{2})=t_{F}(\underline{\mathcal{E}}_{1})\ast t_{F}(\underline{\mathcal{E}}_{3})$ if and only if for every $\gamma\in\mathbb{Q}$,
\[
0\rightarrow\mathcal{F}_{F}^{\gamma}(\underline{\mathcal{E}}_{1})\rightarrow\mathcal{F}_{F}^{\gamma}(\underline{\mathcal{E}}_{2})\rightarrow\mathcal{F}_{F}^{\gamma}(\underline{\mathcal{E}}_{3})\rightarrow0
\]
is exact.
\end{prop}

\begin{proof}
Again, see~\cite[Proposition~21]{Co16} or~\cite[4.4.4]{An09}.
\end{proof}

\begin{prop}\label{prop:OnAdmModif_CompFarguesNewton}
For every admissible modification $\underline{\mathcal{E}}$ of rank $r\in\mathbb{N}$,
\[
t_{F}(\underline{\mathcal{E}})\leq t_{N}(\underline{\mathcal{E}})\quad\text{in}\quad\mathbb{Q}_{\geq}^{r}.
\]
\end{prop}
\begin{proof}
The breaks of the concave polygon $t_{F}(\underline{\mathcal{E}})$ have coordinates
\[
\left(\rank,\deg\right)\left(\mathcal{F}_{F}^{\gamma}(\underline{\mathcal{E}})_{2}\right)\in\left\{ 0,\dots,r\right\} \times\mathbb{Z}
\]
for $\gamma\in\mathbb{Q}$, where $\mathcal{F}_{F}^{\gamma}(\underline{\mathcal{E}})_{2}$ is a strict subobject of $\mathcal{E}_{2}$ in $\Bun_{X}$, equal to $\mathcal{E}_{2}$ for $\gamma\ll0$. Thus by definition of $\mathcal{F}_{N}(\mathcal{E}_{2})$, we find that $t_{F}(\underline{\mathcal{E}})$ lies below $t_{N}(\mathcal{E}_{2})=t_{N}(\underline{\mathcal{E}})$ and both polygons have the same terminal points, which proves the proposition.
\end{proof}

\subsubsection{}

Let $\underline{\mathcal{E}}=(\mathcal{E}_{1},\mathcal{E}_{2},\alpha)$ be an admissible modification and set $V=\Gamma(X,\mathcal{E}_{1})$, so that $\mathcal{E}_{1,\eta}=V_{E(X)}$ and $\mathcal{E}_{1}(\infty)=V_{C}$. We view $\mathcal{F}_{H}=\mathcal{F}_{H}(\underline{\mathcal{E}})$ as an element of $\mathbf{F}(V_{C})$, $\mathcal{F}_{N}^{\ast}:=\alpha_{\eta}^{-1}(\mathcal{F}_{N}(\mathcal{E}_{2})_{\eta})$ as an element of $\mathbf{F}(V_{E(X)})$ and $\mathcal{F}_{F}^{\ast}:=\Gamma(X,\mathcal{F}_{F}(\underline{\mathcal{E}})_{1})$ as an element of $\mathbf{F}(V)$. For every $\mathcal{F}\in\mathbf{F}(V_{E(X)})$, define
\[
\left\langle \mathcal{E}_{1},\mathcal{F}\right\rangle \eqd\sum_{\gamma\in\mathbb{R}}\gamma\deg\Gr_{\mathcal{F}}^{\gamma}\left(\mathcal{E}_{1}\right)\quad\text{and}\quad\left\langle \mathcal{E}_{2},\mathcal{F}\right\rangle \eqd\sum_{\gamma\in\mathbb{R}}\gamma\deg\Gr_{\mathcal{F}}^{\gamma}\left(\mathcal{E}_{2}\right).
\]
Here $\Gr_{\mathcal{F}}^{\gamma}(\mathcal{E}_{i}):=\mathcal{F}^{\geq\gamma}(\mathcal{E}_{i})/\mathcal{F}^{>\gamma}(\mathcal{E}_{i})$ where $\mathcal{F}^{\geq\gamma}(\mathcal{E}_{i})$ and $\mathcal{F}^{>\gamma}(\mathcal{E}_{i})$ are the strict subobjects of $\mathcal{E}_{i}$ with generic fiber $\mathcal{F}^{\geq\gamma}$ and $\mathcal{F}^{>\gamma}$ in $V_{E(X)}=\mathcal{E}_{1,\eta}$ if $i=1$, or $\alpha_{\eta}(\mathcal{F}^{\geq\gamma})$ and $\alpha_{\eta}(\mathcal{F}^{>\gamma})$ in $\mathcal{E}_{2,\eta}$ if $i=2$. Thus whenever $\{\gamma_{s}>\dots>\gamma_{0}\}\subset\mathbb{R}$ contains
\[
\Jump(\mathcal{F})\eqd\left\{ \gamma\in\mathbb{R}:\Gr_{\mathcal{F}}^{\gamma}\neq0\right\},
\]
we have for any $i\in\{1,2\}$ the following equality:
\[
\left\langle \mathcal{E}_{i},\mathcal{F}\right\rangle =\gamma_{0}\deg(\mathcal{E}_{i})+\sum_{j=1}^{s}\left(\gamma_{j}-\gamma_{j-1}\right)\deg\mathcal{F}^{\gamma_{j}}\left(\mathcal{E}_{i}\right).
\]
Since $\mathcal{E}_{1}$ is semi-stable of slope $0$, $\left\langle
\mathcal{E}_{1},\mathcal{F}\right\rangle \leq0$ with equality if and only if each $\mathcal{F}^{\gamma_{j}}(\mathcal{E}_{1})$ is of degree~$0$. We thus obtain: for every $\mathcal{F}\in\mathbf{F}(V_{E(X)})$,
\begin{align*}
\left\langle \mathcal{E}_{1},\mathcal{F}\right\rangle \geq0 & \iff \left\langle \mathcal{E}_{1},\mathcal{F}\right\rangle =0,\\
& \iff {\forall\gamma\in\mathbb{R}},\,\mathcal{F}^{\gamma}(\mathcal{E}_{1})\text{ is semi-stable of slope }0,\\
& \iff \mathcal{F}\in\mathbf{F}(V).
\end{align*}

\begin{prop}\label{prop:EquivCond4Fargues=00003DNewton}
With notations as above, the following conditions are equivalent:
\begin{align*}
\mathcal{F}_{F}^{\ast}=\mathcal{F}_{N}^{\ast} & \iff \mathcal{F}_{N}^{\ast}\in\mathbf{F}(V),\\
& \iff \left\langle \mathcal{E}_{1},\mathcal{F}_{N}^{\ast}\right\rangle \geq0,\\
& \iff \left\langle \mathcal{E}_{1},\mathcal{F}_{N}^{\ast}\right\rangle =0,\\
& \iff {\forall\gamma\in\mathbb{R}},\,(\mathcal{F}_{N}^{\ast})^{\gamma}(\mathcal{E}_{1})\text{ is semi-stable of slope }0.
\end{align*}
If $\underline{\mathcal{E}}=(\mathcal{E}_{1},\mathcal{E}_{2},\alpha)$ is effective, $(\mathcal{F}_{N}^{\ast})^{\gamma}(\mathcal{E}_{1})=\alpha^{-1}(\mathcal{F}_{N}^{\gamma}(\mathcal{E}_{2}))$ thus also
\[
\mathcal{F}_{F}^{\ast}=\mathcal{F}_{N}^{\ast}\iff{\forall\gamma\in\mathbb{R}},\,\alpha^{-1}(\mathcal{F}_{N}^{\gamma}(\mathcal{E}_{2}))\text{ is semi-stable of slope }0.
\]
\end{prop}

\begin{proof}
By~\cite[Proposition~6]{Co16}:
\begin{enumerate}
\item $\mathcal{F}_{N}^{\ast}$ is the unique element $\mathcal{F}$ of $\mathbf{F}(V_{E(X)})$ such that $\left\langle\mathcal{E}_{2},\mathcal{G}\right\rangle \leq\left\langle\mathcal{F},\mathcal{G}\right\rangle $ for every $\mathcal{G}\in\mathbf{F}(V_{E(X)})$ with equality for $\mathcal{G}=\mathcal{F}$, and;
\item $\mathcal{F}_{F}^{\ast}$ is the unique element $f$ of $\mathbf{F}(V)$ such that $\left\langle\mathcal{E}_{2},g\right\rangle \leq\left\langle f,g\right\rangle $ for every $g\in\mathbf{F}(V)$ with equality for $g=f$.
\end{enumerate}
Thus $\mathcal{F}_{F}^{\ast}=\mathcal{F}_{N}^{\ast}\Leftrightarrow\mathcal{F}_{N}^{\ast}\in\mathbf{F}(V)$ and the proposition follows.
\end{proof}

\subsection{Hodge--Tate modules}\label{subsec:HodgeTateModDef}

\subsubsection{}

Let $\HT_{E}^{B_{\dR}}$ be the category of pairs $(V,\Xi)$ where $V$ is a finite $E$-vector space and $\Xi$ is a $B_{\dR}^{+}$-lattice in $V_{\dR}=V\otimes_{E}B_{\dR}$. A morphism $F:(V,\Xi)\rightarrow(V',\Xi')$ is an $E$-linear morphism $f:V\rightarrow V'$ whose $B_{\dR}$-linear extension $f_{\dR}:V_{\dR}\rightarrow V_{\dR}^{\prime}$ satisfies $f_{\dR}(\Xi)\subset\Xi'$. The kernel and cokernel of $F$ are given by
\[
\ker(F)=\left(\ker(f),\ker(f_{\dR})\cap\Xi\right)\quad\text{and}\;\;\coker(F)=\left(V'/\mathrm{im}(f),\Xi'/\mathrm{im}(f_{\dR})\cap\Xi'\right).
\]
This defines a quasi-abelian rigid $E$-linear $\otimes$-category with tensor product
\[
(V_{1},\Xi_{1})\otimes(V_{2},\Xi_{2})\eqd(V_{1}\otimes_{E}V_{2},\Xi_{1}\otimes_{B_{\dR}^{+}}\Xi_{2}),
\]
neutral object $(E,B_{\dR}^{+})$ and duals, symmetric and exterior powers given by
\[
(V,\Xi)^{\vee}\eqd(V^{\vee},\Xi^{\vee}),\quad\Sym^{k}(V,\Xi)\eqd(\Sym^{k}V,\Sym^{k}\Xi),\quad\Lambda^{k}(V,\Xi)\eqd(\Lambda^{k}V,\Lambda^{k}\Xi)
\]
where the tensor product constructions are over $E$ or $B_{\dR}^{+}$.

\subsubsection{}

There is an (exact) $\otimes$-equivalence of $\otimes$-categories~\cite[4.2.3]{Fa15}
\[
\HTT:\Modif_{X}^{ad}\rightarrow\HT_{E}^{B_{\dR}},\qquad\HTT(\mathcal{E}_{1},\mathcal{E}_{2},f)\eqd\left(\Gamma(X,\mathcal{E}_{1}),f_{\dR}^{-1}(\mathcal{E}_{2,\dR}^{+})\right).
\]
with inverse $(V,\Xi)\mapsto(\mathcal{E}_{1},\mathcal{E}_{2},f)$ where $(\mathcal{E}_{2},f)$ is the modification of $\mathcal{E}_{1}=V\otimes_{E}\mathcal{O}_{X}$ at $\infty\in X$ corresponding to the $(\mathcal{O}_{X,\infty}^{\wedge}=B_{\dR}^{+})$-lattice $\Xi$ of $\mathcal{E}_{1,\infty}^{\wedge}[\xi^{-1}]=V\otimes_{E}B_{\dR}$ under the Beauville--Laszlo correspondence of~\cite[5.3.1]{FaFo15} (see also Section~$4.1$ of Colmez's preface to~\cite{FaFo15}). Under this equivalence of categories, the Hodge filtration $\mathcal{F}_{H}(V,\Xi)$ is the $\mathbb{Z}$-filtration which is induced by $\Xi$ on the residue $V_{C}=V\otimes_{E}C$ of the standard lattice $V_{\dR}^{+}=V\otimes_{E}B_{\dR}^{+}$ of $V_{\dR}$: for every $\gamma\in\mathbb{Z}$,
\[
\mathcal{F}_{H}^{\gamma}(V,\Xi)\eqd\frac{V_{\dR}^{+}\cap\xi^{\gamma}\Xi+\xi V_{\dR}^{+}}{\xi V_{\dR}^{+}}\quad\text{ in}\quad V_{C}=\frac{V_{\dR}^{+}}{\xi V_{\dR}^{+}}.
\]
We denote by $t_{H}(V,\Xi)$ the type of $\mathcal{F}_{H}(V,\Xi)$. The rank and degree functions
\[
\rank:\sk\,\HT_{E}^{B_{\dR}}\rightarrow\mathbb{N}\quad\text{ and}\quad\deg:\sk\,\HT_{E}^{B_{\dR}}\rightarrow\mathbb{Z}
\]
are respectively given by
\[
\begin{aligned}\rank(V,\Xi) & \eqd\dim_{E}(V)=\rank_{B_{\dR}^{+}}(\Xi),\\
\deg(V,\Xi) & \eqd\nu(V_{\dR}^{+},\Xi)=\deg\mathcal{F}_{H}(V,\Xi).
\end{aligned}
\]
We denote by $\mathcal{F}_{F}(V,\Xi)$ the corresponding Fargues $\mathbb{Q}$-filtration, with type $t_{F}(V,\Xi)$ in $\mathbb{Q}_{\geq}^{r}$ if $r=\dim_{E}V$. The Tate object is $\HTT(\underline{\mathcal{O}}_{X}\{1\})=\left(E(1),\xi^{-1}E(1)_{\dR}^{+}\right)$.

\begin{prop}
Let $f:(V_{1},\Xi_{1})\rightarrow(V_{2},\Xi_{2})$ be a mono-epi in $\HT_{E}^{B_{\dR}}$, so that $f:V_{1}\rightarrow V_{2}$ is an isomorphism and $f_{\dR}:\Xi_{1}\rightarrow\Xi_{2}$ is injective with cokernel $Q$ of finite length. If $r=\dim_{E}V_{1}=\dim_{E}V_{2}$, then for every $s\in[0,r]$,
\[
0\leq t_{F}(V_{2},\Xi_{2})(s)-t_{F}(V_{1},\Xi_{1})(s)\leq\length_{B_{\dR}^{+}}(Q).
\]
with equality on the left (resp.~right) for $s=0$ (resp.~$s=r$). In particular,
\[
0\leq\left\{\!
\begin{array}{c}
t_{F}^{\max}(V_{2},\Xi_{2})-t_{F}^{\max}(V_{1},\Xi_{1})\\
t_{F}^{\min}(V_{2},\Xi_{2})-t_{F}^{\min}(V_{1},\Xi_{1})
\end{array}\!\right\} \leq\length_{B_{\dR}^{+}}(Q).
\]
\end{prop}

\begin{proof}
This is analogous to Proposition~\ref{prop:PseudoIsogInModt}.
\end{proof}

\subsubsection{}

There is also an exact and fully faithful $\otimes$-functor from the category $\HT_{E}^{B_{\dR}}$ to the quasi-abelian $\otimes$-category denoted by $\Norm_{E}^{B_{\dR}}$ in~\cite[\S 5.2]{Co16}, which maps $(V,\Xi)$ to $(V,\alpha_{\Xi})$ where $\alpha_{\Xi}:V_{\dR}\rightarrow\mathbb{R}_{+}$ is the gauge norm of the $B_{\dR}^{+}$-lattice $\Xi\subset V_{\dR}$. This functor is plainly compatible with the rank and degree functions of both categories (for the appropriate normalization of the valuation on $B_{\dR}$), and its essential image is stable under strict subobjects. It is therefore also compatible with the corresponding Harder--Narasimhan filtrations. Since the Harder--Narasimhan filtration on $\Norm_{E}^{B_{\dR}}$ is compatible with tensor products, duals, symmetric and exterior powers by~\cite[Proposition~28]{Co16}, we obtain the following proposition:

\begin{prop}\label{prop:OnHT+AdModif_Farguestensor}
The Fargues filtrations $\mathcal{F}_{F}$ on $\HT_{E}^{B_{\dR}}$ and $\Modif_{X}^{ad}$ are compatible with tensor products, duals, symmetric and exterior powers.
\end{prop}

\subsubsection{}

Fix an admissible modification $\underline{\mathcal{E}}$ of rank $r$ and set $\HTT(\underline{\mathcal{E}})=\left(V,\Xi\right)$. Then
\[
\mathcal{F}_{H}=\mathcal{F}_{H}(\underline{\mathcal{E}})=\mathcal{F}_{H}(V,\Xi)
\]
is the $\mathbb{Z}$-filtration on $V_{C}=\mathcal{E}_{1}(\infty)$ which is denoted by $\loc(\alpha_{\Xi})$ in~\cite[6.4]{Co16}, where $\alpha_{\Xi}$ is the gauge norm of the $B_{\dR}^{+}$-lattice $\Xi\subset V_{\dR}$. For any $\mathcal{F}\in\mathbf{F}(V_{E(X)})$, we set
\[
\left\langle \underline{\mathcal{E}},\mathcal{F}\right\rangle \eqd\left\langle \mathcal{E}_{2},\mathcal{F}\right\rangle -\left\langle \mathcal{E}_{1},\mathcal{F}\right\rangle.
\]
Thus if $\Jump(\mathcal{F})\subset\{\gamma_{s}>\dots>\gamma_{0}\}\subset\mathbb{R}$ for some $s\in\mathbb{N}$, then
\[
\left\langle \underline{\mathcal{E}},\mathcal{F}\right\rangle =\gamma_{0}\deg(\mathcal{E}_{2})+\sum_{i=1}^{s}(\gamma_{i}-\gamma_{i-1})\left(\deg\mathcal{F}^{\gamma_{i}}(\mathcal{E}_{2})-\deg\mathcal{F}^{\gamma_{i}}(\mathcal{E}_{1})\right).
\]
Suppose first that $\underline{\mathcal{E}}=(\mathcal{E}_{1},\mathcal{E}_{2},f)$ is effective. Then $\mathcal{F}^{\gamma}(\mathcal{E}_{1})=f^{-1}(\mathcal{F}^{\gamma}(\mathcal{E}_{2}))$ and
\begin{align*}
\left\langle \underline{\mathcal{E}},\mathcal{F}\right\rangle & = \gamma_{0}\deg(\mathcal{Q})+\sum_{i=1}^{s}(\gamma_{i}-\gamma_{i-1})\left(\deg\mathcal{F}^{\gamma_{i}}(\mathcal{Q})\right)\\
& = \sum_{\gamma\in\mathbb{R}}\gamma\deg\Gr_{\mathcal{F}}^{\gamma}(\mathcal{Q})
\end{align*}
where $\mathcal{F}^{\gamma}(\mathcal{Q})$ is the image of $\mathcal{F}^{\gamma}(\mathcal{E}_{2})$ in the torsion sheaf $\mathcal{Q}=\mathcal{E}_{2}/f(\mathcal{E}_{1})$ on $X$ and $\Gr_{\mathcal{F}}^{\gamma}(\mathcal{Q})=\mathcal{F}^{\geq\gamma}(\mathcal{Q})/\mathcal{F}^{>\gamma}(\mathcal{Q})$. These are skyscraper sheaves supported at $\infty$, with
\[
\Gamma(X,\mathcal{Q})=\Xi/V_{\dR}^{+}\quad\text{ and}\quad\Gamma(X,\mathcal{F}^{\gamma}(Q))=\Xi\cap\mathcal{F}_{\dR}^{\gamma}/V_{\dR}^{+}\cap\mathcal{F}_{\dR}^{\gamma}
\]
where $\mathcal{F}_{\dR}\in\mathbf{F}(V_{\dR})$ is the base change of $\mathcal{F}$ through $E(X)\hookrightarrow B_{\dR}$. Therefore
\[
\left\langle \underline{\mathcal{E}},\mathcal{F}\right\rangle =\sum\gamma\nu\left(\Gr_{\mathcal{F}_{\dR}}^{\gamma}V_{\dR}^{+},\Gr_{\mathcal{F}_{\dR}}^{\gamma}\Xi\right)=\left\langle \overrightarrow{\circ\alpha_{\Xi}},\mathcal{F}_{\dR}\right\rangle
\]
where $\circ$ is the gauge norm of $V_{\dR}^{+}\subset V_{\dR}$ and the right-hand side term is the Busemann scalar product, see~\cite[6.4.15]{Co14}. This formula still holds true for a non-necessarily effective admissible modification $\underline{\mathcal{E}}$, since indeed for every $i\in\mathbb{Z}$,
\[
\left\langle \underline{\mathcal{E}}\{i\},\mathcal{F}\right\rangle =\left\langle \underline{\mathcal{E}},\mathcal{F}\right\rangle +i\deg\mathcal{F}\quad\text{ and}\quad\left\langle \overrightarrow{\circ\alpha_{\xi^{-i}\Xi}},\mathcal{F}_{\dR}\right\rangle =\left\langle \overrightarrow{\circ\alpha_{\Xi}},\mathcal{F}_{\dR}\right\rangle +i\deg(\mathcal{F}_{\dR}).
\]
Returning thus to the general case, we now obtain:
\[
\left\langle \underline{\mathcal{E}},\mathcal{F}\right\rangle =\left\langle \overrightarrow{\circ\alpha_{\Xi}},\mathcal{F}_{\dR}\right\rangle \leq\left\langle \loc(\alpha_{\Xi}),\loc(\mathcal{F}_{\dR})\right\rangle =\left\langle \mathcal{F}_{H},\mathcal{F}_{C}\right\rangle \leq\left\langle t_{H},\mathbf{t}(\mathcal{F})\right\rangle.
\]
Here $t_{H}=\mathbf{t}(\mathcal{F}_{H})$ is the Hodge type of $\underline{\mathcal{E}}$ and $\mathcal{F}_{C}=\loc(\mathcal{F}_{\dR})$ is the $\mathbb{R}$-filtration on $V_{C}=V_{\dR}^{+}/\xi V_{\dR}^{+}$ which is induced by the $\mathbb{R}$-filtration $\mathcal{F}_{\dR}$ on $V_{\dR}$, so that
\[
\mathbf{t}(\mathcal{F}_{C})=\mathbf{t}(\mathcal{F}_{\dR})=\mathbf{t}(\mathcal{F})
\]
in $\mathbb{R}_{\geq}^{r}$. The last pairing is the standard scalar product on $\mathbb{R}_{\geq}^{r}\subset\mathbb{R}^{r}$, and the two inequalities come from~\cite[4.2 \& 5.5]{Co14}. For $\mathcal{F}=\mathcal{F}_{N}^{\ast}=\mathcal{F}_{N}^{\ast}(\underline{\mathcal{E}})$, we obtain
\[
\left\langle \underline{\mathcal{E}},\mathcal{F}_{N}^{\ast}\right\rangle \leq\left\langle t_{H},t_{N}\right\rangle
\]
where $t_{N}=\mathbf{t}(\mathcal{F}_{N})$ is the Newton type of $\underline{\mathcal{E}}$. Now we have already seen that
\[
\left\langle \underline{\mathcal{E}},\mathcal{F}_{N}^{\ast}\right\rangle =\left\langle \mathcal{E}_{2},\mathcal{F}_{N}^{\ast}\right\rangle -\left\langle \mathcal{E}_{1},\mathcal{F}_{N}^{\ast}\right\rangle =\left\Vert t_{N}\right\Vert ^{2}-\left\langle \mathcal{E}_{1},\mathcal{F}_{N}^{\ast}\right\rangle
\]
with $\left\langle \mathcal{E}_{1},\mathcal{F}_{N}^{\ast}\right\rangle
\leq0$, and we thus obtain the following inequalities:
\[
\left\Vert t_{N}\right\Vert ^{2}-\left\langle t_{H},t_{N}\right\rangle \leq\left\langle \mathcal{E}_{1},\mathcal{F}_{N}^{\ast}\right\rangle \leq0.
\]

\begin{prop}\label{prop:Crit4Fargues=00003DNewton}
With notations as above, $\left\Vert t_{N}\right\Vert ^{2}\leq\left\langle t_{H},t_{N}\right\rangle $ and
\[
\left\Vert t_{N}\right\Vert ^{2}=\left\langle t_{H},t_{N}\right\rangle \Longrightarrow\left\langle \mathcal{E}_{1},\mathcal{F}_{N}^{\ast}\right\rangle =0\iff\mathcal{F}_{N}^{\ast}=\mathcal{F}_{F}.
\]
\end{prop}
\begin{proof}
This now follows from~Proposition~\ref{prop:EquivCond4Fargues=00003DNewton}.
\end{proof}

\subsubsection{}

Let $\HT_{\mathcal{O}_{E}}^{B_{\dR}}$ be the category of pairs $(T,\Xi)$ where $T$ is a finite free $\mathcal{O}_{E}$-module and $\Xi$ is a $B_{\dR}^{+}$-lattice in $V_{\dR}=T\otimes_{\mathcal{O}_{E}}B_{\dR}=V\otimes_{E}B_{\dR}$, where $V=T\otimes_{\mathcal{O}_{E}}E$. A morphism $F:(T,\Xi)\rightarrow(T',\Xi')$ in $\HT_{\mathcal{O}_{E}}^{B_{\dR}}$ is an $\mathcal{O}_{E}$-linear morphism $f:T\rightarrow T'$ whose $B_{\dR}$-linear extension $f_{\dR}:V_{\dR}\rightarrow V_{\dR}^{\prime}$ satisfies $f_{\dR}(\Xi)\subset\Xi'$. Any such morphism has a kernel and a cokernel, which are respectively given by
\[
\left(\ker(f),\ker(f_{\dR})\cap\Xi\right)\quad\text{ and}\quad\left(T'/f(T)^{sat},\Xi'/f_{\dR}(V_{\dR})\cap\Xi'\right)
\]
where $f(T)^{sat}/f(T)$ is the torsion submodule of $T'/f(T)$. This defines a quasi-abelian rigid $\mathcal{O}_{E}$-linear $\otimes$-category with tensor product
\[
(T_{1},\Xi_{1})\otimes(T_{2},\Xi_{2})\eqd(T_{1}\otimes_{\mathcal{O}_{E}}T_{2},\Xi_{1}\otimes_{B_{\dR}^{+}}\Xi_{2}),
\]
neutral object $(\mathcal{O}_{E},B_{\dR}^{+})$ and duals, symmetric and exterior powers given by
\[
(T,\Xi)^{\vee}\eqd(T^{\vee},\Xi^{\vee}),\quad\Sym^{k}(T,\Xi)\eqd(\Sym^{k}T,\Sym^{k}\Xi),\quad\Lambda^{k}(T,\Xi)\eqd(\Lambda^{k}T,\Lambda^{k}\Xi)
\]
where the tensor product constructions are over $\mathcal{O}_{E}$ or $B_{\dR}^{+}$. There is also a Tate twist in $\HT_{\mathcal{O}_{E}}^{B_{\dR}}$, corresponding to the Tate object $(\mathcal{O}_{E}(1),\xi^{-1}E(1)_{\dR}^{+})$.

\subsubsection{}

The exact and faithful $\otimes$-functor
\[
\HT_{\mathcal{O}_{E}}^{B_{\dR}}\rightarrow\HT_{E}^{B_{\dR}}\qquad(T,\Xi)\mapsto(V,\Xi)\quad\text{with}\quad V=T\otimes_{\mathcal{O}_{E}}E
\]
induces a $\otimes$-equivalence of $\otimes$-categories
\[
\HT_{\mathcal{O}_{E}}^{B_{\dR}}\otimes E\rightarrow\HT_{E}^{B_{\dR}}.
\]


\subsection{The Bhatt--Morrow--Scholze functor}\label{subsec:The-Bhatt--Morrow--Scholze-functor}

\subsubsection{}

Let $(M,\varphi_{M})$ be a finite free BKF-module over $A$. Then $M\otimes_{A}A(C)$ is a finite free \'etale $\varphi$-module over $A(C)=W_{\mathcal{O}_{E}}(C^{\flat})$, thus by~\ref{subsec:RisW(Kflat)},
\[
T\eqd\left\{ x\in M\otimes_{A}A(C):\phi_{M\otimes A(C)}(x)=x\right\}
\]
is finite free over $\mathcal{O}_{E}$ and $T\hookrightarrow M\otimes A(C)$ extends to a $\varphi$-equivariant isomorphism
\[
T\otimes_{\mathcal{O}_{E}}A(C)\stackrel{\simeq}{\longrightarrow}M\otimes_{A}A(C).
\]
By~\cite[4.26]{BaMoSc16}, the latter descends to the subring $A[\frac{1}{\mu}]\subset A(C)$, giving an isomorphism
\[
\eta_{M}:T\otimes_{\mathcal{O}_{E}}A\!\left[\frac{1}{\mu}\right]\stackrel{\simeq}{\longrightarrow}M\!\left[\frac{1}{\mu}\right].
\]
Note that since $\mu=[\epsilon]-1$ has residue $\epsilon-1\neq0$ in $C^{\flat}$, it is indeed invertible in $A(C)=W_{\mathcal{O}_{E}}(C^{\flat})$. Tensoring with $A[\frac{1}{\mu}]\hookrightarrow B_{\dR}$, we obtain an isomorphism
\[
\eta_{M,\dR}:T\otimes_{\mathcal{O}_{E}}B_{\dR}\stackrel{\simeq}{\longrightarrow}M\otimes_{A}B_{\dR}.
\]
This yields a Hodge--Tate module $(T,\Xi)$ over $\mathcal{O}_{E}$, with
\[
\Xi\eqd\eta_{M,\dR}^{-1}(M\otimes_{A}B_{\dR}^{+}).
\]
We have thus defined an $\mathcal{O}_{E}$-linear $\otimes$-functor
\[
\HTT^{\prime}:\Mod_{A,f}^{\varphi}\rightarrow\HT_{\mathcal{O}_{E}}^{B_{\dR}},\qquad M\mapsto(T,\Xi).
\]
With $V=T\otimes_{\mathcal{O}_{E}}E$ as usual, we also denote by
\[
\HTT^{\prime}:\Mod_{A,f}^{\varphi}\otimes E\rightarrow\HT_{E}^{B_{\dR}},\qquad M\otimes E\mapsto(V,\Xi)
\]
the induced $E$-linear $\otimes$-functor.

\subsubsection{Compatibility with Hodge filtrations}

Since $\xi'=\frac{\varphi(\mu)}{\mu}$ is already invertible in $A[\frac{1}{\varphi(\mu)}]$, there is a commutative diagram whose first square is made of isomorphisms,
\[
\xymatrix{T\otimes_{\mathcal{O}_{E}}A[{\frac{1}{\mu}}]\ar[dd]_{\mathrm{Id}\otimes\varphi}\ar[r]^{\eta_{M}} & M[{\frac{1}{\mu}}]\ar[d]^{\varphi}\ar@{^{(}->}[r] & M\otimes A(C)\ar[dd]^{\phi_{M\otimes A(C)}}\\
& \varphi^{\ast}M[{\frac{1}{\varphi(\mu)}}]\ar[d]^{\varphi_{M}}\\
T\otimes_{\mathcal{O}_{E}}A[{\frac{1}{\varphi(\mu)}]}\ar[r]^{\eta_{M}} & M[{\frac{1}{\varphi(\mu)}}]\ar@{^{(}->}[r] & M\otimes A(C) }
\]
This first square induces yet another commutative diagram of isomorphisms
\[
\xymatrix{T\otimes_{\mathcal{O}_{E}}B_{\dR}\ar[dd]_{\mathrm{Id}\otimes\varphi}\ar[r]^{\eta_{M,\dR}} & M\otimes_{A}B_{\dR}\ar[d]^{\varphi\otimes\varphi}\\
& \varphi^{\ast}M\otimes_{A}B_{\dR}^{\prime}\ar[d]^{\varphi_{M}}\\
T\otimes_{\mathcal{O}_{E}}B_{\dR}^{\prime}\ar[r]^{\eta_{M,\dR}^{\prime}} & M\otimes_{A}B_{\dR}^{\prime} }
\]
with notations as in~\ref{subsec:RisA=00005B1/pi=00005D}. Restricting to lattices, we obtain the following commutative diagrams of isomorphisms (for the second diagram, note that $\mu\in(B_{\dR}^{\prime+})^{\times}$):
\[
\xymatrix@C=1pc{\Xi\ar[dd]_{\mathrm{Id}\otimes\varphi}\ar[r] & M\otimes_{A}B_{\dR}^{+}\ar[d]^{\varphi\otimes\varphi} & & V\otimes_{E}B_{\dR}^{+}\ar[dd]_{\mathrm{Id}\otimes\varphi}\ar[r] & \eta_{M,\dR}(V\otimes_{E}B_{\dR}^{+})\ar[d]^{\varphi\otimes\varphi}\\
& \varphi^{\ast}M\otimes_{A}B_{\dR}^{\prime+}\ar[d]^{\varphi_{M}} & \text{and} & & \varphi_{M}^{-1}(M\otimes_{A}B_{\dR}^{\prime+})\ar[d]^{\varphi_{M}}\\
(\mathrm{Id}\otimes\varphi)(\Xi)\ar[r] & \varphi_{M}(\varphi^{\ast}M\otimes_{A}B_{\dR}^{\prime+}) & & V\otimes_{E}B_{\dR}^{\prime+}\ar[r] & M\otimes_{A}B_{\dR}^{\prime+} }
\]
It follows that our various Hodge $\mathbb{Z}$-filtrations
\[
\begin{aligned}
\mathcal{F}_{H}(M\otimes E)\;\;&\text{on}\;\; M\otimes_{A}C^{\prime},\\
\mathcal{F}_{H}^{\prime}(M\otimes E)\;\;&\text{on}\;\; M\otimes_{A}C,
\end{aligned}\quad\text{ and}\quad
\begin{aligned}
\mathcal{F}_{H}(V,\Xi)=\mathcal{F}(V\otimes B_{\dR}^{+},\Xi)\;\;&\text{on}\;\; V\otimes_{E}C,\\
\mathcal{F}_{H}^{\prime}(V,\Xi)=\mathcal{F}(\Xi,V\otimes B_{\dR}^{+})\;\;&\text{on}\;\;\Xi\otimes_{B_{\dR}^{+}}C.
\end{aligned}
\]
are related as follows:
\[
\begin{aligned}
\mathcal{F}_{H}(M\otimes E) & = \eta_{M,C}^{\prime}\left(\mathcal{F}_{H}(V,\Xi)\otimes_{C,\varphi}C'\right) && \text{on}\;\; M\otimes_{A}C^{\prime},\\
\mathcal{F}_{H}^{\prime}(M\otimes E) & = \eta_{M,C}\left(\mathcal{F}_{H}^{\prime}(V,\Xi)\right) && \text{on}\;\; M\otimes_{A}C
\end{aligned}
\]
where $\varphi:C\rightarrow C'$ is the residue of $\varphi:B_{\dR}^{+}\rightarrow B_{\dR}^{\prime+}$ and the isomorphisms
\[
\eta_{M,C}:\Xi\otimes_{B_{\dR}^{+}}C\stackrel{\simeq}{\longrightarrow}M\otimes_{A}C\quad\text{and}\quad\eta_{M,C}^{\prime}:T\otimes_{\mathcal{O}_{E}}C'\stackrel{\simeq}{\longrightarrow}M\otimes_{A}C'
\]
are respectively induced by
\[
\eta_{M,\dR}:\Xi\stackrel{\simeq}{\longrightarrow}M\otimes_{A}B_{\dR}^{+}\quad\text{and}\quad\eta_{M,\dR}^{\prime}:T\otimes_{\mathcal{O}_{E}}B_{\dR}^{\prime+}\stackrel{\simeq}{\longrightarrow}M\otimes_{A}B_{\dR}^{\prime+}.
\]


\subsubsection{Compatibility with Tate objects}

The Tate object of $\Mod_{A,f}^{\varphi}$ is given by
\[
A\{1\}=\left({\frac{1}{\mu}}A\otimes\mathcal{O}_{E}(1),\varphi\otimes\mathrm{Id}\right).
\]
Thus since $\mu$ is invertible in $A(C)$,
\[
A\{1\}(C)=\left(A(C)\otimes\mathcal{O}_{E}(1),\varphi\otimes\mathrm{Id}\right).
\]
Since $\mathcal{O}_{E}=A(C)^{\varphi=\mathrm{Id}}$ and $t_{H}(A\{1\})=1$, it follows that
\[
\HTT^{\prime}\left(A\{1\}\right)=(\mathcal{O}_{E}(1),\xi^{-1}E(1)_{\dR}^{+})
\]
is the Tate object of $\HT_{\mathcal{O}_{E}}^{B_{\dR}}$.

\subsubsection{Fargues's theorem}\label{subsec:Fargue's-theorem}

The following theorem was conjectured by Fargues in~\cite{Fa15}.

\begin{theo}[Fargues, Scholze]\label{thm:FarguesScholze}
The $\otimes$-functors
\[
\HTT^{\prime}:\Mod_{A,f}^{\varphi}\rightarrow\HT_{\mathcal{O}_{E}}^{B_{\dR}}\quad\text{and}\quad\HTT^{\prime}:\Mod_{A,f}^{\varphi}\otimes E\rightarrow\HT_{E}^{B_{\dR}}
\]
are equivalences of $\otimes$-categories.
\end{theo}

The full faithfulness is established in~\cite[4.29]{BaMoSc16}. A proof of the essential surjectivity is sketched in Scholze's Berkeley lectures~\cite{ScWe15}, where it is mostly attributed to Fargues. An expanded and referenced version of this sketch is given in Section~\ref{subsec:AnalConstruction} below.

\begin{coro}
The categories $\Mod_{A,f}^{\varphi}$ and $\Mod_{A,f}^{\varphi}\otimes E$ are quasi-abelian.
\end{coro}

In particular, any morphism in these categories has a kernel and a cokernel. But we have no explicit and manageable formulas for them. Note also that we have two structures of exact category on $\Mod_{A,f}^{\varphi}$ and $\Mod_{A,f}^{\varphi}\otimes E$: the canonical structure which any quasi-abelian category has, and the naive structure inherited from the abelian category $\Mod_{A}^{\varphi}$. A three term complex which is naively exact is also canonically exact, but the converse is not true. We will investigate this in Section~\ref{subsec:Exactness}.
\begin{rema}
We have defined in~\ref{subsec:DefFullSubCatHNType} a strictly full subcategory $\Mod_{A,f}^{\varphi,\ast}$ of $\Mod_{A,f}^{\varphi}$, as those finite free BKF-modules $M$ such that the Fargues filtrations on the $M_{n}$'s induces a well-behaved Fargues filtration on $M$. We do not know how to describe its essential image in $\HT_{\mathcal{O}_{E}}^{B_{\dR}}$. We will later define a strictly full subcategory $\Modif_{X}^{ad,\ast}$ of $\Modif_{X}^{ad}$, corresponding to the admissible modifications $(\mathcal{E}_{1},\mathcal{E}_{2},f_{\mathcal{E}})$ whose Fargues filtration induces a well-behaved (=split) filtration on $\mathcal{E}_{2}$. We hope that Fargues's functor restricts to an equivalence of $\otimes$-categories $\underline{\mathcal{E}}:\Mod_{A,f}^{\varphi,\ast}\otimes E\rightarrow\Modif_{X}^{ad,\ast}$. This would imply that for any $M\in\Mod_{A,f}^{\varphi,\ast}$, $\underline{\mathcal{E}}$ maps the Fargues filtration of $M$ to the Fargues filtration of $\underline{\mathcal{E}}(M)$, and thus also $t_{F,\infty}(M)=t_{F}(\underline{\mathcal{E}}(M))$ (see~Proposition~\ref{prop:SpecialSubCats}).
\end{rema}

\subsection{The analytic construction}\label{subsec:AnalConstruction}

\subsubsection{}

In a category $\sfC$ with duals and effective objects, let us say that an object $X$ is anti-effective if its dual $X^{\vee}$ is effective. We denote by $\sfC^{\geq}$ and $\sfC^{\leq}$ the full subcategories of effective and anti-effective objects in $\sfC$.

\subsubsection{}

We equip $A$ with its $(\pi,[\varpi])$-topology. Following~\cite[12.2]{ScWe15}, we give names to four special points of $\Spa(A)=\Spa(A,A)$, labeled by their residue fields: $y_{\mathbb{F}}$, $y_{C^{\flat}}$, $y_{L}$ and $y_{C}$, corresponding respectively to the trivial valuation on the residue field $\mathbb{F}$ of $A$ and to the fixed valuations on the $A$-algebras $C^{\flat}$, $L$ and $C$. Then $y_{\mathbb{F}}$ is the unique non-analytic point of $\Spa(A)$ and the complement $\mathscr{Y}=\Spa(A)\setminus\{y_{\mathbb{F}}\}$ is equipped with a continuous surjective map $\kappa:\mathscr{Y}\rightarrow[-\infty,+\infty]$ defined by
\[
\kappa(y)\eqd\log_{q}\left(\frac{\log\left|[\varpi](\tilde{y})\right|}{\log\left|\pi(\tilde{y})\right|}\right)
\]
where $\tilde{y}$ is the maximal generalization of $y$, see~\cite[12.2]{ScWe15}. We have
\[
\kappa(y_{C^{\flat}})=-\infty,\qquad\kappa(y_{C})=0,\qquad\kappa(y_{L})=+\infty.
\]
The Frobenius $\varphi$ of $A$ induces an automorphism $\Spa(\varphi)$ of $\Spa(A)$ and $\mathscr{Y}$, which we still denote by $\varphi$. It fixes $y_{\mathbb{F}}$, $y_{C^{\flat}}$ and $y_{L}$, but not $y_{C}$. We set $y_{i}=\varphi^{i}(y_{C})$ for every $i\in\mathbb{Z}$, so that $\kappa(y_{i})=i$ since more generally $\kappa(\varphi(y))=\varphi(\kappa(y))$ for every $y\in\mathscr{Y}$, where the automorphism $\varphi$ of $[-\infty,+\infty]$ maps $x$ to $x+1$ (and fixes $\pm\infty$). Thus $y_{0}=y_{C}$ while $y_{-1}$ corresponds to $A\twoheadrightarrow\mathcal{O}_{C^{\prime}}\hookrightarrow C'$. For any interval $I\subset[-\infty,+\infty]$, we denote by $\mathscr{Y}_{I}$ the interior of the pre-image of $I$ under $\kappa$. We set
\[
\mathscr{Y}^{+}\eqd\mathscr{Y}_{]-\infty,+\infty]},\quad\mathscr{Y}^{-}\eqd\mathscr{Y}_{[-\infty,+\infty[}\quad\text{and}\quad\mathscr{Y}^{\circ}\eqd\mathscr{Y}^{+}\cap\mathscr{Y}^{-}=\mathscr{Y}_{]-\infty,+\infty[}.
\]

\subsubsection{}

By~\cite[13.1.1]{ScWe15}, $\mathscr{Y}$ is an honest (or sheafy) adic space. This means that the presheaf $\mathscr{O}_{\mathscr{Y}}$ of analytic functions on $\mathscr{Y}$ is a sheaf on $\mathscr{Y}$. Thus there is a well-defined $\otimes$-category $\Bun_{\mathscr{Y}_{I}}$ of vector bundles on $\mathscr{Y}_{I}$. A $\varphi$-equivariant bundle on $\mathscr{Y}_{I}$ is a pair $(\mathscr{E},\varphi_{\mathscr{E}})$ where $\mathscr{E}$ is a vector bundle on $\mathscr{Y}_{I}$ and $\varphi_{\mathscr{E}}:\varphi^{\ast}\mathscr{E}\vert_{\mathscr{Y}_{\varphi^{-1}(I)\cap I}}\rightarrow\mathscr{E}\vert_{\mathscr{Y}_{\varphi^{-1}(I)\cap I}}$ is an isomorphism. This defines a $\otimes$-category $\Bun_{\mathscr{Y}_{I}}^{\varphi}$. By~\cite{Ke16}, the adic subspace $\mathscr{Y}^{\circ}$ of $\mathscr{Y}$ is strongly Noetherian. Thus for any interval $I\subset\mathopen]-\infty,+\infty\mathclose[$, there is also a well-behaved abelian category $\Coh_{\mathscr{Y}_{I}}$ of coherent sheaves on $\mathscr{Y}_{I}$. A modification of vector bundles on $\mathscr{Y}_{I}$ is a monomorphism $f:\mathscr{E}_{1}\hookrightarrow\mathscr{E}_{2}$ of vector bundles on $\mathscr{Y}_{I}$ whose cokernel is a coherent sheaf supported at $\{y_{i}:i\in\mathbb{Z}\}\cap\mathscr{Y}_{I}$. Similarly, there is a notion of $\varphi$-equivariant modification of $\varphi$-equivariant vector bundles on $\mathscr{Y}_{I}$.

\subsubsection{}\label{subsec:FromBKF2Modif*}

By~\cite[3.9]{Ke16b}, the global section functor yields an equivalence of $\otimes$-categories $\Gamma(\mathscr{Y},-):\Bun_{\mathscr{Y}}\rightarrow\Mod_{A,f}$ with inverse $M\mapsto M\otimes_{A}\mathscr{O}_{\mathscr{Y}}$. In particular, every vector bundle $\mathscr{E}$ over $\mathscr{Y}$ is actually finite and free. Let $\Modif_{\mathscr{Y}}^{\star}$ be the $\otimes$-category of pairs $(\mathscr{E},\psi_{\mathscr{E}})$ where $\mathscr{E}$ is a vector bundle on $\mathscr{Y}$ and $\psi_{\mathscr{E}}:\mathscr{E}\rightarrow\varphi^{\ast}\mathscr{E}$ is a modification supported at $\{y_{-1}\}$, i.e.~$\psi_{\mathscr{E}}$ is an isomorphism over $\mathscr{Y}\setminus\{y_{-1}\}$. Then plainly
\[
\xymatrix@C=2pc@R=0.2pc{\Mod_{A,f}^{\varphi,\leq}\ar@{<->}[r] & \Modif_{\mathscr{Y}}^{\star}\\
(M,\varphi_{M})\ar@{|->}[r] & \left(\varphi_{M}^{-1}:M\rightarrow\varphi^{\ast}M\right)\otimes_{A}\mathscr{O}_{\mathscr{Y}}\\
\left(\Gamma(\mathscr{Y},\mathscr{E}),\Gamma(\mathscr{Y},\psi_{\mathscr{E}})^{-1}\right) & (\mathscr{E},\psi_{\mathscr{E}})\ar@{|->}[l] }
\]
are mutually inverse equivalences of $\otimes$-categories.

\subsubsection{}\label{subsec:FromModif*2ModifAn}

Let $\Modif_{\mathscr{Y}^{-},\mathscr{Y}^{+}}^{\varphi,\geq}$ be the $\otimes$-category of triples $(\mathscr{E}^{-},\mathscr{E}^{+},f_{\mathscr{E}})$ where $\mathscr{E}^{-}$ and $\mathscr{E}^{+}$ are $\varphi$-equivariant bundles over respectively $\mathscr{Y}^{-}$ and $\mathscr{Y}^{+}$ while $f_{\mathscr{E}}:\mathscr{E}^{-}\vert_{\mathscr{Y}^{\circ}}\rightarrow\mathscr{E}^{+}\vert_{\mathscr{Y}^{\circ}}$ is a $\varphi$-equivariant modification between their restriction to $\mathscr{Y}^{\circ}=\mathscr{Y}^{+}\cap\mathscr{Y}^{-}$. We claim that there are mutually inverse equivalences of $\otimes$-categories
\[
\xymatrix@R=0.2pc{\Modif_{\mathscr{Y}}^{\star}\ar@{<->}[r] & \Modif_{\mathscr{Y}^{-},\mathscr{Y}^{+}}^{\varphi,\geq}\\
(\mathscr{E},\psi_{\mathscr{E}})\ar@{<->}[r] & \left(\mathscr{E}^{-},\mathscr{E}^{+},f_{\mathscr{E}}\right) }
\]
Starting on the left hand side, set $\mathscr{E}(i)=(\varphi^{i})^{\ast}\mathscr{E}$ and define $\theta_{i}:\mathscr{E}(i)\hookrightarrow\mathscr{E}(i+1)$ by $\theta_{i}=(\varphi^{i})^{\ast}(\theta_{0})$ for $i\in\mathbb{Z}$ with $\theta_{0}=\psi_{\mathscr{E}}:\mathscr{E}(0)\rightarrow\mathscr{E}(1)$. Note that $\theta_{0}$ is a modification supported at $\{y_{-1}\}$, thus $\theta_{i}$ is a modification supported at $\{y_{-i-1}\}$ for all $i\in\mathbb{Z}$. As in~\cite[\S 4.4]{Fa15}, the following commutative diagram of vector bundles on $\mathscr{Y}$
\[
\xymatrix@R=2pc@C=1.6pc{\cdots\ar[r] & \varphi^{\ast}\mathscr{\mathscr{E}}(-2)\ar[r]^{\theta_{-1}} & \varphi^{\ast}\mathscr{\mathscr{E}}(-1)\ar[r]^{\theta_{0}} & \varphi^{\ast}\mathscr{E}(0)\ar[r]^{\theta_{1}} & \varphi^{\ast}\mathscr{\mathscr{E}}(1)\ar[r]^{\theta_{2}} & \varphi^{\ast}\mathscr{\mathscr{E}}(2)\ar[r]^{\theta_{3}} & \cdots\\
\cdots\ar[r] & \mathscr{\mathscr{E}}(-2)\ar[r]^{\theta_{-2}}\ar[u]^{\theta_{-2}} & \mathscr{\mathscr{E}}(-1)\ar[r]^{\theta_{-1}}\ar[u]^{\theta_{-1}} & \mathscr{E}(0)\ar[r]^{\theta_{0}}\ar[u]^{\theta_{0}} & \mathscr{\mathscr{E}}(1)\ar[r]^{\theta_{1}}\ar[u]^{\theta_{1}} & \mathscr{\mathscr{E}}(2)\ar[r]^{\theta_{2}}\ar[u]^{\theta_{2}} & \cdots }
\]
defines two $\varphi$-equivariant sheaves on $\mathscr{Y}$, namely
\begin{align*}
\mathscr{\mathscr{E}}(-\infty) & \eqd \underleftarrow{\lim}_{i\geq0}\mathscr{\mathscr{E}}(-i) = \cap_{i\geq0}\mathscr{\mathscr{E}}(-i)\\
\mathscr{\mathscr{E}}(+\infty) & \eqd \underrightarrow{\lim}_{i\geq0}\mathscr{E}(+i) = \cup_{i\geq0}\mathscr{\mathscr{E}}(+i)
\end{align*}
whose inverse Frobenius mappings
\[
\varphi_{\mathscr{\mathscr{E}}(-\infty)}^{-1}:\mathscr{\mathscr{E}}(-\infty)\rightarrow\varphi^{\ast}\mathscr{\mathscr{E}}(-\infty)\quad\text{and}\quad\varphi_{\mathscr{\mathscr{E}}(+\infty)}^{-1}:\mathscr{\mathscr{E}}(+\infty)\rightarrow\varphi^{\ast}\mathscr{\mathscr{E}}(+\infty)
\]
are induced by the vertical maps of the above diagram. Moreover,
\[
\mathscr{\mathscr{E}}(-\infty)\hookrightarrow\mathscr{\mathscr{E}}(i)\quad\text{and}\quad\mathscr{\mathscr{E}}(i)\hookrightarrow\mathscr{\mathscr{E}}(+\infty)
\]
are respectively isomorphisms outside $\{y_{j}:j\geq-i\}$ and $\{y_{j}:j<-i\}$, thus
\[
\mathscr{E}^{-}\eqd\mathscr{E}(-\infty)\vert_{\mathscr{Y}^{-}}\quad\text{and}\quad\mathscr{E}^{+}\eqd\mathscr{E}(+\infty)\vert_{\mathscr{Y}^{+}}
\]
are $\varphi$-equivariant vector bundles over respectively $\mathscr{Y}^{-}$ and $\mathscr{Y}^{+}$, and
\[
\left(f_{\mathscr{E}}:\mathscr{E}^{-}\vert_{\mathscr{Y}^{\circ}}\rightarrow\mathscr{E}^{+}\vert_{\mathscr{Y}^{\circ}}\right)\eqd\left(\mathscr{\mathscr{E}}(-\infty)\vert_{\mathscr{Y}^{\circ}}\rightarrow\mathscr{E}(0)\vert_{\mathscr{Y}^{\circ}}\rightarrow\mathscr{E}(+\infty)\vert_{\mathscr{Y}^{\circ}}\right)
\]
is a $\varphi$-equivariant modification as desired.

Conversely, starting from $(\mathscr{E}^{-},\mathscr{E}^{+},f_{\mathscr{E}})$ on the right hand side, we define a vector bundle $\mathscr{E}$ on $\mathscr{Y}$ by gluing $\mathscr{E}^{-}\vert_{\mathscr{Y}_{[-\infty,0[}}$ and $\mathscr{E}^{+}\vert_{\mathscr{Y}_{]-1,+\infty]}}$ along the isomorphism induced by the restriction of $f_{\mathscr{E}}$ to $\mathscr{Y}_{]-1,0[}$. Thus $\mathscr{E}\vert_{\mathscr{Y}^{\circ}}$ is the subsheaf of $\mathscr{E}^{+}\vert_{\mathscr{Y}^{\circ}}$ made of those sections whose restriction to $\mathscr{Y}_{]-\infty,0[}$ belong to the image of $f_{\mathscr{E}}$. Since $f_{\mathscr{E}}$ is a $\varphi$-equivariant modification, it follows that there is a commutative diagram
\[
\xymatrix{\mathscr{E}^{-}\vert_{\mathscr{Y}^{\circ}}\ar@{^{(}->}[r]\ar[d]^{\varphi_{\mathscr{E}^{-}}^{-1}} & \mathscr{E}\vert_{\mathscr{Y}^{\circ}}\ar@{^{(}->}[r]\ar[d]^{\psi_{\mathscr{E}}} & \mathscr{E}^{+}\vert_{\mathscr{Y}^{\circ}}\ar[d]^{\varphi_{\mathscr{E}^{+}}^{-1}}\\
\varphi^{\ast}\mathscr{E}^{-}\vert_{\mathscr{Y}^{\circ}}\ar@{^{(}->}[r] & \varphi^{\ast}\mathscr{E}\vert_{\mathscr{Y}^{\circ}}\ar@{^{(}->}[r] & \varphi^{\ast}\mathscr{E}^{+}\vert_{\mathscr{Y}^{\circ}} }
\]
We extend $\psi_{\mathscr{E}}$ to $\mathscr{Y}$ by setting $\psi_{\mathscr{E}}:=\varphi_{\mathscr{E}^{-}}^{-1}$ on $\mathscr{E}_{[-\infty,-1[}$ and $\psi_{\mathscr{E}}:=\varphi_{\mathscr{E}^{+}}^{-1}$ on $\mathscr{Y}_{]-1,+\infty]}$. Therefore $\psi_{\mathscr{E}}:\mathscr{E}\rightarrow\varphi^{\ast}\mathscr{E}$ is an isomorphism away from $\kappa^{-1}(-1)\cap\{y_{i}\}=\{y_{-1}\}$, i.e.~$\psi_{\mathscr{E}}$ is indeed a modification supported at $y_{-1}$.

One checks easily that these constructions yield mutually inverse $\otimes$-functors.

\subsubsection{}\label{subsec:VariantA=00005B1/pi=00005D}

Starting with $\Mod_{A[\frac{1}{\pi}],f}^{\varphi}$, we may analogously define $\otimes$-functors
\[
\xymatrix@R=0.5pc{\Mod_{A[\frac{1}{\pi}],f}^{\varphi,\leq}\ar[r] & \Modif_{\mathscr{Y}^{+}}^{\star}\ar@{<->}[r] & \Modif_{\mathscr{Y}^{\circ},\mathscr{Y}^{+}}^{\varphi,\geq}\\
(N,\varphi_{N})\ar@{|->}[r] & (\mathscr{E},\psi_{\mathscr{E}})\ar@{|->}[r] & (\mathscr{E}^{-},\mathscr{E}^{+},f_{\mathscr{E}}) }
\]
with the obvious definitions for the $\otimes$-categories $\Modif_{\mathscr{Y}^{+}}^{\star}$ and $\Modif_{\mathscr{Y}^{\circ},\mathscr{Y}^{+}}^{\varphi,\geq}$, where
\[
(\mathscr{E},\psi_{\mathscr{E}})\eqd(N,\varphi_{N}^{-1})\otimes_{A[\frac{1}{\pi}]}\mathscr{O}_{\mathscr{Y}^{+}}
\]
and $\mathscr{E}^{-}\in\Bun_{\mathscr{Y}^{\circ}}^{\varphi}$, $\mathscr{E}^{+}\in\Bun_{\mathscr{Y}^{+}}^{\varphi}$. While $\Modif_{\mathscr{Y}^{+}}^{\star}\leftrightarrow\Modif_{\mathscr{Y}^{\circ},\mathscr{Y}^{+}}^{\varphi,\geq}$ are still mutually inverse equivalences of $\otimes$-categories, it is not clear that the first functor is an equivalence. Indeed, the functor $\Mod_{A[\frac{1}{\pi}],f}\rightarrow\Bun_{\mathscr{Y}^{+}}$ is already not full.

\subsubsection{}\label{subsec:PhiBunOverY-}

As in~\cite[12.3.4]{ScWe15} and~\cite[8.5.3]{KeLi15}, there are equivalences of $\otimes$-categories
\[
\xymatrix@R=0pc{\Bun_{\mathscr{Y}^{-}}^{\varphi}\ar[r] & \Bun_{\mathscr{R}_{-}^{\mathrm{int}}}^{\varphi}\ar[r] & \Bun_{A(C)}^{\varphi}\ar[r] & \Bun_{\mathcal{O}_{E}}\\
\mathscr{E}\ar@{|->}[r] & \mathscr{E}_{y_{C^{\flat}}}\ar@{|->}[r] & \mathscr{E}_{y_{C^{\flat}}}^{\wedge}\ar@{|->}[r] & (\mathscr{E}_{y_{C^{\flat}}}^{\wedge})^{\varphi_{\mathscr{E}}=1} }
\]
where $\mathscr{R}_{-}^{\mathrm{int}}:=\lim_{r\mapsto-\infty}\Gamma(\mathscr{Y}_{[-\infty,r]},\mathscr{O}_{\mathscr{Y}})$ is the local ring of $\mathscr{Y}$ at $y_{C^{\flat}}$; this is the integral Robba ring, a Henselian discrete valuation ring with uniformizer $\pi$, residue field $C^{\flat}$ and completion $A(C)=W_{\mathcal{O}_{E}}(C^{\flat})$~\cite[1.8.2]{FaFo15}. The objects of the middle two categories are the finite free \'etale $\varphi$-modules $(N,\varphi_{N})$ over the indicated local rings, and the functor between them is the base change map (or $\pi$-adic completion) with respect to $\mathscr{R}_{-}^{\mathrm{int}}\hookrightarrow A(C)$. We have already encountered the last functor in~\ref{subsec:RisW(Kflat)}: it maps $(N,\varphi_{N})$ to $T=N^{\phi_{N}=1}$. The inverse $\otimes$-functor maps the finite free $\mathcal{O}_{E}$-module $T$ to the ``constant'' $\varphi$-bundle $(\mathscr{E}^{-},\varphi_{\mathscr{E}^{-}})=(T\otimes_{\mathcal{O}_{E}}\mathscr{O}_{\mathscr{Y}^{-}},\mathrm{Id}\otimes\varphi)$ over $\mathscr{Y}^{-}$. In particular, every $\varphi$-bundle over $\mathscr{Y}^{-}$ is actually finite free.

\subsubsection{}\label{subsec:FFDiagram}

There is also a commutative diagram of $\otimes$-categories~\cite[\S 11.4]{FaFo15}
\[
\xymatrix@R=2.5pc@C=5pc{\Bun_{\mathscr{R}_{+}^{\mathrm{int}}}^{\varphi}\ar@{.>}[d]_{{\,\cdot\,}\otimes\overline{B}} & \Bun_{\mathscr{Y}^{+}}^{\prime,\varphi}\ar@{.>}[l]_{({\,\cdot\,})_{y_{L}}}\ar[r]^{{\,\cdot\,}\vert\mathscr{Y}^{\circ}}\ar@<1ex>[d]^{\Gamma(\mathscr{Y}^{+},{\,\cdot\,})} & \Bun_{\mathscr{Y}^{\circ}}^{\varphi}\ar@<.7ex>[r]^{{\,\cdot\,}/\varphi}\ar@<.7ex>[d]^{\Gamma(\mathscr{Y}^{\circ},{\,\cdot\,})} & \Bun_{\mathscr{X}}\ar@<1ex>[l]^{\pi^{\ast}}\ar@<1ex>[d]^{({\,\cdot\,})^{alg}}\\
\Bun_{\overline{B}}^{\varphi} & \Bun_{B^{+}}^{\prime,\varphi}\ar[l]_{{\,\cdot\,}\otimes\overline{B}}\ar[r]^{{\,\cdot\,}\otimes B}\ar@<1ex>[u]^{{\,\cdot\,}\otimes\mathscr{O}_{\mathscr{Y}^{+}}} & \Bun_{B}^{\varphi}\ar[r]^{\mathcal{E}}\ar@<1ex>[u]^{{\,\cdot\,}\otimes\mathscr{O}_{\mathscr{Y}^{\circ}}} & \Bun_{X}\ar@<1ex>[u]^{({\,\cdot\,})^{an}} }
\]
in which all solid arrows are equivalences of $\otimes$-categories. In the first line,
\[
\mathscr{R}_{+}^{\mathrm{int}}\eqd\lim_{r\rightarrow\text{+\ensuremath{\infty}}}\,\Gamma(\mathscr{Y}_{[r,+\infty]},\mathscr{O}_{\mathscr{Y}})
\]
is the analog of the integral Robba ring $\mathscr{R}_{-}^{\mathrm{int}}$ with $y_{C^{\flat}}$ replaced by $y_{L}$, and
\[
\mathscr{X}\eqd\mathscr{Y}^{\circ}/\varphi^{\mathbb{Z}}
\]
is the adic version of the Fargues--Fontaine curve $X$, a strongly noetherian analytic space. There is a morphism of locally ringed space $\mathscr{X}\rightarrow X$ which induces pull-back $\otimes$-functors $({\,\cdot\,})^{an}:\Coh_{X}\rightarrow\Coh_{\mathscr{X}}$ and $({\,\cdot\,})^{an}:\Bun_{X}\rightarrow\Bun_{\mathscr{X}}$. The equivalence of $\otimes$-categories $\Bun_{\mathscr{Y}^{\circ}}^{\varphi}\leftrightarrow\Bun_{\mathscr{X}}$ maps a vector bundle on $\mathscr{X}$ to its pull-back through the $\varphi$-invariant morphism $\pi:\mathscr{Y}^{\circ}\rightarrow\mathscr{X}$, and maps a $\varphi$-bundle $\mathscr{E}$ on $\mathscr{Y}^{\circ}$ to the sheaf $\mathscr{E}/\varphi_{\mathscr{E}}$ of $\varphi_{\mathscr{E}}$-invariant sections of $\pi_{\ast}\mathscr{E}$. We denote by $\mathscr{E}\mapsto\mathscr{E}(d)$ the Tate twists on $\Bun_{\mathscr{X}}$ and $\Bun_{\mathscr{Y}^{\circ}}^{\varphi}$ corresponding to the Tate objects $\mathscr{O}_{\mathscr{X}}(1)=\mathcal{O}_{X}(1)^{an}$ and $\mathscr{O}_{\mathscr{Y}^{\circ}}(1)=\pi^{\ast}\mathscr{O}_{\mathscr{X}}(1)$. In the second line, the $A${[}$\frac{1}{\pi}]$-algebras
\[
B\hookleftarrow B^{+}\twoheadrightarrow\overline{B}
\]
are defined in~\cite[1.10]{FaFo15}. They are related to the adic space $\mathscr{Y}$ by
\[
B^{+}=\Gamma(\mathscr{Y}^{+},\mathscr{O}_{\mathscr{Y}})\quad\text{and}\quad B=\Gamma(\mathscr{Y}^{\circ},\mathscr{O}_{\mathscr{Y}}).
\]
Moreover, $\overline{B}$ is a local domain with residue field $L$ which is also a quotient of $\mathscr{R}_{+}^{\mathrm{int}}$. The Fargues--Fontaine curve $X$ equals $\Proj(P)$ where $P:=\bigoplus_{d\geq0}P_{d}$ with
\[
P_{d}\eqd\Gamma(X,\mathcal{O}_{X}(d))=\Gamma(\mathscr{X},\mathscr{O}_{\mathscr{X}}(d))=B^{\varphi=\pi^{d}}=(B^{+})^{\varphi=\pi^{d}}.
\]
The $\otimes$-functor $\mathcal{E}:\Bun_{B}^{\varphi}\rightarrow\Bun_{X}$ maps a finite projective \'etale $\varphi$-module $(N,\varphi_{N})$ to the quasi-coherent sheaf on $X$ associated with the graded $P$-module $\bigoplus_{d\geq0}N^{\varphi_{N}=\pi^{d}}$. The $\otimes$-functor $({\,\cdot\,})^{alg}:\Bun_{\mathscr{X}}\rightarrow\Bun_{X}$ maps an adic vector bundle $\mathscr{E}$ on $\mathscr{X}$ to the quasi-coherent sheaf on $X$ associated with the graded $P$-module $\bigoplus_{d\geq0}\Gamma(\mathscr{X},\mathscr{E}(d))$. In the second column of our diagram, the primes refer to the full $\otimes$-subcategories of finite free objects in the relevant $\otimes$-categories of $\varphi$-bundles. Thus plainly,
\[
\Gamma(\mathscr{Y}^{+},{\cdot\,}):\Bun_{\mathscr{Y}^{+}}^{\prime,\varphi}\longleftrightarrow\Bun_{B^{+}}^{\prime,\varphi}:({\,\cdot\,}\otimes\mathscr{O}_{\mathscr{Y}^{+}})
\]
are mutually inverse equivalences of $\otimes$-categories. The $\otimes$-functor
\[
({\,\cdot\,})_{y_{L}}:\Bun_{\mathscr{Y}^{+}}^{\prime,\varphi}\rightarrow\Bun_{\mathscr{R}_{+}^{\mathrm{int}}}^{\varphi}
\]
maps a $\varphi$-bundle $(\mathscr{E},\varphi_{\mathscr{E}})$ on $\mathscr{Y}^{+}$ to the $\varphi$-bundle $(\mathscr{E}_{y_{L}},\varphi_{\mathscr{E}_{y_{L}}})$ over $\mathscr{R}_{+}^{\mathrm{int}}$ defined by
\[
\mathscr{E}_{y_{L}}\eqd\lim_{r\rightarrow\text{+\ensuremath{\infty}}}\,\Gamma(\mathscr{Y}_{[r,+\infty]},\mathscr{E})\quad\text{with}\quad\varphi_{\mathscr{E}_{y_{L}}}\eqd\lim_{r\rightarrow\text{+\ensuremath{\infty}}}\,\varphi_{\mathscr{E}}\vert_{\mathscr{Y}_{[r,+\infty]}}
\]
and the $\otimes$-functors labelled $-\otimes\overline{B}$ or $-\otimes B$ are induced by tensorisation with the relevant $\varphi$--equivariant morphisms of $B\hookleftarrow B^{+}\rightarrow\mathscr{R}_{+}^{\mathrm{int}}\twoheadrightarrow\overline{B}$. The $\otimes$-functors
\[
\Bun_{\overline{B}}^{\varphi}\stackrel{1.7}{\longleftarrow}\Bun_{B^{+}}^{\prime,\varphi}\stackrel{1.9}{\longrightarrow}\Bun_{X}\stackrel{3.1}{\longrightarrow}\Bun_{\mathscr{X}}\quad\text{and}\quad\Bun_{\mathscr{Y}^{\circ}}^{\varphi}\stackrel{2.2}{\longrightarrow}\Bun_{B}^{\varphi}
\]
are equivalence of $\otimes$-categories by the indicated references in~\cite[\S 11]{FaFo15}, and so are therefore also all of the above solid arrow functors. In particular, every $\varphi$-bundle on $\mathscr{Y}^{\circ}$ is finite free and extends uniquely to a finite free $\varphi$-bundle on $\mathscr{Y}^{+}$.

\subsubsection{}\label{subsec:PhiBunY=0000B0}

This is in sharp contrast with what happens at $y_{C^{\flat}}$: not every $\varphi$-bundle on $\mathscr{Y}^{\circ}$ extends to $\mathscr{Y}^{-}$, and those who do have many extensions. This is related to semi-stability as follows. Let $({\,\cdot\,})^{+}:\Bun_{\mathscr{Y}^{\circ}}^{\varphi}\rightarrow\Bun_{\mathscr{Y}^{+}}^{\prime,\varphi}$ be a chosen $\otimes$-inverse of the restriction $\otimes$-functor $\Bun_{\mathscr{Y}^{+}}^{\prime,\varphi}\rightarrow\Bun_{\mathscr{Y}^{\circ}}^{\varphi}$. We then have three $\otimes$-functors
\[
\xymatrix@C=1pc@R=1pc{ & \Bun_{\mathscr{R}_{-}}^{\varphi} & \text{given by} & \underset{r\mapsto-\infty}{\lim}\Gamma\left(\mathscr{Y}_{]-\infty,r]},{\cdot\,}\right)\simeq\Gamma\left(\mathscr{Y}^{\circ},{\cdot\,}\right)\otimes_{B}\mathscr{R}_{-}\\
\Bun_{\mathscr{Y}^{\circ}}^{\varphi}\ar[ru]\ar[r]\ar[rd] & \Bun_{X} & \text{given by} & ({\,\cdot\,}/\varphi)^{alg}\simeq\mathcal{E}\circ\Gamma(\mathscr{Y}^{\circ},{\cdot\,})\\
\Bun_{\mathscr{Y}^{+}}^{\prime,\varphi}\ar[r]\ar@{<->}[u] & \Bun_{L}^{\varphi} & \text{given by} & ({\,\cdot\,})_{y_{L}}^{+}\otimes_{\mathscr{R}_{+}^{\mathrm{int}}}L\simeq\Gamma(\mathscr{Y}^{+},({\,\cdot\,})^{+})\otimes_{B^{+}}L }
\]
where $\mathscr{R}_{-}:=\lim_{r\mapsto-\infty}\Gamma(\mathscr{Y}_{]-\infty,r]},\mathscr{O}_{\mathscr{Y}})$ is the Robba ring; this is a B\'ezout ring by~\cite[3.5.8]{FaFo15} or~\cite[2.9.6]{Ke05}. The first $\otimes$-functor is an equivalence of categories by~\cite[11.2.22]{FaFo15}, and we have already seen that so is the second. The third one is not: $\Bun_{L}^{\varphi}$ is abelian semi-simple while $\Bun_{X}$ (along with $\Bun_{\mathscr{Y}^{\circ}}^{\varphi}$ and $\Bun_{\mathscr{R}_{-}}^{\varphi}$) is only quasi-abelian, and not at all semi-simple. The three target categories are quasi-abelian $\otimes$-categories, with a Harder--Narasimhan formalism compatible with $\otimes$-products: this is due respectively to Kedlaya~\cite{Ke05}, Fargues and Fontaine~\cite{FaFo15} (see~\ref{subsec:Newton4BundlesOnCurve}), and to the Dieudonn\'e-Manin classification of isocrystals, which actually gives rise to a pair of opposed Newton slope filtrations $\mathcal{F}_{N}$ and $\mathcal{F}_{N}^{\iota}$ (see~\ref{subsec:RisL}). These formalisms are compatible, provided that we pick the opposed Newton filtration $\mathcal{F}_{N}^{\iota}$ on $\Bun_{L}^{\varphi}$.

The compatibility of the slope filtrations along $\Bun_{X}\simeq\Bun_{\mathscr{Y}^{\circ}}^{\varphi}\simeq\Bun_{\mathscr{R}_{-}}^{\varphi}$ is build up in the proof of~\cite[11.2.22]{FaFo15}. Their compatibility along $\Bun_{X}\simeq\Bun_{\mathscr{Y}^{\circ}}^{\varphi}\rightarrow\Bun_{L}^{\varphi}$ can be seen as follows. Starting with a $\varphi$-bundle $(\mathscr{E},\varphi_{\mathscr{E}})$ on $\mathscr{Y}^{\circ}$, set
\[
(M,\varphi_{M})=\Gamma\left(\mathscr{Y}^{+},(\mathscr{E},\varphi_{\mathscr{E}})^{+}\right).
\]
This is a finite free \'etale $\varphi$-module over $B^{+}$ and $(N,\varphi_{N})=(M,\varphi_{M})\otimes_{B^{+}}L$ is the image of $(\mathscr{E},\varphi_{\mathscr{E}})$ in $\Bun_{L}^{\varphi}$. Fix a section $s:\mathbb{F}\hookrightarrow\mathcal{O}_{C^{\flat}}$ of the quotient map $\mathcal{O}_{C^{\flat}}\twoheadrightarrow\mathbb{F}$. This gives rise to sections of $A\twoheadrightarrow\mathcal{O}_{L}$ and $A[\frac{1}{\pi}]\hookrightarrow B^{+}\rightarrow\mathscr{R}_{+}^{\mathrm{int}}\twoheadrightarrow\overline{B}\twoheadrightarrow L$, which we still denote by $s$. Then $(M,\varphi_{M})$ is non-canonically isomorphic to $(N,\varphi_{N})\otimes_{L,s}B^{+}$ by~\cite[\S 11.1]{FaFo15}, thus $(\mathscr{E}/\varphi_{\mathscr{E}})^{alg}\simeq\mathcal{E}((M,\varphi_{M})\otimes_{B^{+}}B)$ is non-canonically isomorphic to $\mathcal{E}^{s}(N,\varphi_{N}):=\mathcal{E}((N,\varphi_{N})\otimes_{L,s}B)$. Our claim now follows from~\cite[\S 8.2.4]{FaFo15}, where this $\otimes$-functor $\mathcal{E}^{s}:\Bun_{L}^{\varphi}\rightarrow\Bun_{X}$ is denoted by $\mathcal{E}$.

Now by Kedlaya's theory, we have equivalences of $\otimes$-categories
\[
\xymatrix@C=4pc{\Vect_{E} & \Bun_{\mathscr{R}_{-}^{\mathrm{int}}}^{\varphi}\otimes E\ar@{<->}[l]\ar[r]^-{-\otimes\mathscr{R}_{-}} & \Bun_{\mathscr{R}_{-}}^{\varphi,0}}
\]
where $\Bun_{\mathscr{R}_{-}}^{\varphi,0}$ is the full $\otimes$-subcategory of slope $0$ semi-stable objects in $\Bun_{\mathscr{R}_{-}}^{\varphi}$. The composite $\otimes$-functor $\Vect_{E}\rightarrow\Bun_{\mathscr{R}_{-}}^{\varphi,0}$ maps $V$ to $(V\otimes_{E}\mathscr{R}_{-},\mathrm{Id}\otimes\varphi)$ with inverse $(N,\varphi_{N})\mapsto N^{\varphi_{N}=1}$. It follows that we have equivalences of $\otimes$-categories
\[
\xymatrix@C=4pc{\Vect_{E} & \Bun_{\mathscr{Y}^{-}}^{\varphi}\otimes E\ar@{<->}[l]\ar[r]^-{-\vert\mathscr{Y}^{\circ}} & \Bun_{\mathscr{Y}^{\circ}}^{\varphi,0}}
\]
where $\Bun_{\mathscr{Y}^{\circ}}^{\varphi,0}$ is the full $\otimes$-subcategory of slope $0$ semi-stable objects in $\Bun_{\mathscr{Y}^{\circ}}^{\varphi}$. The composite functor $\Vect_{E}\rightarrow\Bun_{\mathscr{Y}^{\circ}}^{\varphi,0}$ maps $V$ to $(V\otimes_{E}\mathscr{O}_{\mathscr{Y}^{\circ}},\mathrm{Id}\otimes\varphi)$ with inverse $(\mathscr{E},\varphi_{\mathscr{E}})\mapsto\Gamma(\mathscr{Y}^{\circ},\mathscr{E})^{\varphi_{\mathscr{E}}=1}$. In other words, a $\varphi$-bundle $(\mathscr{E},\varphi_{\mathscr{E}})$ over $\mathscr{Y}^{\circ}$ extends to a $\varphi$-bundle over $\mathscr{Y}^{-}$ if and only if it is semi-stable of slope $0$ and then, there is a functorial bijective correspondence between the set of all such extensions and the set of all $\mathcal{O}_{E}$-lattices $T$ in $V=\Gamma(\mathscr{Y}^{\circ},\mathscr{E})^{\varphi_{\mathscr{E}}=1}$, given by $T\mapsto\left(T\otimes\mathscr{O}_{\mathscr{Y}^{-}},\mathrm{Id}\otimes\varphi\right)$.

\subsubsection{}\label{subsec:ComputeAnalFunct}

We shall now compute the equivalence of $\otimes$-categories
\[
\xymatrix@R=0.2pc{\Mod_{A,f}^{\varphi,\leq}\ar[r] & \Modif_{\mathscr{Y}}^{\star}\ar[r] & \Modif_{\mathscr{Y}^{-},\mathscr{Y}^{+}}^{\varphi,\geq}\\
(M,\varphi_{M})\ar@{|->}[r] & \left(\mathscr{E},\psi_{\mathscr{E}}\right)\ar@{|->}[r] & \left(\mathscr{E}^{-},\mathscr{E}^{+},f_{\mathscr{E}}\right) }
\]
Starting with the anti-effective finite free BKF-module $(M,\varphi_{M})$ over $A$, set
\[
T=(M\otimes_{A}A(C))^{\varphi_{M}\otimes\varphi=1}\quad\text{and}\quad(\overline{M},\varphi_{\overline{M}})=(M,\varphi_{M})\otimes_{A}\overline{B}.
\]
Thus $T$ is a finite free $\mathcal{O}_{E}$-module and $(\overline{M},\varphi_{\overline{M}})$ is a finite free \'etale $\varphi$-module over $\overline{B}$ (since $\xi'$ is invertible in $\overline{B}$). By~\cite[4.26]{BaMoSc16} and its proof, the canonical isomorphism
\[
\left(T\otimes_{\mathcal{O}_{E}}A(C),\mathrm{Id}\otimes\varphi\right)\simeq\left(M\otimes_{A}A(C),\varphi_{M}\otimes\varphi\right)
\]
descends to an isomorphism over the subring $A[\frac{1}{\mu}]\subset A(C)$,
\[
\eta_{M}^{-}\!\left[\frac{1}{\mu}\right]:T\otimes_{\mathcal{O}_{E}}A\!\left[{\frac{1}{\mu}}\right]\stackrel{\simeq}{\longrightarrow}M\otimes_{A}A\!\left[{\frac{1}{\mu}}\right]
\]
which is induced by a $\varphi^{-1}$-invariant $\mathcal{O}_{E}$-linear morphism
\[
\eta_{M}^{-}:T\hookrightarrow M.
\]
The latter gives a morphism of modifications of vector bundles on $\mathscr{Y}$,
\[
\eta_{M}^{-}\otimes\mathscr{O}_{\mathscr{Y}}:\left(T\otimes\mathscr{O}_{\mathscr{Y}},\mathrm{Id}\otimes\varphi^{-1}\right)\hookrightarrow\left(\mathscr{E},\psi_{\mathscr{E}}\right)
\]
whose restriction to $\mathscr{Y}^{-}$ factors through a morphism of $\varphi$-bundles over $\mathscr{Y}^{-}$,
\[
f_{M}^{-}:\left(T\otimes\mathscr{O}_{\mathscr{Y}^{-}},\mathrm{Id}\otimes\varphi\right)\hookrightarrow\left(\mathscr{E}^{-},\varphi_{\mathscr{E}^{-}}\right).
\]
Since $\mu$ is invertible on $\mathscr{Y}_{[-\infty,0[}$, both $\eta_{M}^{-}\otimes\mathscr{O}_{\mathscr{Y}}$ and $f_{M}^{-}$ are isomorphisms over $\mathscr{Y}_{[-\infty,0[}$. In particular, the localization of $f_{M}^{-}$ at $y_{C^{\flat}}$ is an isomorphism, and so is therefore $f_{M}^{-}$ itself by~\ref{subsec:PhiBunOverY-}.

On the other hand, pick any finite free \'etale $\varphi$-module $(D^{+},\varphi_{D^{+}})$ over $B^{+}$ reducing to $(\overline{M},\varphi_{\overline{M}})$ over $\overline{B}$. By~\cite[4.26]{Fa15} applied to the effective dual BKF-module $(M,\varphi_{M})^{\vee}=(M^{\vee},\varphi_{M}^{\vee-1})$, there is a unique $\varphi^{-1}$-equivariant morphism
\[
\eta_{M}^{+}:\left(M\otimes_{A}B^{+},\varphi_{M}^{-1}\otimes\varphi^{-1}\right)\rightarrow\left(D^{+},\varphi_{D^{+}}^{-1}\right)
\]
reducing to the given isomorphism $(\overline{M},\varphi_{\overline{M}}^{-1})\simeq\left(D^{+},\varphi_{D^{+}}^{-1}\right)\otimes_{B^{+}}\overline{B}$. As above, the latter yields a morphism of modifications of vector bundles over $\mathscr{Y}^{+}$,
\[
\eta_{M}^{+}\otimes\mathscr{O}_{\mathscr{Y}}:\left(\mathscr{E},\psi_{\mathscr{E}}\right)\vert_{\mathscr{Y}^{+}}\rightarrow\left(D^{+},\varphi_{D^{+}}^{-1}\right)\otimes_{B^{+}}\mathscr{O}_{\mathscr{Y}^{+}}
\]
which induces a morphism of $\varphi$-bundles over $\mathscr{Y}^{+}$,
\[
f_{M}^{+}:\left(\mathscr{E}^{+},\varphi_{\mathscr{E}^{+}}\right)\rightarrow\left(D^{+},\varphi_{D^{+}}\right)\otimes_{B^{+}}\mathscr{O}_{\mathscr{Y}^{+}}.
\]
By~\cite[4.31]{Fa15}, $\eta_{M}^{+}\otimes\mathscr{O}_{\mathscr{Y}}$ restricts to an isomorphism over $\mathscr{Y}_{[r,+\infty]}$ for $r\gg0$. Since also $\mathscr{E}\vert_{\mathscr{Y}^{+}}\hookrightarrow\mathscr{E}^{+}$ is an isomorphism over $\mathscr{Y}_{]-1,\infty]}$, it follows that $f_{M}^{+}$ is an isomorphism over $\mathscr{Y}_{[r,+\infty]}$. But then $(\varphi^{i})^{\ast}(f_{M}^{+})\simeq f_{M}^{+}$ is an isomorphism over $\mathscr{Y}_{[r-i,+\infty]}$ for all $i\geq0$, thus $f_{M}^{+}$ is an isomorphism over the whole of $\mathscr{Y}^{+}$.

Finally, let $f_{M}:T\otimes_{\mathcal{O}_{E}}\mathscr{O}_{\mathscr{Y}^{\circ}}\rightarrow D^{+}\otimes_{B^{+}}\mathscr{O}_{\mathscr{Y}^{\circ}}$ be the $\varphi$-equivariant morphism
\[
\xymatrix@C=45pt{T\otimes_{\mathcal{O}_{E}}\mathscr{O}_{\mathscr{Y}^{\circ}}\ar[r]^{\eta_{M}^{-}\otimes\mathscr{O}_{\mathscr{Y}^{\circ}}} & M\otimes_{A}\mathscr{O}_{\mathscr{Y}^{\circ}}\ar[r]^{\eta_{M}^{+}\otimes\mathscr{O}_{\mathscr{Y}^{\circ}}} & D^{+}\otimes_{B^{+}}\mathscr{O}_{\mathscr{Y}^{\circ}}}
\]
We thus have shown that $f_{M}^{-}$ and $f_{M}^{+}$ induce an isomorphism
\[
\left(\mathscr{E}^{-},\mathscr{E}^{+},f_{\mathscr{E}}\right)\simeq\left(T\otimes_{\mathcal{O}_{E}}\mathscr{O}_{\mathscr{Y}^{-}},D^{+}\otimes_{B^{+}}\mathscr{O}_{\mathscr{Y}^{+}},f_{M}\right).
\]
In particular, $\mathscr{E}^{+}\simeq D^{+}\otimes_{B^{+}}\mathscr{O}_{\mathscr{Y}^{+}}$ is finite free (we did not know this so far) and
\[
\eta_{M}^{+}\otimes\mathscr{O}_{\mathscr{Y}^{+}}:M\otimes_{A}\mathscr{O}_{\mathscr{Y}^{+}}\rightarrow D^{+}\otimes_{B^{+}}\mathscr{O}_{\mathscr{Y}^{+}}
\]
is an isomorphism over $\mathscr{Y}_{]-1,+\infty]}$. The freeness of $\mathscr{E}^{+}$ also yields a canonical choice for the finite free lift $D^{+}$ of $\overline{M}=M\otimes_{A}\overline{B}$: we may take $D^{+}=\Gamma(\mathscr{Y}^{+},\mathscr{E}^{+})$ with the isomorphism $\overline{M}\simeq D^{+}\otimes_{B^{+}}\overline{B}$ induced by $\mathscr{E}\vert_{\mathscr{Y}^{+}}\hookrightarrow\mathscr{E}^{+}$.

\subsubsection{}\label{subsec:ConstruFarguesFunctEff}

The discussion above shows that the $\otimes$-functor
\[
\xymatrix@C=2pc{\Mod_{A[\frac{1}{\pi}],f}^{\varphi,\leq}\ar[r] & \Modif_{\mathscr{Y}^{+}}^{\star}\ar[r] & \Modif_{\mathscr{Y}^{\circ},\mathscr{Y}^{+}}^{\varphi,\geq}}
\]
induces an equivalence of $\otimes$-categories
\[
\xymatrix{\Mod_{A,f}^{\varphi,\leq}\otimes E\ar[r] & \Modif_{\mathscr{Y}^{\circ},\mathscr{Y}^{+}}^{\varphi,ad,\geq}}
\]
where $\Modif_{\mathscr{Y}^{\circ},\mathscr{Y}^{+}}^{\varphi,ad,\geq}$ is the full $\otimes$-subcategory of objects $(\mathscr{E}^{-},\mathscr{E}^{+},f_{\mathscr{E}})$ in $\Modif_{\mathscr{Y}^{\circ},\mathscr{Y}^{+}}^{\varphi,\geq}$ such that $\mathscr{E}^{-}\in\Bun_{\mathscr{Y}^{\circ}}^{\varphi}$ belongs to $\Bun_{\mathscr{Y}^{\circ}}^{\varphi,0}$. Moreover for any such object, $\mathscr{E}^{+}$ is actually finite free, and ${\,\cdot\,}\vert_{\mathscr{Y}^{\circ}}:\Bun_{\mathscr{Y}^{+}}^{\prime,\varphi}\stackrel{\simeq}{\rightarrow}\Bun_{\mathscr{Y}^{\circ}}^{\varphi}$ thus induces an equivalence
\[
\xymatrix@R=0.5pc{\Modif_{\mathscr{Y}^{\circ},\mathscr{Y}^{+}}^{\varphi,ad,\geq}\ar[r] & \Modif_{\mathscr{Y}^{\circ},\mathscr{Y}^{\circ}}^{\varphi,ad,\geq}\\
(\mathscr{E}^{-},\mathscr{E}^{+},f_{\mathscr{E}})\ar@{|->}[r] & (\mathscr{E}^{-},\mathscr{E}^{+}\vert_{\mathscr{Y}^{\circ}},f_{\mathscr{E}}) }
\]
of $\otimes$-categories. Finally, the equivalence of $\otimes$-categories
\[
\left(\xymatrix{\Bun_{\mathscr{Y}^{\circ}}^{\varphi}\ar[r]^{{\,\cdot\,}/\varphi} & \Bun_{\mathscr{X}}\ar[r]^{({\,\cdot\,})^{\mathrm{alg}}} & \Bun_{X}}
\right)=\left(\xymatrix{\Bun_{\mathscr{Y}^{\circ}}^{\varphi}\ar[r]^{\Gamma(\mathscr{Y}^{\circ},{\,\cdot\,})} & \Bun_{B}^{\varphi}\ar[r]^{\mathcal{E}} & \Bun_{X}}
\right)
\]
induces equivalences of $\otimes$-categories
\[
\xymatrix{\Modif_{\mathscr{Y}^{\circ},\mathscr{Y}^{\circ}}^{\varphi,\geq}\ar[r] & \Modif_{X}^{\geq}}
\quad\text{and}\quad\xymatrix{\Modif_{\mathscr{Y}^{\circ},\mathscr{Y}^{\circ}}^{\varphi,ad,\geq}\ar[r] & \Modif_{X}^{ad,\geq}}.
\]
Putting everything together, we obtain an equivalence of $\otimes$-categories
\[
\xymatrix{\underline{\mathcal{E}}:\Mod_{A,f}^{\varphi,\leq}\otimes E\ar[r] & \Modif_{X}^{ad,\geq}.}
\]
This is of course the restriction of the $\otimes$-functor
\[
\xymatrix{\underline{\mathcal{E}}:\Mod_{A[\frac{1}{\pi}]}^{\varphi,\leq}\ar[r] & \Modif_{\mathscr{Y}^{\circ},\mathscr{Y}^{+}}^{\varphi,\geq}\ar[r] & \Modif_{\mathscr{Y}^{\circ},\mathscr{Y}^{\circ}}^{\varphi,\geq}\ar[r] & \Modif_{X}^{\geq}}
\]
but the first two components of the latter may not be equivalences.

\subsubsection{Compatibility with Hodge filtrations}\label{subsec:CompHodge}

The morphisms of locally ringed space
\[
\mathscr{Y}^{\circ}\rightarrow\mathscr{X}\rightarrow X\quad\text{and}\quad\mathscr{Y}^{\circ}\rightarrow\Spa(A)\rightarrow\Spec(A)
\]
map $y_{i}\in\left|\mathscr{Y}^{\circ}\right|$ to respectively $\infty\in\left|X\right|$ and $\mathfrak{m}_{i}=A\varphi^{-i}(\xi)\in\left|\Spec(A)\right|$. Moreover, they induce isomorphism between the corresponding completed local rings $\mathscr{O}_{\mathscr{Y},y_{i}}^{\wedge}$, $\mathcal{O}_{X,\infty}^{\wedge}$ and $A_{\mathfrak{m}_{i}}^{\wedge}=B_{\dR,\mathfrak{m}_{i}}^{+}$. For $i=0$, the latter is just $B_{\dR}^{+}$. For $(N,\varphi_{N})$ in $\Mod_{A[\frac{1}{\pi}]}^{\varphi,\leq}$ mapping to $(\mathscr{E},\psi_{\mathscr{E}})$ in $\Modif_{\mathscr{Y}^{+}}^{\ast}$ and $\underline{\mathcal{E}}=(\mathcal{E}_{1}\hookrightarrow\mathcal{E}_{2})$ in $\Modif_{X}^{\geq}$, we thus find
\begin{align*}
\left(\mathcal{E}_{1,\infty}^{\wedge}\hookrightarrow\mathcal{E}_{2,\infty}^{\wedge}\right) & \simeq \left(\mathscr{\mathscr{E}}(-\infty)_{y_{0}}^{\wedge}\hookrightarrow\mathscr{\mathscr{E}}(+\infty)_{y_{0}}^{\wedge}\right)\\
& = \left(\mathscr{\mathscr{E}}(-1)_{y_{0}}^{\wedge}\hookrightarrow\mathscr{\mathscr{E}}(0)_{y_{0}}^{\wedge}\right)\\
& \simeq \left((\varphi^{-1})^{\ast}(\varphi_{N}^{-1}):(\varphi^{-1})^{\ast}N\otimes B_{\dR}^{+}\hookrightarrow N\otimes B_{\dR}^{+}\right).
\end{align*}
It follows that
\begin{align*}
\mathcal{F}_{H,2}(\underline{\mathcal{E}}) & = \mathcal{F}_{H}^{\iota}(N,\varphi_{N})\quad\text{ on}\quad\mathcal{E}_{2}(\infty)=N\otimes_{A}C\\
\mathcal{F}_{H,1}(\underline{\mathcal{E}})\otimes_{C,\varphi}C' & = \mathcal{F}_{H}(N,\varphi_{N})\quad\text{ on}\quad\mathcal{E}_{1}(\infty)\otimes_{C}C'=N\otimes_{A}C'
\end{align*}


\subsubsection{Compatibility with Tate objects}\label{subsec:CompTate}

The Tate object over $A$ is anti-effective,
\[
A\{1\}=\left({\frac{1}{\mu}}A\otimes\mathcal{O}_{E}(1),\varphi\otimes\mathrm{Id}\right).
\]
The corresponding sequence $\dots\rightarrow\mathscr{E}(i)\rightarrow\mathscr{E}(i+1)\rightarrow\cdots$ is obtained from
\[
\xymatrix@C=2pc{\cdots\ar[r] & \frac{1}{\varphi^{-2}(\mu)}A\,\ar@{^{(}->}[r] & \frac{1}{\varphi^{-1}(\mu)}A\,\ar@{^{(}->}[r] & \frac{1}{\mu}A\,\ar@{^{(}->}[r] & \frac{1}{\varphi(\mu)}A\,\ar@{^{(}->}[r] & \frac{1}{\varphi^{2}(\mu)}A\,\ar@{^{(}->}[r] & \cdots}
\]
by tensoring with ${\,\cdot\,}\otimes_{A}(\mathscr{O}_{\mathscr{Y}}\otimes\mathcal{O}_{E}(1))$. Thus by~\cite[3.23]{BaMoSc16},
\[
\left(\mathscr{E}^{-}\vert_{\mathscr{Y}^{\circ}}\hookrightarrow\mathscr{E}^{+}\vert_{\mathscr{Y}^{\circ}}\right)=\left(\mathscr{O}_{\mathscr{Y}^{\circ}}\otimes E(1)\hookrightarrow\mathscr{O}_{\mathscr{Y}^{\circ}}(1)\otimes E(1)\right)
\]
where $\mathscr{O}_{\mathscr{Y}^{\circ}}\hookrightarrow\mathscr{O}_{\mathscr{Y}^{\circ}}(1)$ maps to $\mathcal{O}_{X}\hookrightarrow\mathcal{O}_{X}(1)$, therefore
\[
\underline{\mathcal{E}}\left(A\Mk\left[{\frac{1}{\pi}}\right]\{1\}\right)=\left(\mathcal{O}_{X}\otimes E(1)\hookrightarrow\mathcal{O}_{X}(1)\otimes E(1)\right)=\underline{\mathcal{O}}_{X}\{1\}.
\]
The $\otimes$-functor constructed in~\ref{subsec:ConstruFarguesFunctEff} thus extends to a $\otimes$-functor
\[
\underline{\mathcal{E}}:\Mod_{A[\frac{1}{\pi}],f}^{\varphi}\rightarrow\Modif_{X}
\]
mapping $N$ to $\underline{\mathcal{E}}(N\{i\})\{-i\}$ for $i\gg0$. The latter is still compatible with Hodge filtrations by~\ref{subsec:Hodge-and-NewtonForModif} and Proposition~\ref{prop:t_H(twist)4A=00005B1/pi=00005D}, and it induces an equivalence of $\otimes$-categories
\[
\underline{\mathcal{E}}:\Mod_{A,f}^{\varphi}\otimes E\rightarrow\Modif_{X}^{ad}.
\]

\subsubsection{Compatibility with Newton types}\label{subsec:CompNewton}

For $(N,\varphi_{N})$ in $\Mod_{A,f}^{\varphi}\otimes E$ of rank $r\in\mathbb{N}$ mapping to $\underline{\mathcal{E}}=(\mathcal{E}_{1},\mathcal{E}_{2},f)$ in $\Modif_{X}^{ad}$ and to $(D,\varphi_{D})=(M,\varphi_{M})\otimes_{A}L$ in $\Bun_{L}^{\varphi}$,
\[
t_{N}(\underline{\mathcal{E}})=t_{N}(\mathcal{E}_{2})\quad\text{equals}\quad t_{N}^{\iota}(D,\varphi_{D})\quad\text{in}\quad\mathbb{Q}_{\geq}^{r}.
\]
Indeed, we may assume that $(N,\varphi_{N})=(M,\varphi_{M})\otimes E$ for an anti-effective finite free BKF-module $(M,\varphi_{M})$ over $A$ by~\ref{subsec:RisL} and~\ref{subsec:Hodge-and-NewtonForModif}. If $(M,\varphi_{M})$ maps to $(\mathscr{E},\psi_{\mathscr{E}})$ and $(\mathscr{E}^{+},\mathscr{E}^{-},f_{\mathscr{E}})$ as above, then $\mathcal{E}_{2}$ is the image of $(\mathscr{E}^{+},\varphi_{\mathscr{E}^{+}})$ under
\[
\mathcal{E}\circ\Gamma(\mathscr{Y}^{\circ},{\cdot\,}):\Bun_{\mathscr{Y}^{+}}^{\varphi}\rightarrow\Bun_{X}
\]
thus $t_{N}(\mathcal{E}_{2})=t_{N}^{\iota}(D',\varphi_{D'})$ by~\ref{subsec:PhiBunY=0000B0} where $(D^{\prime},\varphi_{D^{\prime}})$ is the image of $(\mathscr{E}^{+},\varphi_{\mathscr{E}^{+}})$ under
\[
({\,\cdot\,})_{y_{L}}\otimes_{\mathscr{R}_{+}^{\mathrm{int}}}L:\Bun_{\mathscr{Y}^{+}}^{\varphi}\rightarrow\Bun_{L}^{\varphi}.
\]
Since $(\mathscr{E},\psi_{\mathscr{E}})\vert_{\mathscr{Y}^{+}}\hookrightarrow(\mathscr{E}^{+},\varphi_{\mathscr{E}^{+}}^{-1})$ is an isomorphism over $\mathscr{Y}_{]-1,\infty]}$, it induces
\[
(M,\varphi_{M}^{-1})\otimes_{A}\mathscr{R}_{+}^{\mathrm{int}}=(\mathscr{E},\psi_{\mathscr{E}})_{y_{L}}\stackrel{\simeq}{\longrightarrow}(\mathscr{E}^{+},\varphi_{\mathscr{E}^{+}}^{-1})_{y_{L}}
\]
therefore also $(D,\varphi_{D})=(M,\varphi_{M})\otimes_{A}L\simeq(D',\varphi_{D'})$, which proves our claim.

\subsubsection{Compatibility with Bhatt--Morrow--Scholze}\label{subsec:CompScholzeFargues}

We now claim that the $\otimes$-functor
\[
\HTT\circ\underline{\mathcal{E}}:\Mod_{A,f}^{\varphi}\otimes E\rightarrow\Modif_{X}^{ad}\rightarrow\HT_{E}^{B_{\dR}}
\]
is canonically isomorphic to the Bhatt--Morrow--Scholze $\otimes$-functor
\[
\HTT':\Mod_{A,f}^{\varphi}\otimes E\rightarrow\HT_{E}^{B_{\dR}}
\]
of Section~\ref{subsec:The-Bhatt--Morrow--Scholze-functor}. Since both functors are compatible with Tate twists, it is sufficient to establish that they have canonically isomorphic restrictions to the full $\otimes$-subcategory of anti-effective objects in $\Mod_{A,f}^{\varphi}\otimes E$, and this immediately follows from the computations in Section~\ref{subsec:ComputeAnalFunct} and~\ref{subsec:CompHodge}.

\subsubsection{Proof of Theorem~\ref{thm:FarguesScholze}}

It remains to establish that the $\otimes$-functor
\[
\HTT':\Mod_{A,f}^{\varphi}\rightarrow\HT_{\mathcal{O}_{E}}^{B_{\dR}}
\]
is an equivalence of $\otimes$-categories. Consider the ($2-)$commutative diagram
\[
\xymatrix@C=3pc@R=2pc{\Mod_{A,f}^{\varphi}\ar[r]^{\HTT'}\ar[d]_{{\,\cdot\,}\otimes E} & \HT_{\mathcal{O}_{E}}^{B_{\dR}}\ar[r]^{T}\ar[d]_{{\,\cdot\,}\otimes E} & \Bun_{\mathcal{O}_{E}}\ar[d]^{{\,\cdot\,}\otimes E}\\
\Mod_{A,f}^{\varphi}\otimes E\ar[r]^{\HTT'} & \HT_{E}^{B_{\dR}}\ar[r]^{V} & \Vect_{E} }
\]
Since the second square is cartesian, it is sufficient to establish that the outer rectangle is cartesian, for then so will be the first square, and its top row will thus be an equivalence of categories since so is the second row. We may again restrict our attention to anti-effective objects. The outer rectangle then factors as
\[
\xymatrix@C=3pc@R=2pc{\Mod_{A,f}^{\varphi,\leq}\ar[r]^{\underline{\mathscr{E}}}\ar[d]_{{\,\cdot\,}\otimes E} & \Modif_{\mathscr{Y}^{-},\mathscr{Y}^{+}}^{\varphi,\geq}\ar[r]^{\mathscr{E}^{-}}\ar[d]_{{\,\cdot\,}\vert_{\mathscr{Y}^{\circ}}} & \Bun_{\mathscr{Y}^{-}}^{\varphi}\ar[d]_{{\,\cdot\,}\vert_{\mathscr{Y}^{\circ}}}\ar[r] & \Bun_{\mathcal{O}_{E}}\ar[d]^{{\,\cdot\,}\otimes E}\\
\Mod_{A,f}^{\varphi,\leq}\otimes E\ar[r]^{\underline{\mathscr{E}}} & \Modif_{\mathscr{Y}^{\circ},\mathscr{Y}^{+}}^{\varphi,ad,\geq}\ar[r]^{\mathscr{E}^{-}} & \Bun_{\mathscr{Y}^{\circ}}^{\varphi,0}\ar[r] & \Vect_{E} }
\]
In this commutative diagram, the first square is cartesian since the two $\underline{\mathscr{E}}$'s are equivalences of $\otimes$-categories, the second square is obviously cartesian, and the third square is cartesian by Kedlaya's theory as explained in~\ref{subsec:PhiBunY=0000B0}. So the outer rectangle is indeed cartesian. This finishes the proof of Theorem~\ref{thm:FarguesScholze}.

\subsubsection{Final questions}

Is it true that every $\varphi$-bundle over $\mathscr{Y}^{+}$ is finite and free? Is there an integral version of the Fargues--Fontaine curve $X$ corresponding to $\mathscr{Y}^{-}/\varphi^{\mathbb{Z}}$? And is it true that $\underline{\mathcal{E}}:\Mod_{A[\frac{1}{\pi}],f}^{\varphi}\rightarrow\Modif_{X}$ is an equivalence of $\otimes$-categories?

\subsection{Exactness}\label{subsec:Exactness}

\subsubsection{}

Recall that Theorem~\ref{thm:FarguesScholze} implies that $\Mod_{A,f}^{\varphi}$ is a quasi-abelian category. In particular, any morphism $f$ in $\Mod_{A,f}^{\varphi}$ has a kernel and a cokernel, whose underlying $A$-modules may however not be the kernel and cokernel of $\alpha$ as computed in the abelian category of all $A$-modules. Accordingly, we say that a three term complex
\[
0\rightarrow(M_{1},\varphi_{1})\stackrel{f}{\longrightarrow}(M_{2},\varphi_{2})\stackrel{g}{\longrightarrow}(M_{3},\varphi_{3})\rightarrow0
\]
in $\Mod_{A,f}^{\varphi}$ is canonically (resp.~naively) exact if $f=\ker(g)$ and $g=\coker(f)$ in $\Mod_{A,f}^{\varphi}$ (resp.~in $\Mod_{A}$). We want to investigate the difference between canonical and naive short exact sequences in $\Mod_{A,f}^{\varphi}$.

\subsubsection{}

Plainly, a naive short exact sequence is also canonically exact. Conversely, let us start with a fixed canonical short exact sequence in $\Mod_{A,f}^{\varphi}$,
\[
0\rightarrow(M_{1},\varphi_{1})\rightarrow(M_{2},\varphi_{2})\rightarrow(M_{3},\varphi_{3})\rightarrow0.
\]
The corresponding complex of Hodge--Tate module is a short exact sequence
\[
0\rightarrow(T_{1},\Xi_{1})\rightarrow(T_{2},\Xi_{2})\rightarrow(T_{3},\Xi_{3})\rightarrow0
\]
and we now know what it means: the underlying complexes of $\mathcal{O}_{E}$ and $B_{\dR}^{+}$-modules are both exact. Since $M_{i}[\frac{1}{\mu}]\simeq T_{i}\otimes A[\frac{1}{\mu}]$, it follows that
\[
0\rightarrow M_{1}\!\left[{\frac{1}{\mu}}\right]\rightarrow M_{2}\!\left[{\frac{1}{\mu}}\right]\rightarrow M_{3}\!\left[{\frac{1}{\mu}}\right]\rightarrow0
\]
is exact. In particular, $M_{1}\rightarrow M_{2}$ is injective.

\subsubsection{}

Let $(Q,\varphi_{Q})$ be the cokernel of $(M_{1},\varphi_{1})\hookrightarrow(M_{2},\varphi_{2})$ in the abelian category $\Mod_{A}^{\varphi}$. Then $Q$ is the cokernel of $M_{1}\hookrightarrow M_{2}$ in $\Mod_{A}$, therefore $Q$ is a perfect $A$-module of projective dimension $\leq1$ with $Q[\frac{1}{\mu}]\simeq M_{3}[\frac{1}{\mu}]$ finite free over $A[\frac{1}{\mu}]$.

\begin{lemm}
If $Q[\pi^{\infty}]$ is finitely presented over $A$, then $Q[\pi^{\infty}]=0$.
\end{lemm}

\begin{proof}
Suppose that $Q[\pi^{\infty}]$ is finitely presented over $A$. Its inverse image $M_{1}^{\prime}$ in $M_{2}$ is then a finitely presented $A$-module with $M_{1}^{\prime}[\frac{1}{\pi}]\simeq M_{1}[\frac{1}{\pi}]$ free over $A[\frac{1}{\pi}]$, so $M_{1}^{\prime}$ is a torsion-free BKF-module. Then $M_{1}\subset M_{1}^{\prime}\subset M_{1,f}^{\prime}\subset M_{2}$ with $M_{1,f}^{\prime}/M_{1}$ killed by $\pi^{n}$ for $n\gg0$, so $M_{1,f}^{\prime}$ is contained in the kernel of $(M_{2},\varphi_{2})\rightarrow(M_{3},\varphi_{3})$ in $\Mod_{A,f}^{\varphi}$, i.e.~$M_{1,f}^{\prime}\hookrightarrow M_{2}$ factors through $M_{1}\hookrightarrow M_{2}$, which means that actually $M_{1}=M_{1}^{\prime}=M_{1,f}^{\prime}$, and indeed $Q[\pi^{\infty}]=M_{1}^{\prime}/M_{1}=0$.
\end{proof}

\subsubsection{}

Recall that $B_{crys}^{+}=A_{crys}[\frac{1}{\pi}]$ where $A_{crys}$ is the $\pi$-adic completion of the $A$-subalgebra of $A[\frac{1}{\pi}]$ generated by $\frac{\xi^{m}}{m!}$ for all $m\geq0$.

\begin{prop}\label{prop:ConditionsExactAlgebraic}
The following conditions are equivalent:
\begin{enumerate}
\item \label{3.14.1} Our complex induces an exact sequence of $B_{crys}^{+}$-modules
\[
0\rightarrow M_{1}\otimes_{A}B_{crys}^{+}\rightarrow M_{2}\otimes_{A}B_{crys}^{+}\rightarrow M_{3}\otimes_{A}B_{crys}^{+}\rightarrow0.
\]
\item \label{3.14.2} Our complex induces an exact sequence of $B_{crys}^{+}$-modules
\[
M_{1}\otimes_{A}B_{crys}^{+}\rightarrow M_{2}\otimes_{A}B_{crys}^{+}\rightarrow M_{3}\otimes_{A}B_{crys}^{+}\rightarrow0.
\]
\item \label{3.14.3} $Q[\frac{1}{\pi}]$ is free over $A[\frac{1}{\pi}]$.
\item \label{3.14.4} $Q[\frac{1}{\pi}]$ is projective over $A[\frac{1}{\pi}].$
\item \label{3.14.5} Our complex induces an exact sequence of $A$-modules
\[
0\rightarrow M_{1}\rightarrow M_{2}\rightarrow M_{3}\rightarrow\overline{Q}\rightarrow0
\]
with $\overline{Q}$ supported at $\{\mathfrak{m}\}$, i.e.~$\overline{Q}\in\Mod_{A,\mathfrak{m}^{\infty}}$.
\item \label{3.14.6} Our complex induces an exact sequence
\[
0\rightarrow\tilde{M}_{1}\rightarrow\tilde{M}_{2}\rightarrow\tilde{M}_{3}\rightarrow0
\]
of quasi-coherent sheaves on $U=\Spec(A)\setminus\{\mathfrak{m}\}$.
\item \label{3.14.7} Our complex induces an exact sequence of $A[\frac{1}{\pi}]$-modules
\[
0\rightarrow M_{1}\!\left[{\frac{1}{\pi}}\right]\rightarrow M_{2}\!\left[{\frac{1}{\pi}}\right]\rightarrow M_{3}\!\left[{\frac{1}{\pi}}\right]\rightarrow0
\]
\item \label{3.14.8} Our complex is isogeneous to a complex which is naively exact.
\end{enumerate}
\end{prop}

\begin{proof}
$\eqref{3.14.1}\Rightarrow\eqref{3.14.2}$ is obvious. $\eqref{3.14.2}\Rightarrow\eqref{3.14.3}$ follows from~\cite[4.19]{BaMoSc16}. $\eqref{3.14.3}\Leftrightarrow\eqref{3.14.4}$ is~\cite[4.12]{BaMoSc16}.

$\eqref{3.14.3}\Rightarrow\eqref{3.14.5}$ The assumption says that $Q$ is a BKF-module. Then $Q[\pi^{\infty}]$ is finitely presented, hence trivial by the previous lemma. It is then obvious that
\[
M_{2}\twoheadrightarrow Q\hookrightarrow Q_{f}
\]
is a cokernel of $M_{1}\hookrightarrow M_{2}$ in $\Mod_{A,f}^{\varphi}$, which proves~\eqref{3.14.5} with $M_{3}=Q_{f}$.

$\eqref{3.14.5}\Rightarrow\eqref{3.14.6}\Rightarrow\eqref{3.14.7}\Rightarrow\eqref{3.14.1}$ and $\eqref{3.14.8}\Rightarrow\eqref{3.14.7}$ are obvious.

$\eqref{3.14.5}\Rightarrow\eqref{3.14.8}$: if $\pi^{n}\overline{Q}=0$, the pull-back through multiplication by $\pi^{n}$ on $M_{3}$ yields an exact sequence
\[
\xymatrix@C=1.5pc@R=1.5pc{
0 \ar[r]& M_{1} \ar[r]\ar@{=}[d]& M_{2}^{\prime} \ar[r]\ar[d]& M_{3} \ar[r]\ar[d]^{\pi^n}& 0\\
0 \ar[r]& M_{1} \ar[r]& M_{2} \ar[r]& M_{3} \ar[r]& \overline{Q} \ar[r]& 0
}
\]
of the desired form, i.e.~isogeneous to the initial sequence and naively exact.
\end{proof}

\subsubsection{}

Suppose that our BKF-modules are anti-effective and let
\begin{gather*}
0\rightarrow(\mathscr{E}_{1},\psi_{1})\rightarrow(\mathscr{E}_{2},\psi_{2})\rightarrow(\mathscr{E}_{3},\psi_{3})\rightarrow0
\\
0\rightarrow(\mathscr{E}_{1}^{-}\hookrightarrow\mathscr{E}_{1}^{+})\rightarrow(\mathscr{E}_{2}^{-}\hookrightarrow\mathscr{E}_{2}^{+})\rightarrow(\mathscr{E}_{3}^{-}\hookrightarrow\mathscr{E}_{3}^{+})\rightarrow0
\end{gather*}
be the corresponding complexes in $\Modif_{\mathscr{Y}}^{\star}$ and $\Modif_{\mathscr{Y}^{-},\mathscr{Y}^{+}}^{\varphi,\geq}$. Note that
\[
0\rightarrow\mathscr{E}_{1}^{-}\rightarrow\mathscr{E}_{2}^{-}\rightarrow\mathscr{E}_{3}^{-}\rightarrow0
\]
is a (split) short exact sequence of sheaves on $\mathscr{Y}^{-}$ since $\mathscr{E}_{i}^{-}=T_{i}\otimes_{\mathcal{O}_{E}}\mathscr{O}_{\mathscr{Y}^{-}}$.

\begin{prop}\label{prop:ConditionsExactnessAdic}
The conditions of Proposition~\ref{prop:ConditionsExactAlgebraic} are equivalent to:
\begin{enumerate}
\item \label{1} Anyone of the following complexes is exact:
\begin{enumerate}
\item \label{1a}
$0\rightarrow\mathscr{E}_{1}\rightarrow\mathscr{E}_{2}\rightarrow\mathscr{E}_{3}\rightarrow0$ in $\Bun_{\mathscr{Y}}$.
\item \label{1b} $0\rightarrow\mathscr{E}_{1}\vert_{\mathscr{Y}^{+}}\rightarrow\mathscr{E}_{2}\vert_{\mathscr{Y}^{+}}\rightarrow\mathscr{E}_{3}\vert_{\mathscr{Y}^{+}}\rightarrow0$ in $\Bun_{\mathscr{Y}^{+}}$.
\item \label{1c} $0\rightarrow\mathscr{E}_{1}^{+}\rightarrow\mathscr{E}_{2}^{+}\rightarrow\mathscr{E}_{3}^{+}\rightarrow0$ in $\Bun_{\mathscr{Y}^{+}}$.
\item \label{1d} $0\rightarrow M_{1}\otimes\mathscr{R}_{+}^{\mathrm{int}}\rightarrow M_{2}\otimes\mathscr{R}_{+}^{\mathrm{int}}\rightarrow M_{3}\otimes\mathscr{R}_{+}^{\mathrm{int}}\rightarrow0$ in $\Mod_{\mathscr{R}_{+}^{\mathrm{int}}}$.
\item \label{1e} $0\rightarrow M_{1}\otimes\overline{B}\rightarrow M_{2}\otimes\overline{B}\rightarrow M_{3}\otimes\overline{B}\rightarrow0$ in $\Mod_{\overline{B}}$.
\end{enumerate}
\item \label{2} Anyone of the following complexes is split exact.
\begin{enumerate}
\item \label{2a} $0\rightarrow\mathscr{E}_{1}^{+}\rightarrow\mathscr{E}_{2}^{+}\rightarrow\mathscr{E}_{3}^{+}\rightarrow0$ in $\Bun_{\mathscr{Y}^{+}}^{\varphi}$.
\item \label{2b} $0\rightarrow\mathscr{E}_{1}^{+}\vert_{\mathscr{Y}^{\circ}}\rightarrow\mathscr{E}_{2}^{+}\vert_{\mathscr{Y}^{\circ}}\rightarrow\mathscr{E}_{3}^{+}\vert_{\mathscr{Y}^{\circ}}\rightarrow0$ in $\Bun_{\mathscr{Y}^{\circ}}^{\varphi}$.
\item \label{2c} $0\rightarrow M_{1}\otimes\overline{B}\rightarrow M_{2}\otimes\overline{B}\rightarrow M_{3}\otimes\overline{B}\rightarrow0$ in $\Bun_{\overline{B}}^{\varphi}$.
\end{enumerate}
\end{enumerate}
\end{prop}

\begin{proof}
In~\eqref{1}, plainly $\eqref{1a}\Rightarrow\eqref{1b}$, moreover $\eqref{1b}\Rightarrow\eqref{1c}$ by construction of $\mathscr{E}\mapsto\mathscr{E}^{+}$, $\eqref{1c}\Rightarrow\eqref{1d}$ by localization at $y_{L}$ and $\eqref{1d}\Rightarrow\eqref{1e}$ by base change along $\mathscr{R}_{+}^{\mathrm{int}}\twoheadrightarrow\overline{B}$ (using that $M_{3}$ is free over $A$). Moreover, $\eqref{1c}\Rightarrow\eqref{1a}$ since $\mathscr{E}_{i}=\mathscr{E}_{i}^{-}$ on $\mathscr{Y}_{[-\infty,0[}$ and $\mathscr{E}_{i}=\mathscr{E}_{i}^{+}$ on $\mathscr{Y}_{]-1,+\infty]}$. Since $\Bun_{\overline{B}}^{\varphi}\simeq\Bun_{\mathscr{Y}^{+}}^{\prime,\varphi}\simeq\Bun_{\mathscr{Y}^{\circ}}^{\varphi}$, the three conditions of \eqref{2} are equivalent. Obviously $\eqref{2a}\Rightarrow\eqref{1c}$, and $\eqref{1e}\Rightarrow\eqref{2c}$ by~\cite[\S 11.1]{FaFo15}. Condition~\eqref{3.14.7} of Proposition~\ref{prop:ConditionsExactAlgebraic} implies~\eqref{1e}. Finally~\eqref{1a} implies condition~\eqref{3.14.5} of Proposition~\ref{prop:ConditionsExactAlgebraic} by the next proposition (since indeed $M_{i}=\Gamma(\mathscr{Y},\mathscr{E}_{i})$ and $\mathscr{E}_{1}=M_{1}\otimes_{A}\mathscr{O}_{\mathscr{Y}}\simeq\mathscr{O}_{\mathscr{Y}}^{r_{1}}$ with $r_{1}=\rank_{A}M_{1}$).
\end{proof}

\begin{prop}
We have $H^{1}(\mathscr{Y}^{+},\mathscr{O}_{\mathscr{Y}})=0=H^{1}(\mathscr{Y}^{-},\mathscr{O}_{\mathscr{Y}})$ and
\[
H^{1}(\mathscr{Y},\mathscr{O}_{\mathscr{Y}})\!\left[{\frac{1}{\pi}}\right]=H^{1}(\mathscr{Y},\mathscr{O}_{\mathscr{Y}})\!\left[{\frac{1}{[\varpi]}}\right]=0.
\]
\end{prop}

\begin{proof}
The following proof was indicated to us by Fargues. First since
\[
\mathscr{Y}^{-}=\cup_{s}\mathscr{Y}_{[-\infty,s]}\quad\text{and}\quad\mathscr{Y}^{+}=\cup_{r}\mathscr{Y}_{[r,+\infty]},
\]
we have exact sequences of $A$-modules
\[
0\rightarrow R^{1}\underleftarrow{\lim}H^{0}\left(\mathscr{Y}_{[-\infty,s]},\mathscr{O}_{\mathscr{Y}}\right)\rightarrow H^{1}\left(\mathscr{Y}^{-},\mathscr{O}_{\mathscr{Y}}\right)\rightarrow\underleftarrow{\lim}H^{1}\left(\mathscr{Y}_{[-\infty,s]},\mathscr{O}_{\mathscr{Y}}\right)\rightarrow0
\]
\[
0\rightarrow R^{1}\underleftarrow{\lim}H^{0}\left(\mathscr{Y}_{[r,+\infty]},\mathscr{O}_{\mathscr{Y}}\right)\rightarrow H^{1}\left(\mathscr{Y}^{+},\mathscr{O}_{\mathscr{Y}}\right)\rightarrow\underleftarrow{\lim}H^{1}\left(\mathscr{Y}_{[r,+\infty]},\mathscr{O}_{\mathscr{Y}}\right)\rightarrow0
\]
By~\cite[2.7.7]{KeLi15}, $H^{1}\left(\mathscr{Y}_{[-\infty,s]},\mathscr{O}_{\mathscr{Y}}\right)=0=H^{1}\left(\mathscr{Y}_{[r,+\infty]},\mathscr{O}_{\mathscr{Y}}\right)$, thus
\[
\underleftarrow{\lim}H^{1}\left(\mathscr{Y}_{[-\infty,s]},\mathscr{O}_{\mathscr{Y}}\right)=0=\underleftarrow{\lim}H^{1}\left(\mathscr{Y}_{[r,+\infty]},\mathscr{O}_{\mathscr{Y}}\right).
\]
On the other hand the images of the restriction maps
\[
H^{0}\left(\mathscr{Y}^{-},\mathscr{O}_{\mathscr{Y}}\right)\rightarrow H^{0}\left(\mathscr{Y}_{[-\infty,s]},\mathscr{O}_{\mathscr{Y}}\right)\quad\text{and}\quad H^{0}\left(\mathscr{Y}^{+},\mathscr{O}_{\mathscr{Y}}\right)\rightarrow H^{0}\left(\mathscr{Y}_{[r,+\infty]},\mathscr{O}_{\mathscr{Y}}\right)
\]
are dense in their complete codomain, hence by the Mittag--Leffler lemma,
\[
R^{1}\underleftarrow{\lim}H^{0}\left(\mathscr{Y}_{[-\infty,s]},\mathscr{O}_{\mathscr{Y}}\right)=0=R^{1}\underleftarrow{\lim}H^{0}\left(\mathscr{Y}_{[r,+\infty]},\mathscr{O}_{\mathscr{Y}}\right).
\]
Thus indeed $H^{1}(\mathscr{Y}^{+},\mathscr{O}_{\mathscr{Y}})=0=H^{1}(\mathscr{Y}^{-},\mathscr{O}_{\mathscr{Y}})$. The Mayer--Vietoris sequence gives
\[
H^{1}\left(\mathscr{Y},\mathscr{O}_{\mathscr{Y}}\right)=\coker\left(H^{0}\left(\mathscr{Y}^{-},\mathscr{O}_{\mathscr{Y}}\right)\oplus H^{0}\left(\mathscr{Y}^{+},\mathscr{O}_{\mathscr{Y}}\right)\rightarrow H^{0}\left(\mathscr{Y}^{\circ},\mathscr{O}_{\mathscr{Y}}\right)\right).
\]
One checks that this cokernel is indeed annihilated by $-[\frac{1}{\pi}]$ and $-[\frac{1}{[\varpi]}]$.
\end{proof}

\subsubsection{}

Returning to the general case, let
\[
0\rightarrow(\mathcal{E}_{1,s},\mathcal{E}_{1,t},f_{1})\rightarrow(\mathcal{E}_{2,s},\mathcal{E}_{2,t},f_{2})\rightarrow(\mathcal{E}_{3,s},\mathcal{E}_{3,t},f_{3})\rightarrow0
\]
be the image of our canonical short exact sequence in $\Modif_{X}^{ad}$. Then
\[
0\rightarrow\mathcal{E}_{1,s}\rightarrow\mathcal{E}_{2,s}\rightarrow\mathcal{E}_{3,s}\rightarrow0\quad\text{and}\quad0\rightarrow\mathcal{E}_{1,t}\rightarrow\mathcal{E}_{2,t}\rightarrow\mathcal{E}_{3,t}\rightarrow0
\]
are short exact sequences in $\Bun_{X}$, and the first one is even split.

\begin{prop}\label{prop:CritExactOnModif}
The conditions of Proposition~\ref{prop:ConditionsExactAlgebraic} are equivalent to:

The exact sequence $0\rightarrow\mathcal{E}_{1,t}\rightarrow\mathcal{E}_{2,t}\rightarrow\mathcal{E}_{3,t}\rightarrow0$ is split.
\end{prop}

\begin{proof}
Using the compatibility of $\underline{\mathcal{E}}:\Mod_{A,f}^{\varphi}\otimes E\rightarrow\Modif_{X}^{ad}$ with Tate twists, we may assume that our BKF-modules are anti-effective. Our claim then follows from the criterion~\eqref{2b} of Proposition~\ref{prop:ConditionsExactnessAdic} since $\mathcal{E}_{i,t}=({\,\cdot\,}/\varphi)^{alg}(\mathscr{E}_{i}^{+}\vert_{\mathscr{Y}^{\circ}})$ with the local notations, and $({\,\cdot\,}/\varphi)^{alg}:\Bun_{\mathscr{Y}^{\circ}}^{\varphi}\rightarrow\Bun_{X}$ is an exact equivalence of categories.
\end{proof}

\subsubsection{Application}

Let $\Modif_{X}^{ad,\ast}$ be the strictly full subcategory of $\Modif_{X}^{ad}$ whose objects are the admissible modifications $(\mathcal{E}_{1},\mathcal{E}_{2},f_{\mathcal{E}})$ such that the $\mathbb{Q}$-filtration on $\mathcal{E}_{2}$ induced by the Fargues $\mathbb{Q}$-filtration of $(\mathcal{E}_{1},\mathcal{E}_{2},f_{\mathcal{E}})$ is split. Recall from~\ref{subsec:DefFullSubCatHNType} that $\Mod_{A,f}^{\varphi,\ast}$ is the strictly full subcategory of $\Mod_{A,f}^{\varphi}$ whose objects are the finite free BKF-modules of HN-type (Definition~\ref{Def:HNType}).

\begin{prop}\label{prop:SpecialSubCats}
Fix $(M,\varphi_{M})\in\Mod_{A,f}^{\varphi}\otimes E$ with image $(\mathcal{E}_{1},\mathcal{E}_{2},f_{\mathcal{E}})\in\Modif_{X}^{ad}$ and rank $r\in\mathbb{N}$. Then $t_{F,\infty}(M,\varphi_{M})(r)=t_{F}(\mathcal{E}_{1},\mathcal{E}_{2},f_{\mathcal{E}})(r)$ and for every $s\in[0,r]$,
\begin{align*}
(M,\varphi_{M})\in\Mod_{A,f}^{\varphi,\ast}\otimes E & \Longrightarrow t_{F,\infty}(M,\varphi_{M})(s)\leq t_{F}(\mathcal{E}_{1},\mathcal{E}_{2},f_{\mathcal{E}})(s),\\
(\mathcal{E}_{1},\mathcal{E}_{2},f_{\mathcal{E}})\in\Modif_{X}^{ad,\ast} & \Longrightarrow t_{F}(\mathcal{E}_{1},\mathcal{E}_{2},f_{\mathcal{E}})(s)\leq t_{F,\infty}(M,\varphi_{M})(s).
\end{align*}
If both condition holds, then $\underline{\mathcal{E}}$ maps the Fargues filtration $\mathcal{F}_{F}$ on $(M,\varphi_{M})$ (from Proposition~\ref{prop:TypeHN}) to the Fargues filtration $\mathcal{F}_{F}$ on $(\mathcal{E}_{1},\mathcal{E}_{2},f_{\mathcal{E}})$ (defined in Section~\ref{subsec:DefFarguesFiltrOnModif}).
\end{prop}

\begin{proof}
The first claim follows from~\ref{subsec:CompNewton}.

$(1)$
Suppose that $(M,\varphi_{M})$ belongs to $\Mod_{A,f}^{\varphi,\ast}$, so that $t_{F,\infty}(M,\varphi_{M})=t_{F}(M,\varphi_{M})$ by Proposition~\ref{prop:TypeHN}. The graph of $t_{F}(M,\varphi_{M})$ (resp. $t_{F}(\mathcal{E}_{1},\mathcal{E}_{2},f_{\mathcal{E}})$) is the concave upper bound of the convex hull of $\mathcal{A}$ (resp.~$\mathcal{B}$) where
\begin{align*}
\mathcal{A} & = \left\{ (\rank,\deg)\left(\mathcal{F}_{F}^{\gamma}(M,\varphi_{M})\right):\gamma\in\mathbb{R}\right\},\\
\mathcal{B} & = \left\{ (\rank,\deg)\left(\underline{\mathcal{E}}'\right):\underline{\mathcal{E}}'\text{ strict subobject of }(\mathcal{E}_{1},\mathcal{E}_{2},f_{\mathcal{E}})\text{ in }\Modif_{X}^{ad}\right\}.
\end{align*}
Now for every $\gamma\in\mathbb{R}$, the naively exact sequence
\[
0\rightarrow\mathcal{F}_{F}^{\gamma}M\rightarrow M\rightarrow M/\mathcal{F}_{F}^{\gamma}M\rightarrow0
\]
in $\Mod_{A,f}^{\varphi}$ induces a canonically exact sequence
\[
0\rightarrow\underline{\mathcal{E}}\left(\mathcal{F}_{F}^{\gamma}M\right)\rightarrow(\mathcal{E}_{1},\mathcal{E}_{2},f_{\mathcal{E}})\rightarrow\underline{\mathcal{E}}\left(M/\mathcal{F}_{F}^{\gamma}M\right)\rightarrow0
\]
in $\Modif_{X}^{ad}$. Thus $\mathcal{A}\subset\mathcal{B}$ and our claim easily follows.

$(2)$
Suppose that $(\mathcal{E}_{1},\mathcal{E}_{2},f_{\mathcal{E}})$ belongs to $\Modif_{X}^{ad,\ast}$. We need to show that for all $\gamma\in\mathbb{R}$, $d\leq f(s)$ where $f=t_{F,\infty}(M,\varphi_{M})$ and $(s,d)=(\rank,\deg)\left(\mathcal{F}_{F}^{\gamma}(\mathcal{E}_{1},\mathcal{E}_{2},f_{\mathcal{E}})\right)$. By assumption, Propositions~\ref{prop:CritExactOnModif} and~\ref{prop:DeftFinfty}, we may assume that the exact sequence
\[
0\rightarrow\mathcal{F}_{F}^{\gamma}(\mathcal{E}_{1},\mathcal{E}_{2},f_{\mathcal{E}})\rightarrow(\mathcal{E}_{1},\mathcal{E}_{2},f_{\mathcal{E}})\rightarrow\frac{(\mathcal{E}_{1},\mathcal{E}_{2},f_{\mathcal{E}})}{\mathcal{F}_{F}^{\gamma}(\mathcal{E}_{1},\mathcal{E}_{2},f_{\mathcal{E}})}\rightarrow0
\]
in $\Modif_{X}^{ad}$ arises from a naively exact sequence
\[
0\rightarrow M'\rightarrow M\rightarrow M''\rightarrow0
\]
in $\Mod_{A,f}^{\varphi}$, which gives rise to exact sequences
\[
0\rightarrow M'_{n}\rightarrow M_{n}\rightarrow M''_{n}\rightarrow0
\]
in $\Mod_{A,t}^{\varphi}$ for all $n\geq0$. Then by definition of $f'=t_{F,\infty}(M')$ and $f=t_{F,\infty}(M)$,
\[
d=\deg(M')=f'(s)=\lim_{n\rightarrow\infty}{\frac{1}{n}}t_{F}(M'_{n})(ns)\leq\lim_{n\rightarrow\infty}{\frac{1}{n}}t_{F}(M{}_{n})(ns)=f(s)
\]
using Proposition~\ref{prop:ExactSeqInModt} for the middle inequality.

$(3)$
Suppose now that both conditions hold. For $\gamma\in\mathbb{R}$, consider the image of the (naively) exact sequence
\[
0\rightarrow\mathcal{F}_{F}^{\gamma}M\rightarrow M\rightarrow M/\mathcal{F}_{F}^{\gamma}M\rightarrow0
\]
from~Proposition~\ref{prop:TypeHN}, which is an exact sequence in $\Modif_{X}^{ad}$,
\[
0\rightarrow\underline{\mathcal{E}}(\mathcal{F}_{F}^{\gamma}M)\rightarrow(\mathcal{E}_{1},\mathcal{E}_{2},f_{\mathcal{E}})\rightarrow\underline{\mathcal{E}}(M/\mathcal{F}_{F}^{\gamma}M)\rightarrow0.
\]
Set $(r_{\gamma},d_{\gamma})\mk=\mk(\rank,\deg)(\mathcal{F}_{F}^{\gamma}M)$, so that $f(r_{\gamma})\mk=\mk d_{\gamma}$ where $f\mk=\mk t_{F}(M)$. By~$(1)$~and~$(2)$, we know that $f=t_{F}(\mathcal{E}_{1},\mathcal{E}_{2},f_{\mathcal{E}})$, thus also $(r_{\gamma},d_{\gamma})=(\rank,\deg)(\mathcal{F}_{F}^{\gamma}(\mathcal{E}_{1},\mathcal{E}_{2},f_{\mathcal{E}}))$. It then follows from Proposition~\ref{prop:OnAdmModif_FarguesExactness} that $\underline{\mathcal{E}}(\mathcal{F}_{F}^{\gamma}M)=\mathcal{F}_{F}^{\gamma}\underline{\mathcal{E}}(\mathcal{F}_{F}^{\gamma}M)$. By functoriality of $\mathcal{F}_{F}$ on $\Modif_{X}^{ad}$, we find that $\underline{\mathcal{E}}(\mathcal{F}_{F}^{\gamma}M)\hookrightarrow(\mathcal{E}_{1},\mathcal{E}_{2},f_{\mathcal{E}})$ induces a monomorphism $\underline{\mathcal{E}}(\mathcal{F}_{F}^{\gamma}M)\hookrightarrow\mathcal{F}_{F}^{\gamma}(\mathcal{E}_{1},\mathcal{E}_{2},f_{\mathcal{E}})$. Since its domain and codomain have the same rank and degree, this monomorphism is indeed an isomorphism.
\end{proof}

\begin{thebibliography}{MMMMMM}
\bibitem[And09]{An09} André, Yves Slope filtrations, Confluentes Math., Volume 1 (2009) no. 1, pp. 1-85 \href{https://ahl.centre-mersenne.org/item/AHL_2019__2__415_0/#An09}{Publisher reference}.
\bibitem[BMS16]{BaMoSc16} Bhatt, Bhargav; Morrow, Matthew; Scholze, Peter Integral p-adic Hodge theory (2016) \url{https://arxiv.org/abs/1602.03148} (preprint) \href{https://ahl.centre-mersenne.org/item/AHL_2019__2__415_0/#BaMoSc16}{Publisher reference}.
\bibitem[Bou64]{BoAC56} Bourbaki, Nicolas Éléments de mathématique. Fasc. XXX. Algèbre commutative. Chapitre 5: Entiers. Chapitre 6: Valuations, Actualités Scientifiques et Industrielles, 1308, Hermann, 1964, 207 pages \href{https://ahl.centre-mersenne.org/item/AHL_2019__2__415_0/#BoAC56}{Publisher reference}.
\bibitem[BvdB03]{BoVdB03} Bondal, Alexey; van den Bergh, Michel Generators and representability of functors in commutative and noncommutative geometry, Mosc. Math. J., Volume 3 (2003) no. 1, pp. 1-36 \href{https://ahl.centre-mersenne.org/item/AHL_2019__2__415_0/#BoVdB03}{Publisher reference}.
\bibitem[CN16]{CoNi16} Cornut, Christophe; Nicole, Marc-Hubert Cristaux et immeubles, Bull. Soc. Math. Fr., Volume 144 (2016) no. 1, pp. 125-143 \href{https://ahl.centre-mersenne.org/item/AHL_2019__2__415_0/#CoNi16}{Publisher reference}.
\bibitem[Cor]{Co14} Cornut, Christophe Filtrations and buildings \url{http://webusers.imj-prg.fr/~christophe.cornut/papers.html} (to appear in Mem. Am. Math. Soc.) \href{https://ahl.centre-mersenne.org/item/AHL_2019__2__415_0/#Co14}{Publisher reference}.
\bibitem[Cor18]{Co16} Cornut, Christophe On Harder–Narasimhan filtrations and their compatibility with tensor products, Confluentes Math., Volume 10 (2018) no. 2, pp. 3-49 \href{https://ahl.centre-mersenne.org/item/AHL_2019__2__415_0/#Co16}{Publisher reference}.
\bibitem[Cor19]{Co15} Cornut, Christophe Mazur’s Inequality and Laffaille’s theorem, Math. Ann., Volume 374 (2019) no. 3-4, pp. 1657-1679 \href{https://ahl.centre-mersenne.org/item/AHL_2019__2__415_0/#Co15}{Publisher reference}.
\bibitem[Dem86]{De86} Demazure, Michel Lectures on p-divisible groups, Lecture Notes in Mathematics, 302, Springer, 1986, vi+98 pages (Reprint of the 1972 original) \href{https://ahl.centre-mersenne.org/item/AHL_2019__2__415_0/#De86}{Publisher reference}.
\bibitem[DOR10]{DaOrRa10} Dat, Jean-François; Orlik, Sascha; Rapoport, Michael Period domains over finite and p-adic fields, Cambridge Tracts in Mathematics, 183, Cambridge University Press, 2010, xxii+372 pages \href{https://ahl.centre-mersenne.org/item/AHL_2019__2__415_0/#DaOrRa10}{Publisher reference}.
\bibitem[Far10]{Fa10} Fargues, Laurent La filtration de Harder–Narasimhan des schémas en groupes finis et plats, J. Reine Angew. Math., Volume 645 (2010), pp. 1-39 \href{https://ahl.centre-mersenne.org/item/AHL_2019__2__415_0/#Fa10}{Publisher reference}.
\bibitem[Far15]{Fa15} Fargues, Laurent Quelques résultats et conjectures concernant la courbe, De la géométrie algébrique aux formes automorphes (I) (Astérisque), Société Mathématique de France, 2015 no. 369, pp. 325-374 \href{https://ahl.centre-mersenne.org/item/AHL_2019__2__415_0/#Fa15}{Publisher reference}.
\bibitem[Far16]{Fa16} Fargues, Laurent Geometrization of the local Langlands correspondence: an overview (2016) \url{https://arxiv.org/abs/1602.00999} (preprint) \href{https://ahl.centre-mersenne.org/item/AHL_2019__2__415_0/#Fa16}{Publisher reference}.
\bibitem[Far19]{Fa12} Fargues, Laurent Théorie de la réduction pour les groupes p-divisibles (2019) \url{https://arxiv.org/abs/1901.08270} (preprint) \href{https://ahl.centre-mersenne.org/item/AHL_2019__2__415_0/#Fa12}{Publisher reference}.
\bibitem[FF18]{FaFo15} Fargues, Laurent; Fontaine, Jean-Marc Courbes et fibrés vectoriels en théorie de Hodge p-adique, Astérisque, Société Mathématique de France, 2018 \href{https://ahl.centre-mersenne.org/item/AHL_2019__2__415_0/#FaFo15}{Publisher reference}.
\bibitem[Fon90]{Fo90} Fontaine, Jean-Marc Représentations p-adiques des corps locaux. I, The Grothendieck Festschrift, Vol. II (Progress in Mathematics), Volume 87, Birkhäuser, 1990, pp. 249-309 \href{https://ahl.centre-mersenne.org/item/AHL_2019__2__415_0/#Fo90}{Publisher reference}.
\bibitem[GR03]{GaRa03} Gabber, Ofer; Ramero, Lorenzo Almost ring theory, Lecture Notes in Mathematics, 1800, Springer, 2003, vi+307 pages \href{https://ahl.centre-mersenne.org/item/AHL_2019__2__415_0/#GaRa03}{Publisher reference}.
\bibitem[Gro05]{SGA2r} Grothendieck, Alexander Cohomologie locale des faisceaux cohérents et théorèmes de Lefschetz locaux et globaux (SGA 2), Documents Mathématiques, 4, Société Mathématique de France, 2005, x+208 pages (Séminaire de Géométrie Algébrique du Bois Marie, 1962, Augmenté d’un exposé de Michèle Raynaud) \href{https://ahl.centre-mersenne.org/item/AHL_2019__2__415_0/#SGA2r}{Publisher reference}.
\bibitem[HM07]{HaMa07} Hamilton, Tracy Dawn; Marley, Thomas Non-Noetherian Cohen–Macaulay rings, J. Algebra, Volume 307 (2007) no. 1, pp. 343-360 \href{https://ahl.centre-mersenne.org/item/AHL_2019__2__415_0/#HaMa07}{Publisher reference}.
\bibitem[Iri16]{Imp16} Irissarry, Macarena Peche The reduction of G-ordinary crystalline representations with G-structures (2016) \url{https://arxiv.org/abs/1611.08529} (preprint) \href{https://ahl.centre-mersenne.org/item/AHL_2019__2__415_0/#Imp16}{Publisher reference}.
\bibitem[Kap49]{Ka66} Kaplansky, Irving Elementary divisors and modules, Trans. Am. Math. Soc., Volume 66 (1949), pp. 464-491 \href{https://ahl.centre-mersenne.org/item/AHL_2019__2__415_0/#Ka66}{Publisher reference}.
\bibitem[Ked05]{Ke05} Kedlaya, Kiran S. Slope filtrations revisited, Doc. Math., Volume 10 (2005), pp. 447-525 \href{https://ahl.centre-mersenne.org/item/AHL_2019__2__415_0/#Ke05}{Publisher reference}.
\bibitem[Ked16a]{Ke16} Kedlaya, Kiran S. Noetherian properties of Fargues–Fontaine curves, Int. Math. Res. Not. (2016) no. 8, pp. 2544-2567 \href{https://ahl.centre-mersenne.org/item/AHL_2019__2__415_0/#Ke16}{Publisher reference}.
\bibitem[Ked16b]{Ke16b} Kedlaya, Kiran S. Some ring-theoretic properties of A inf  (2016) \url{https://arxiv.org/abs/1602.09016} (preprint) \href{https://ahl.centre-mersenne.org/item/AHL_2019__2__415_0/#Ke16b}{Publisher reference}.
\bibitem[KL15]{KeLi15} Kedlaya, Kiran S.; Liu, Ruochuan Relative p-adic Hodge theory: foundations, Astérisque, 371, Société Mathématique de France, 2015 \href{https://ahl.centre-mersenne.org/item/AHL_2019__2__415_0/#KeLi15}{Publisher reference}.
\bibitem[Kot85]{Ko85} Kottwitz, Robert E. Isocrystals with additional structure, Compos. Math., Volume 56 (1985) no. 2, pp. 201-220 \href{https://ahl.centre-mersenne.org/item/AHL_2019__2__415_0/#Ko85}{Publisher reference}.
\bibitem[LWE16]{LeWE16} Levin, Brandon; Wang-Erickson, Carl A Harder–Narasimhan theory for Kisin modules (2016) \url{https://arxiv.org/abs/1606.00914} (preprint) \href{https://ahl.centre-mersenne.org/item/AHL_2019__2__415_0/#LeWE16}{Publisher reference}.
\bibitem[Man63]{Ma63} Manin, Yuri I. Theory of commutative formal groups over fields of finite characteristic, Usp. Mat. Nauk, Volume 18 (1963) no. 6(114), pp. 3-90 \href{https://ahl.centre-mersenne.org/item/AHL_2019__2__415_0/#Ma63}{Publisher reference}.
\bibitem[Nor76]{No76} Northcott, D. G. Finite free resolutions, Cambridge Tracts in Mathematics, 71, Cambridge University Press, 1976, xii+271 pages \href{https://ahl.centre-mersenne.org/item/AHL_2019__2__415_0/#No76}{Publisher reference}.
\bibitem[Sch17]{Sc17} Scholze, Peter p-adic Geometry (2017) \url{https://arxiv.org/abs/1712.03708} (preprint) \href{https://ahl.centre-mersenne.org/item/AHL_2019__2__415_0/#Sc17}{Publisher reference}.
\bibitem[SW17]{ScWe15} Scholze, Peter; Weinstein, Jared Peter Scholze’s lectures on p-adic geometry, Notes by J. Weinstein (2017) \url{https://www.math.uni-bonn.de/people/scholze/Berkeley.pdf} (to appear in Annals of Mathematics Studies) \href{https://ahl.centre-mersenne.org/item/AHL_2019__2__415_0/#ScWe15}{Publisher reference}.
\bibitem[Zin84]{Zi84} Zink, Thomas Cartiertheorie kommutativer formaler Gruppen, Teubner-Texte zur Mathematik, 68, Teubner Verlagsgesellschaft, 1984, 124 pages \href{https://ahl.centre-mersenne.org/item/AHL_2019__2__415_0/#Zi84}{Publisher reference}.
\end{thebibliography}
\end{document}
